% Lectura y comentario : el cuarto poder y la promoción de la cultura de paz — Espagnol Tle A — Cours signé OnBuch+
% Programme officiel MINESEC. Compiler avec : tectonic E24-cuarto-poder-cultura-de-paz.tex
\newcommand{\DOCMATIERE}{Espagnol}
\newcommand{\DOCNIVEAU}{Tle A}
\newcommand{\DOCMODULE}{Módulo 5 : Médias, communication et culture de la paix}
\newcommand{\DOCLECON}{E24}
\newcommand{\DOCTITRE}{Lectura y comentario : el cuarto poder y la promoción de la cultura de paz}
% OnBuch+ — préambule « cours signés » ENRICHI (compilé avec tectonic / XeLaTeX)
% Système visuel « Pop » d'OnBuch+ : fond crème (jamais blanc), bordures encre
% épaisses, ombres « dures » décalées, titres Archivo Black en capitales.
% Version MATHÉMATIQUES : graphes (pgfplots/tikz), boîtes pédagogiques
% (définition, propriété, méthode, attention, le savais-tu, exercice résolu,
% exercice, TP, bilan), page de garde et signature OnBuch+ ancrée en bas.
%
% Macros attendues dans main.tex AVANT \input{preamble} :
%   \DOCMATIERE  \DOCNIVEAU  \DOCMODULE  \DOCLECON (numéro)  \DOCTITRE
\documentclass[11pt]{article}

\usepackage{fontspec}
\usepackage[spanish,french]{babel} % français = langue principale ; l'espagnol via \esp{..} / espagnol
\usepackage[a4paper,margin=2cm,top=3.4cm,bottom=2.8cm,headheight=1.4cm,footskip=1.3cm]{geometry}
\usepackage{amsmath,amssymb,amsfonts}
\usepackage{mathtools} % \xmapsto, \coloneqq... souvent utilisés par les modèles
\usepackage{cancel} % simplifications barrées (bilans de forces)
% \abs{z} (module, valeur absolue) et \norm{u} (norme), utilisés par les modèles
\providecommand{\abs}{}\let\abs\relax\DeclarePairedDelimiter{\abs}{\lvert}{\rvert}
\providecommand{\norm}{}\let\norm\relax\DeclarePairedDelimiter{\norm}{\lVert}{\rVert}
\usepackage[table]{xcolor} % [table] : fournit aussi \rowcolors (alternance de lignes)
\usepackage{pagecolor}
\usepackage{tikz}
\usetikzlibrary{arrows.meta,positioning,calc,shapes.geometric,shapes.misc,shapes.arrows,decorations.pathmorphing,decorations.pathreplacing,decorations.markings,patterns,shadows,mindmap,trees,fit,backgrounds,matrix,intersections,angles,quotes}
\usepackage{pgfplots}
\usepackage[version=4]{mhchem} % équations nucléaires \ce{^{235}_{92}U}
\usepackage[european,siunitx]{circuitikz} % schémas électriques
\pgfplotsset{compat=1.18}
\usepgfplotslibrary{fillbetween,groupplots} % aires entre courbes (calcul intégral)
\usepackage{enumitem}
\usepackage{fancyhdr}
% [most] charge toutes les bibliothèques tcolorbox (listings, minted, raster,
% documentation...) et gonfle inutilement la mémoire TeX sur un document long
% (~300 boîtes) : on ne charge que ce qui est réellement utilisé.
\usepackage[skins,breakable]{tcolorbox}
\usepackage{graphicx}
\usepackage{colortbl}
\usepackage{booktabs}
\usepackage{tabularx}
\usepackage{diagbox} % case d'en-tête diagonale des tableaux à double entrée
% Colonnes extensibles centrée / gauche / droite (C, L, R), souvent utilisées par les modèles
\newcolumntype{C}{>{\centering\arraybackslash}X}
\newcolumntype{L}{>{\raggedright\arraybackslash}X}
\newcolumntype{R}{>{\raggedleft\arraybackslash}X}
\usepackage{array}
\usepackage{multicol}
\usepackage{siunitx}
% parse-numbers=false : accepte aussi \SI{2,5 \times 10^{-3}}{...}
\sisetup{locale=FR,output-decimal-marker={,},group-minimum-digits=4,parse-numbers=false}
\DeclareSIUnit\ohm{\ensuremath{\symup{\Omega}}} % U+2126 absent de la police
\DeclareSIUnit\division{div} % réglages d'oscilloscope (ms/div, V/div)
\DeclareSIUnit\tr{tr} % tours (vitesse de rotation, ex. \SI{1500}{\tr\per\minute})
\DeclareSIUnit\molar{M} % concentration molaire (ex. \SI{3}{\molar}, \SI{1}{\micro\molar})
\DeclareSIUnit\an{an} % année (le modèle confond parfois avec \year, un compteur TeX) — ex. \SI{5}{\kilo\meter\per\an}
\DeclareSIUnit\FCFA{FCFA} % franc CFA (ex. \SI{500}{\FCFA\per\kilo\gram})
\DeclareSIUnit\degreeBrix{\textdegree Brix} % teneur en sucre d'un sirop/jus (réfractométrie)
\DeclareSIUnit\month{mois} % le modèle utilise parfois \month, un compteur TeX, comme unité de durée
\usepackage{qrcode}
% \degree : commande fréquemment utilisée par les modèles pour un angle ou une
% température en dehors de \SI{}{} (non fournie par les paquets chargés).
\providecommand{\degree}{^{\circ}}
% \qed (fin de démonstration), utilisé par les modèles sans amsthm
\providecommand{\qed}{\hfill$\square$}
% \barre : double barre des valeurs interdites dans un tableau de variations (tkz-tab)
\providecommand{\barre}{\ensuremath{\|}}
% environnement proof (amsthm non chargé) utilisé par les modèles
\ifdefined\proof\else\newenvironment{proof}[1][Démonstration]{\par\noindent\textit{#1.}\ }{\qed\par}\fi
\usepackage{lastpage}
\usepackage{hyperref}
\hypersetup{hidelinks}

% ============================================================
% Palette « Pop » OnBuch+ — mêmes hex que la classe P (plus_ui.dart)
% ============================================================
\definecolor{poporange}{HTML}{F59321}
\definecolor{popdark}{HTML}{E07A0C}
\definecolor{popcream}{HTML}{FFF6E8}
\definecolor{popink}{HTML}{1C1714}
\definecolor{poppurple}{HTML}{7A5AE0}
\definecolor{popgreen}{HTML}{2E9E6A}
\definecolor{poppink}{HTML}{E0457B}
\definecolor{popblue}{HTML}{2F7DE1}
\definecolor{popgold}{HTML}{C98A12}
\definecolor{popmuted}{HTML}{8A7F76}
\definecolor{popline}{HTML}{EADFD2}
\definecolor{popgray}{HTML}{8A7F76}
\colorlet{popgrey}{popgray}
% Alias tolérants (le modèle nomme parfois les couleurs différemment)
\colorlet{popwhite}{white}
\colorlet{popblack}{popink}
\colorlet{popred}{poppink}
\colorlet{popyellow}{popgold}
% Teintes claires (fonds de tableaux/encadrés) — le modèle nomme parfois
% les couleurs avec un suffixe L (ex. popblueL) sans les avoir définies.
\colorlet{poporangeL}{poporange!20}
\colorlet{popdarkL}{popdark!20}
\colorlet{poppurpleL}{poppurple!20}
\colorlet{popgreenL}{popgreen!20}
\colorlet{poppinkL}{poppink!20}
\colorlet{popblueL}{popblue!20}
\colorlet{popgoldL}{popgold!20}
\colorlet{popredL}{popred!20}
\colorlet{popgrayL}{popgray!20}
% Teintes claires pour fonds de tableaux / zones
\colorlet{popblueL}{popblue!10!white}
\colorlet{poporangeL}{poporange!14!white}
\colorlet{popgreenL}{popgreen!12!white}
\colorlet{poppurpleL}{poppurple!12!white}
\colorlet{poppinkL}{poppink!12!white}
\colorlet{popgoldL}{popgold!14!white}

\pagecolor{popcream}

% ============================================================
% Typographie
% ============================================================
\defaultfontfeatures{Ligatures=TeX}
\setmainfont{PlusJakartaSans-Medium.ttf}[Path=../../../fonts/,BoldFont=PlusJakartaSans-Bold.ttf]
% Maths en sans-serif (Fira Math) avec chiffres et lettres droites de Plus
% Jakarta, pour que les formules se fondent dans le texte.
\usepackage[math-style=ISO,bold-style=ISO]{unicode-math}
\setmathfont{FiraMath-Regular.otf}[Path=../../../fonts/]
\setmathfont{PlusJakartaSans-Medium.ttf}[Path=../../../fonts/,range={up/num,up/{Latin,latin},\mathup}]
\usepackage{icomma} % virgule décimale sans espace en mode math
% Caractères mathématiques Unicode absents des polices de texte (Plus Jakarta,
% Archivo Black) : rendus par leur équivalent mathématique, en texte comme en
% formule — sinon ils s'affichent en carré vide (ex. « Arithmétique dans ℤ »).
\usepackage{newunicodechar}
\newunicodechar{ℝ}{\ensuremath{\mathbb{R}}}
\newunicodechar{ℤ}{\ensuremath{\mathbb{Z}}}
\newunicodechar{ℂ}{\ensuremath{\mathbb{C}}}
\newunicodechar{ℕ}{\ensuremath{\mathbb{N}}}
\newunicodechar{ℚ}{\ensuremath{\mathbb{Q}}}
\newunicodechar{𝔻}{\ensuremath{\mathbb{D}}}
\newunicodechar{α}{\ensuremath{\alpha}}
\newunicodechar{β}{\ensuremath{\beta}}
\newunicodechar{γ}{\ensuremath{\gamma}}
\newunicodechar{δ}{\ensuremath{\delta}}
\newunicodechar{ε}{\ensuremath{\varepsilon}}
\newunicodechar{θ}{\ensuremath{\theta}}
\newunicodechar{λ}{\ensuremath{\lambda}}
\newunicodechar{μ}{\ensuremath{\mu}}
\newunicodechar{σ}{\ensuremath{\sigma}}
\newunicodechar{τ}{\ensuremath{\tau}}
\newunicodechar{φ}{\ensuremath{\varphi}}
\newunicodechar{ψ}{\ensuremath{\psi}}
\newunicodechar{ω}{\ensuremath{\omega}}
\newunicodechar{Σ}{\ensuremath{\Sigma}}
% majuscules produites par \MakeUppercase dans les titres (α-aminés -> Α)
\newunicodechar{Α}{\ensuremath{\alpha}}
\newunicodechar{Β}{\ensuremath{\beta}}
\newunicodechar{Γ}{\ensuremath{\Gamma}}
\newunicodechar{Δ}{\ensuremath{\Delta}}
\newunicodechar{Ε}{\ensuremath{\varepsilon}}
\newunicodechar{Θ}{\ensuremath{\Theta}}
\newunicodechar{Λ}{\ensuremath{\Lambda}}
\newunicodechar{Μ}{\ensuremath{\mu}}
\newunicodechar{Τ}{\ensuremath{\tau}}
\newunicodechar{Φ}{\ensuremath{\Phi}}
\newunicodechar{Ψ}{\ensuremath{\Psi}}
\newunicodechar{Ω}{\ensuremath{\Omega}}
\newunicodechar{↦}{\ensuremath{\mapsto}}
\newunicodechar{∈}{\ensuremath{\in}}
\newunicodechar{∉}{\ensuremath{\notin}}
\newunicodechar{∧}{\ensuremath{\wedge}}
\newunicodechar{⇔}{\ensuremath{\Leftrightarrow}}
\newunicodechar{⟺}{\ensuremath{\Longleftrightarrow}}
\newunicodechar{⇒}{\ensuremath{\Rightarrow}}
\newunicodechar{⟹}{\ensuremath{\Longrightarrow}}
\newunicodechar{∘}{\ensuremath{\circ}}
\newunicodechar{∪}{\ensuremath{\cup}}
\newunicodechar{∩}{\ensuremath{\cap}}
\newunicodechar{≡}{\ensuremath{\equiv}}
\newunicodechar{⊂}{\ensuremath{\subset}}
\newunicodechar{⊥}{\ensuremath{\perp}}
\newunicodechar{∃}{\ensuremath{\exists}}
\newunicodechar{∀}{\ensuremath{\forall}}
\newunicodechar{∛}{\ensuremath{\sqrt[3]{\,}}}
\newunicodechar{∜}{\ensuremath{\sqrt[4]{\,}}}
\newunicodechar{⃗}{\ensuremath{{}^{\to}}}
\newunicodechar{ⁿ}{\ifmmode^{n}\else\textsuperscript{\ensuremath{n}}\fi}
\newunicodechar{ˣ}{\ifmmode^{x}\else\textsuperscript{\ensuremath{x}}\fi}
\newunicodechar{ᵇ}{\ifmmode^{b}\else\textsuperscript{\ensuremath{b}}\fi}
\newunicodechar{ʳ}{\ifmmode^{r}\else\textsuperscript{\ensuremath{r}}\fi}
\newunicodechar{ᵘ}{\ifmmode^{u}\else\textsuperscript{\ensuremath{u}}\fi}
\newunicodechar{ʸ}{\ifmmode^{y}\else\textsuperscript{\ensuremath{y}}\fi}
\newunicodechar{ᵃ}{\ifmmode^{a}\else\textsuperscript{\ensuremath{a}}\fi}
\newunicodechar{ᵉ}{\ifmmode^{e}\else\textsuperscript{\ensuremath{e}}\fi}
\newunicodechar{ᵏ}{\ifmmode^{k}\else\textsuperscript{\ensuremath{k}}\fi}
\newunicodechar{ᵖ}{\ifmmode^{p}\else\textsuperscript{\ensuremath{p}}\fi}
\newunicodechar{ᴺ}{\ifmmode^{N}\else\textsuperscript{\ensuremath{N}}\fi}
\newunicodechar{ᶜ}{\ifmmode^{c}\else\textsuperscript{\ensuremath{c}}\fi}
\newunicodechar{ʰ}{\ifmmode^{h}\else\textsuperscript{\ensuremath{h}}\fi}
\newunicodechar{ᵠ}{\ifmmode^{q}\else\textsuperscript{\ensuremath{q}}\fi}
\newunicodechar{ₙ}{\ifmmode_{n}\else\textsubscript{\ensuremath{n}}\fi}
\newunicodechar{ₐ}{\ifmmode_{a}\else\textsubscript{\ensuremath{a}}\fi}
\newunicodechar{ᵢ}{\ifmmode_{i}\else\textsubscript{\ensuremath{i}}\fi}
\newunicodechar{ᵣ}{\ifmmode_{r}\else\textsubscript{\ensuremath{r}}\fi}
\newunicodechar{ₖ}{\ifmmode_{k}\else\textsubscript{\ensuremath{k}}\fi}
\newunicodechar{ᵦ}{\ifmmode_{\beta}\else\textsubscript{\ensuremath{\beta}}\fi}
\newunicodechar{ₓ}{\ifmmode_{x}\else\textsubscript{\ensuremath{x}}\fi}
\newfontfamily\popdisplay{ArchivoBlack-Regular.ttf}[Path=../../../fonts/]
\newfontfamily\poplogo{SpaceGrotesk-Bold.ttf}[Path=../../../fonts/]
\newfontfamily\popsemi{PlusJakartaSans-SemiBold.ttf}[Path=../../../fonts/]
\newfontfamily\popxbold{PlusJakartaSans-ExtraBold.ttf}[Path=../../../fonts/]
\newfontfamily\popxboldls{PlusJakartaSans-ExtraBold.ttf}[Path=../../../fonts/,LetterSpace=3.0]

\setlist[enumerate]{leftmargin=1.7em,itemsep=5pt,topsep=4pt}
\setlist[itemize]{leftmargin=1.7em,itemsep=4pt,topsep=4pt,label={\color{poporange}\textbullet}}
\setlength{\parindent}{0pt}
\setlength{\parskip}{5pt}
\renewcommand{\arraystretch}{1.25}

% pgfplots : style Pop par défaut
\pgfplotsset{
  every axis/.append style={axis line style={popink,line width=1pt},tick style={popink},
    grid style={popline,line width=0.5pt},label style={font=\small},tick label style={font=\footnotesize}},
  popcurve/.style={poporange,line width=1.8pt,no markers,smooth},
}
% Même style disponible dans un \draw TikZ simple (hors pgfplots)
\tikzset{popcurve/.style={poporange,line width=1.8pt}}
\tikzset{popgrid/.style={popline,very thin}}
\colorlet{popcurve}{poporange} % le nom du style est parfois utilisé comme couleur

% ============================================================
% Pill / squiggle
% ============================================================
\newcommand{\poppill}[3]{%
  \tikz[baseline=-0.55ex]\node[draw=popink,line width=1.2pt,rounded corners=2.6pt,%
    fill=#2,inner xsep=6.5pt,inner ysep=3pt]{%
    \popxboldls\fontsize{7}{7}\selectfont\color{#3}\MakeUppercase{#1}};%
}
\newcommand{\popsquiggle}[1]{%
  \tikz[baseline=-0.4ex]{
    \draw[#1,line width=2.6pt,line cap=round]
      plot [smooth] coordinates {(0,0.09) (0.34,0.01) (0.68,0.21) (1.02,0.01) (1.36,0.10)};
  }%
}
% Mot-clé surligné (marqueur orange sous le texte)
\newcommand{\cle}[1]{{\popxbold\color{popdark}#1}}

% ============================================================
% En-tête / pied
% ============================================================
\pagestyle{fancy}
\fancyhf{}
\renewcommand{\headrulewidth}{0pt}
\renewcommand{\footrulewidth}{0pt}
\lhead{\raisebox{-0.15ex}{\poplogo\fontsize{13}{13}\selectfont\color{popink}OnBuch\color{poporange}+}}
\rhead{\poppill{\DOCMATIERE\ \textbullet\ \DOCNIVEAU\ \textbullet\ Leçon \DOCLECON}{white}{popink}}
\lfoot{\popxbold\fontsize{7.6}{7.6}\selectfont\color{popmuted}\MakeUppercase{Cours signé OnBuch+}\ \textbullet\ {\popsemi\DOCTITRE}}
\rfoot{\tikz[baseline=-0.6ex]\node[rounded corners=8.5pt,fill=popink,minimum height=17pt,minimum width=17pt,inner xsep=5pt,inner sep=0]{\popxbold\fontsize{8}{8}\selectfont\color{popcream}\thepage\,/\,\pageref*{LastPage}};}

% ============================================================
% Titres
% ============================================================
\newcommand{\chaptitre}[2]{%
  \begin{tcolorbox}[enhanced,colback=poporange,colframe=popink,boxrule=2.6pt,arc=11pt,
      drop shadow=popink,left=15pt,right=15pt,top=12pt,bottom=11pt,before skip=3pt,after skip=18pt]
    \poppill{#2}{popink}{popcream}\par\vspace{9pt}
    {\raggedright\hyphenpenalty=10000\exhyphenpenalty=10000\popdisplay\fontsize{22}{24}\selectfont\color{popcream}\MakeUppercase{#1}\par}
  \end{tcolorbox}
}
% Section numérotée : pastille orange avec le numéro + titre Archivo
\newcounter{popsec}
\newcommand{\coursec}[1]{%
  \stepcounter{popsec}\setcounter{popsub}{0}%
  \par\vspace{16pt}\noindent
  \begin{minipage}{\linewidth}
  \tikz[baseline=-0.7ex]\node[draw=popink,line width=1.6pt,fill=poporange,rounded corners=4pt,minimum size=22pt,inner sep=0]{\popdisplay\fontsize{12}{12}\selectfont\color{popcream}\thepopsec};%
  \hspace{8pt}{\popdisplay\fontsize{15}{17}\selectfont\color{popink}\MakeUppercase{#1}}\par
  \vspace{4pt}\noindent{\color{poporange}\rule{36pt}{3.6pt}}
  \end{minipage}\par\nopagebreak\vspace{8pt}
}
\newcounter{popsub}
\newcommand{\courssub}[1]{%
  \stepcounter{popsub}%
  \par\vspace{9pt}\noindent{\popxbold\fontsize{11.5}{13}\selectfont\color{popink}{\color{poporange}\thepopsec.\thepopsub}\hspace{6pt}#1}\par\nopagebreak\vspace{4pt}
}

% ============================================================
% Boîtes pédagogiques — même recette : fond blanc, bordure épaisse colorée,
% ombre dure encre, badge plein. Titre optionnel [#1].
% ============================================================
\tcbset{popbase/.style={breakable,enhanced,colback=white,boxrule=2.3pt,arc=9pt,
  shadow={3pt}{-3pt}{0pt}{popink},left=14pt,right=14pt,top=11pt,bottom=11pt,before skip=11pt,after skip=15pt,
  fontupper=\popsemi\fontsize{10.6}{14.5}\selectfont\color{popink},
  fontlower=\popsemi\fontsize{10.6}{14.5}\selectfont\color{popink}}}
% Coche dessinée (le glyphe ✓ n'existe pas dans les polices du document)
\newcommand{\popcheck}[1]{\tikz[baseline=-0.5ex]\draw[#1,line width=1.6pt,line cap=round,line join=round](-0.13,0.0)--(-0.04,-0.09)--(0.13,0.1);}
\providecommand{\coche}{\popcheck{popgreen}}
\providecommand{\resultat}[1]{\par\smallskip\noindent\fcolorbox{popgreen}{popgreenL}{\parbox{\dimexpr\linewidth-2\fboxsep-2\fboxrule\relax}{\textbf{Résultat.} #1}}\par\smallskip}
\newcommand{\popboxhead}[3]{\poppill{#1}{#2}{white}\ifx\relax#3\relax\else\hspace{6pt}{\popxbold\fontsize{10.6}{12}\selectfont\color{popink}#3}\fi\par\vspace{7pt}}

\newtcolorbox{aretenir}[1][]{popbase,colframe=popgreen,before upper={\popboxhead{À retenir}{popgreen}{#1}}}
% Texte en espagnol (document de lecture, dialogue, poème, extrait) : cadre
% d'étude. Un titre optionnel (source, auteur) s'affiche dans l'en-tête.
\newtcolorbox{textebox}[1][]{popbase,colframe=popblue,before upper={\popboxhead{Texto}{popblue}{#1}\begin{otherlanguage}{spanish}},after upper={\end{otherlanguage}}}
% Marques ✓ / ✗ (forme correcte / fautive) souvent utilisées par les modèles
\providecommand{\cross}{\textcolor{poppink}{\ensuremath{\times}}}
\providecommand{\xmark}{\cross}\providecommand{\crossmark}{\cross}\providecommand{\cmark}{\ensuremath{\checkmark}}
% Ligne de réponse / justification à compléter (inventée par les modèles)
\providecommand{\justification}[1]{\par\noindent\textit{Justification~:} \dotfill\par}
% \textipa (package tipa non chargé) : la police ne couvre pas l'API -> simple passage
\providecommand{\textipa}[1]{#1}
% Soulignement (ulem non chargé) : \uline, \ul
\providecommand{\uline}[1]{\underline{#1}}\providecommand{\ul}[1]{\underline{#1}}
% Mot ou phrase en espagnol à l'intérieur d'un paragraphe français
\newcommand{\esp}[1]{\ifmmode\text{\textit{#1}}\else\begingroup\selectlanguage{spanish}\itshape #1\endgroup\fi}
\newtcolorbox{exemplebox}[1][]{popbase,colframe=poporange,before upper={\popboxhead{Exemple}{poporange}{#1}}}
\newtcolorbox{definition}[1][]{popbase,colframe=popblue,before upper={\popboxhead{Définition}{popblue}{#1}}}
\newtcolorbox{propriete}[1][]{popbase,colframe=poppurple,before upper={\popboxhead{Propriété}{poppurple}{#1}}}
\newtcolorbox{methode}[1][]{popbase,colframe=popgold,before upper={\popboxhead{Méthode}{popgold}{#1}}}
\newtcolorbox{attention}[1][]{popbase,colframe=poppink,before upper={\popboxhead{Attention}{poppink}{#1}}}
\newtcolorbox{experience}[1][]{popbase,colframe=popblue,colback=popblueL,before upper={\popboxhead{Expérience}{popblue}{#1}}}
% « Le savais-tu ? » : carte violette pleine (contexte, culture, Cameroun)
\newtcolorbox{savaistu}[1][]{popbase,colback=poppurple,colframe=popink,
  fontupper=\popsemi\fontsize{10.4}{14.2}\selectfont\color{white},
  before upper={\poppill{Le savais-tu\,?}{white}{poppurple}\ifx\relax#1\relax\else\hspace{6pt}{\popxbold\color{white}#1}\fi\par\vspace{7pt}}}
% Exercice résolu : énoncé en haut, \tcblower puis la solution sur fond crème
\newcounter{popexo}\newcounter{popexr}
\newtcolorbox{exoresolu}[1][]{popbase,colframe=poporange,colbacklower=popcream,
  segmentation style={popink,line width=1.2pt,dashed},
  before upper={\stepcounter{popexr}\popboxhead{Exercice résolu \thepopexr}{poporange}{#1}},
  before lower={\poppill{Solution}{popink}{popcream}\par\vspace{6pt}}}
% Exercice (non résolu en place) — la correction est dans la partie Corrigés
\newtcolorbox{exercice}[1][]{popbase,colframe=popink,
  before upper={\stepcounter{popexo}\popboxhead{Exercice \thepopexo}{popink}{#1}}}
% Corrigé
\newtcolorbox{corrige}[1][]{popbase,colframe=popgreen,colback=popgreenL,
  before upper={\popboxhead{Corrigé}{popgreen}{#1}}}
% Objectifs / prérequis / bilan
\newtcolorbox{objectifs}{popbase,colframe=popink,colback=poporangeL,
  before upper={\popboxhead{Objectifs de la leçon}{popink}{}}}
\newtcolorbox{prerequis}{popbase,colframe=popink,colback=popblueL,
  before upper={\popboxhead{Prérequis}{popblue}{}}}
\newtcolorbox{bilan}{popbase,colframe=popink,colback=white,boxrule=2.8pt,
  before upper={\popboxhead{Fiche bilan}{poporange}{L'essentiel à connaître pour l'examen}}}
% Figure encadrée avec légende
\newtcolorbox{popfigure}[1][]{enhanced,breakable=false,colback=white,colframe=popline,boxrule=1.4pt,arc=8pt,
  left=10pt,right=10pt,top=10pt,bottom=8pt,before skip=10pt,after skip=12pt,halign=center,
  lower separated=false,halign lower=center,
  fontlower=\popsemi\fontsize{9}{11.5}\selectfont\color{popmuted},
  #1}
\newcommand{\legende}[1]{\par\vspace{6pt}{\popsemi\fontsize{9}{11.5}\selectfont\color{popmuted}#1}}

% ============================================================
% Signature OnBuch+
% ============================================================
\newcommand{\poplogoinline}[1]{{\poplogo\fontsize{#1}{#1}\selectfont\color{popink}OnBuch\color{poporange}+}}
\newcommand{\signatureonbuch}{%
  \begin{tcolorbox}[enhanced,colback=popink,colframe=popink,boxrule=2.2pt,arc=10pt,
      drop shadow=poporange,left=14pt,right=14pt,top=10pt,bottom=10pt,before skip=0pt,after skip=0pt]
    \tikz[baseline=-0.7ex]\node[circle,fill=popgreen,draw=popcream,line width=1.3pt,minimum size=22pt,inner sep=0]{\popcheck{white}};%
    \hspace{9pt}\begin{minipage}[c]{0.62\linewidth}
      {\popxbold\fontsize{12}{13}\selectfont\color{popcream}Signé\ \textbullet\ {\poplogo OnBuch\color{poporange}+}}\par\vspace{2pt}
      {\raggedright\popsemi\fontsize{8}{10.5}\selectfont\color{popline}\MakeUppercase{Équipe pédagogique OnBuch+}\\\MakeUppercase{Conforme au programme officiel MINESEC}\par}
    \end{minipage}\hfill
    {\poplogo\fontsize{22}{22}\selectfont\color{popcream}OnBuch\color{poporange}+}
  \end{tcolorbox}%
}

% ============================================================
% Page de garde
% #1 titre · #2 matière · #3 niveau · #4 module · #5 n° leçon · #6 durée ·
% #7 description / objectifs (texte ou itemize)
% ============================================================
\newcommand{\pagegarde}[7]{%
  \thispagestyle{empty}
  \newgeometry{a4paper,margin=1.7cm,top=1.6cm,bottom=1.4cm}%
  \begin{tikzpicture}[remember picture,overlay]
    \fill[poporange!18] ([xshift=-3.2cm,yshift=-3.6cm]current page.north east) circle (4.6cm);
    \fill[poppurple!14] ([xshift=2.4cm,yshift=7.6cm]current page.south west) circle (3.8cm);
  \end{tikzpicture}%
  {\poplogo\fontsize{30}{30}\selectfont\color{popink}OnBuch\color{poporange}+}\hfill
  \raisebox{0.6ex}{\poppill{Cours signé \textbullet\ Premium}{popink}{popcream}}\par
  \vspace{1pt}\popsquiggle{poppurple}\par
  \vspace{8pt}
  \begin{tcolorbox}[enhanced,colback=poporange,colframe=popink,boxrule=2.8pt,arc=13pt,
      drop shadow=popink,left=18pt,right=18pt,top=12pt,bottom=13pt,after skip=10pt]
    \poppill{#2 \textbullet\ #3}{popink}{popcream}\hspace{5pt}\poppill{#4}{white}{popink}\par\vspace{8pt}
    {\popdisplay\fontsize{14}{14}\selectfont\color{popink}LEÇON #5}\par\vspace{4pt}
    {\raggedright\hyphenpenalty=10000\exhyphenpenalty=10000\popdisplay\fontsize{25}{27}\selectfont\color{popcream}\MakeUppercase{#1}\par}
    \vspace{8pt}
    {\popxbold\fontsize{9.5}{9.5}\selectfont\color{popink}\MakeUppercase{Durée conseillée : #6}}
  \end{tcolorbox}
  \begin{tcolorbox}[enhanced,colback=white,colframe=popink,boxrule=2.2pt,arc=10pt,
      drop shadow=popink,left=16pt,right=16pt,top=10pt,bottom=10pt,after skip=8pt]
    \poppill{Dans cette leçon}{poppurple}{white}\par\vspace{5pt}
    \popsemi\fontsize{9.3}{12.6}\selectfont\color{popink}#7
  \end{tcolorbox}
  \noindent\begin{minipage}[t]{0.487\linewidth}
    \begin{tcolorbox}[enhanced,colback=popgreen,colframe=popink,boxrule=2.2pt,arc=10pt,
        drop shadow=popink,left=11pt,right=11pt,top=10pt,bottom=10pt,height=3.1cm,valign=center]
      \tcbox[colback=white,colframe=popink,boxrule=1.3pt,arc=5pt,boxsep=0pt,nobeforeafter,
        left=3pt,right=3pt,top=3pt,bottom=3pt,valign=center]{\qrcode[height=1.9cm]{https://play.google.com/store/apps/details?id=com.luvvix.onbuch&hl=fr}}%
      \hspace{8pt}\begin{minipage}[c]{0.5\linewidth}
        {\popdisplay\fontsize{11}{12.5}\selectfont\color{white}\MakeUppercase{Télécharge OnBuch}}\par\vspace{4pt}
        {\raggedright\popsemi\fontsize{8.4}{11}\selectfont\color{white}Gratuit \textbullet\ Google Play\\Tuteur IA, annales, concours, résultats\par}
      \end{minipage}
    \end{tcolorbox}
  \end{minipage}\hfill
  \begin{minipage}[t]{0.487\linewidth}
    \begin{tcolorbox}[enhanced,colback=poppurple,colframe=popink,boxrule=2.2pt,arc=10pt,
        drop shadow=popink,left=11pt,right=11pt,top=10pt,bottom=10pt,height=3.1cm,valign=center]
      \tikz[baseline=-0.7ex]\node[circle,fill=white,draw=popink,line width=1.3pt,minimum size=1.6cm,inner sep=0]{\popdisplay\fontsize{24}{24}\selectfont\color{poppurple}+};%
      \hspace{8pt}\begin{minipage}[c]{0.6\linewidth}
        {\popdisplay\fontsize{11}{12.5}\selectfont\color{white}\MakeUppercase{Passe à OnBuch+}}\par\vspace{4pt}
        {\raggedright\popsemi\fontsize{8.4}{11}\selectfont\color{white}Tous les cours signés, Tuteur IA illimité, fiches PDF — onglet OnBuch+ de l'app\par}
      \end{minipage}
    \end{tcolorbox}
  \end{minipage}
  \vspace{8pt}
  \popcontenu
  \vfill
  \signatureonbuch
  \restoregeometry
  \newpage
}

% Rangée « Ce cours contient » (page de garde)
\newcommand{\poptile}[3]{% #1 couleur #2 gros texte #3 légende
  \begin{tcolorbox}[enhanced,colback=#1,colframe=popink,boxrule=2pt,arc=9pt,shadow={2.5pt}{-2.5pt}{0pt}{popink},
      left=6pt,right=6pt,top=9pt,bottom=9pt,halign=center,valign=center,height=1.75cm,width=\linewidth,nobeforeafter]
    {\popdisplay\fontsize{12.5}{13}\selectfont\color{white}\MakeUppercase{#2}}\par\vspace{3pt}
    {\centering\popsemi\fontsize{7.8}{9.5}\selectfont\color{white}#3\par}
  \end{tcolorbox}}
\newcommand{\popcontenu}{%
  \par\noindent{\popxboldls\fontsize{7.5}{7.5}\selectfont\color{popmuted}CE COURS SIGNÉ CONTIENT}\par\vspace{7pt}
  \noindent\begin{minipage}[t]{0.235\linewidth}\poptile{popblue}{Cours}{complet, illustré, pas à pas}\end{minipage}\hfill
  \begin{minipage}[t]{0.235\linewidth}\poptile{poporange}{Méthodes}{et exercices résolus}\end{minipage}\hfill
  \begin{minipage}[t]{0.235\linewidth}\poptile{poppink}{Exercices}{type examen + corrigés}\end{minipage}\hfill
  \begin{minipage}[t]{0.235\linewidth}\poptile{popgreen}{Fiche bilan}{l'essentiel à retenir}\end{minipage}\par}

% Sommaire
\newtcolorbox{sommaire}{enhanced,colback=white,colframe=popink,boxrule=2.2pt,arc=10pt,
  drop shadow=popink,left=18pt,right=18pt,top=13pt,bottom=13pt,before skip=6pt,after skip=20pt,
  before upper={\poppill{Sommaire}{poppurple}{white}\par\vspace{9pt}\popsemi\fontsize{10.6}{14.5}\selectfont\color{popink}}}

% Signature de fin de cours — poussée tout en bas de la dernière page
\newcommand{\findecours}{%
  \par\vfill\nopagebreak
  \begin{center}\popsquiggle{poporange}\end{center}\vspace{4pt}
  \signatureonbuch
}
\AtBeginDocument{\def\llbracket{\mathopen{[\mkern-3.5mu[}}\def\rrbracket{\mathclose{]\mkern-3.5mu]}}} % intervalles d'entiers (glyphe ⟦ absent de la police)
% Glyphes absents de Fira Math : redéfinis à partir de glyphes disponibles
\AtBeginDocument{%
  \def\setminus{\mathbin{\backslash}}\let\smallsetminus\setminus
  \def\vdots{\mathord{\vbox{\baselineskip3.2pt\lineskiplimit0pt\kern4pt\hbox{.}\hbox{.}\hbox{.}}}}%
  \def\ddots{\mathinner{\mkern1mu\raise6.5pt\hbox{.}\mkern2mu\raise3.6pt\hbox{.}\mkern2mu\raise0.7pt\hbox{.}\mkern1mu}}%
  \def\triangle{\mathord{\scalebox{0.9}{$\symup{\Delta}$}}}%
  \def\nabla{\mathord{\raisebox{\depth}{\scalebox{1}[-1]{$\symup{\Delta}$}}}}%
  \def\checkmark{\popcheck{popdark}}%
  \def\star{\mathbin{\ast}}\def\bigstar{\mathord{\ast}}%
  \def\diamond{\mathbin{\scalebox{0.6}{$\Diamond$}}}%
  \def\bigcup{\mathop{\mathchoice{\raisebox{-0.3em}{\scalebox{1.7}{$\cup$}}}{\scalebox{1.25}{$\cup$}}{\cup}{\cup}}\displaylimits}%
  \def\bigcap{\mathop{\mathchoice{\raisebox{-0.3em}{\scalebox{1.7}{$\cap$}}}{\scalebox{1.25}{$\cap$}}{\cap}{\cap}}\displaylimits}%
  \def\lesseqqgtr{\mathrel{\lessgtr}}\def\bigoplus{\mathop{\oplus}\displaylimits}%
}
\providecommand{\ssi}{si et seulement si} % abréviation fréquente
\AtBeginDocument{\providecommand{\wideparen}{\overparen}} % arc de cercle

%% FIGFIX-BEGIN v5 — correctifs globaux des figures (docs/onbuch-generation/tools/figfix)
\usepackage{adjustbox}\usepackage{etoolbox}\tcbuselibrary{raster}
% toute figure TikZ/pgfplots plus large que la ligne est réduite à la largeur disponible
\BeforeBeginEnvironment{tikzpicture}{\begin{adjustbox}{max width=\linewidth}}
\AfterEndEnvironment{tikzpicture}{\end{adjustbox}}
% légendes pgfplots lisibles par-dessus les courbes
\pgfplotsset{every axis/.append style={legend style={fill=white,fill opacity=0.92,text opacity=1,draw=black!15,inner sep=3pt}}}
% étiquettes d'arêtes (syntaxe edge["..."]) avec fond blanc
\tikzset{every edge quotes/.append style={fill=white,inner sep=1.2pt}}
%% FIGFIX-END
\begin{document}

\pagegarde{Lectura y comentario : el cuarto poder y la promoción de la cultura de paz}{Espagnol}{Tle A}{Módulo 5 : Médias, communication et culture de la paix}{E24}{2 h}{Cette leçon de lecture-commentaire propose un texte argumentatif sur le quatrième pouvoir et son rôle dans la construction de la paix. Elle permet d'acquérir le lexique de la presse et de la médiation, et de s'entraîner à la méthode du commentaire de texte exigée au Baccalauréat.
\vspace{4pt}
\begin{multicols}{2}\small
\begin{itemize}
\item À la fin de cette leçon, je sais lire et comprendre un texte argumentatif ou explicatif sur le quatrième pouvoir.
\item Je sais dégager l'idée générale et la thèse d'un texte sur la presse.
\item Je sais employer correctement le lexique du quatrième pouvoir : cuarto poder, libertad de prensa, independencia, censura, opinión pública.
\item Je sais employer le lexique de la paix : mediador, reconciliación, cultura de paz, diálogo.
\item Je sais repérer les procédés argumentatifs d'un texte (exemples, chiffres, citations, connecteurs).
\item Je sais rédiger un commentaire de texte structuré (introduction, développement, conclusion).
\end{itemize}\end{multicols}}

\begin{sommaire}
\begin{multicols}{2}
\begin{enumerate}
\item Texte support : El cuarto poder al servicio de la paz
\item Compréhension du texte
\item Lexique thématique par champs
\item Méthode du commentaire de texte
\item Commentaire modèle
\item Production guidée : défendre la culture de la paix
\item Activité d'intégration
\item Méthodes et exercices résolus
\item Exercices
\item Corrigés des exercices
\item Fiche bilan
\end{enumerate}
\end{multicols}
\end{sommaire}

\begin{objectifs}
\begin{itemize}
\item À la fin de cette leçon, je sais lire et comprendre un texte argumentatif ou explicatif sur le quatrième pouvoir.
\item Je sais dégager l'idée générale et la thèse d'un texte sur la presse.
\item Je sais employer correctement le lexique du quatrième pouvoir : cuarto poder, libertad de prensa, independencia, censura, opinión pública.
\item Je sais employer le lexique de la paix : mediador, reconciliación, cultura de paz, diálogo.
\item Je sais repérer les procédés argumentatifs d'un texte (exemples, chiffres, citations, connecteurs).
\item Je sais rédiger un commentaire de texte structuré (introduction, développement, conclusion).
\item Je sais défendre à l'écrit et à l'oral la promotion de la culture de la paix par les médias.
\item Je sais réagir à une thèse en exprimant mon accord ou mon désaccord avec des arguments.
\end{itemize}
\end{objectifs}

\begin{prerequis}
\begin{itemize}
\item Lexique des médias vu en Première (prensa, periódico, periodista, noticia, reportaje).
\item Connecteurs logiques de l'argumentation (sin embargo, por lo tanto, además, en primer lugar).
\item Expression de l'opinion (pienso que, me parece que, estoy de acuerdo con).
\item Subjonctif après les verbes d'opinion à la négative (no creo que + subjuntivo).
\item Notion de culture de la paix introduite dans le module 5.
\end{itemize}
\end{prerequis}

\newpage

\coursec{Texte support : El cuarto poder al servicio de la paz}

\courssub{Mise en situation et problématique}

Imaginez la scène. À Yaoundé, un journaliste de la CRTV est envoyé dans une localité où deux communautés s'affrontent pour un conflit foncier. Au lieu d'attiser la polémique, il choisit de donner la parole aux deux camps : les anciens du premier village, puis ceux du second, les femmes, les jeunes. Son reportage, diffusé le soir même, apaise les tensions : chacun s'est senti écouté. À des milliers de kilomètres, en Colombie, après des décennies de guerre entre l'État et les FARC, les médias ont accompagné le processus de paix en diffusant les témoignages des victimes, des ex-combattants et des familles réconciliées.

Partout dans le monde, on appelle la presse \esp{el cuarto poder} : après les pouvoirs exécutif, législatif et judiciaire, elle informe, dénonce et influence \esp{la opinión pública}. Mais que se passe-t-il quand la censure la réduit au silence ? Le journaliste peut-il devenir un \esp{mediador} dans les conflits ? Cette leçon nous invite à lire et à commenter un texte argumentatif sur le rôle des médias dans la promotion de la culture de la paix.

\begin{aretenir}[Problématique de la leçon]
\esp{¿Puede la prensa, como cuarto poder, contribuir a la cultura de paz y a la reconciliación entre los pueblos?} — La presse, en tant que quatrième pouvoir, peut-elle contribuer à la culture de la paix et à la réconciliation entre les peuples ?
\end{aretenir}

\courssub{Lecture globale du texte}

\begin{methode}[Comment aborder la première lecture]
\begin{enumerate}
\item \textbf{Lecture silencieuse} (3 minutes) : lis le texte en entier sans dictionnaire. Souligne les mots inconnus, mais ne t'arrête pas dessus.
\item \textbf{Lecture à voix haute par l'enseignant} : écoute la prononciation, l'intonation et les pauses. Note le rythme des phrases argumentatives.
\item \textbf{Première impression} : de quel type de texte s'agit-il ? Qui parle ? À qui ? Pour dire quoi ?
\item \textbf{Deuxième lecture} : avec le glossaire, repère la thèse, les arguments et les exemples.
\end{enumerate}
\end{methode}

\begin{textebox}[El cuarto poder al servicio de la paz]
Se llama \textbf{cuarto poder} a la prensa porque, después de los tres poderes clásicos —ejecutivo, legislativo y judicial—, los medios de comunicación ejercen una influencia decisiva sobre la opinión pública. La expresión nació en el siglo XIX, cuando los periódicos empezaron a vigilar a los gobiernos y a exigirles cuentas. Desde entonces, la prensa cumple tres funciones esenciales: \textbf{informar}, \textbf{denunciar} y \textbf{educar}.

Informar es la primera misión del periodista: contar los hechos con veracidad, sin mentiras ni exageraciones. Un ciudadano bien informado puede votar, decidir y participar en la vida de su país. Denunciar es la segunda función: la pluma del periodista señala las injusticias, la corrupción y los abusos de poder. Gracias a las investigaciones de la prensa, muchos escándalos han salido a la luz. Por último, educar significa formar el espíritu crítico de los lectores y enseñarles valores como la tolerancia, el diálogo y el respeto del otro.

Sin embargo, esta libertad no existe en todas partes. Cuando la \textbf{censura} silencia a los periodistas, cuando los periódicos son cerrados y las radios intervenidas, el pueblo queda ciego y sordo: ya no sabe lo que ocurre en su propio país. La opinión pública, privada de información veraz, se convierte en juguete de los rumores y de la propaganda. Por eso la libertad de prensa es un tesoro que hay que defender cada día.

En los conflictos, además, el periodista puede convertirse en \textbf{mediador}. Dando la palabra a las dos partes enfrentadas, difundiendo los testimonios de las víctimas y rechazando los discursos de odio, favorece el diálogo y la reconciliación. En África y en América Latina, numerosas radios comunitarias trabajan ya por la paz: en Guatemala, por ejemplo, las radios en lenguas mayas ayudaron a curar las heridas de una larga guerra civil; en Colombia, los programas dedicados a las víctimas acompañaron el proceso de paz.

Promover una cultura de paz es, pues, una responsabilidad compartida: los medios deben informar con honestidad, y los ciudadanos deben leer, escuchar y pensar con espíritu crítico. Solo así el cuarto poder será, de verdad, un poder al servicio de la vida y de la reconciliación entre los pueblos.
\end{textebox}

\textbf{Glossaire (mots difficiles du texte)}

\begin{center}
\begin{tabularx}{\linewidth}{|X|X|}
\hline
\rowcolor{popblueL} \textbf{Español} & \textbf{Français} \\
\hline
la censura & la censure \\
\hline
la pluma & la plume (métaphore : l'écriture du journaliste) \\
\hline
exigir cuentas a alguien & demander des comptes à quelqu'un \\
\hline
salir a la luz & être révélé, éclater au grand jour \\
\hline
veraz & véridique, exact \\
\hline
intervenir (una radio) & mettre sous tutelle / séquestrer (une radio) \\
\hline
el juguete & le jouet (ici : la proie, la chose manipulée) \\
\hline
la veracidad & la véracité, l'exactitude \\
\hline
\end{tabularx}
\end{center}

\begin{attention}[Prononciation et faux amis]
\begin{itemize}
\item \esp{veraz} [be-\cle{raθ}] se prononce « bé-rasse » (le \esp{v} espagnol se prononce comme un \esp{b} bilabial ; \cle{z} = \cle{θ} en Espagne, \cle{s} en Amérique).
\item \esp{ejercer} [e-\cle{xer}-\cle{ser}] se prononce « é-hèr-cèr » : le \esp{j} et le \esp{g} devant \esp{e/i} ont le son \cle{x} (comme un \emph{r} grasseyé fort), jamais le \emph{j} français.
\item Faux ami : \esp{la audiencia} ne signifie pas « l'audition » (médical), mais le public, l'auditoire, l'audience (médias).
\item Faux ami : \esp{exigir} = exiger, réclamer ; attention à l'orthographe : un seul \esp{g} (\esp{exigir}, \esp{exijo}, \esp{exigió}).
\item \esp{la imprenta} = l'imprimerie (la machine/le lieu) ; \esp{la prensa} = la presse (ensemble des journaux, des journalistes). Ne les confondez pas.
\end{itemize}
\end{attention}

\courssub{Repérage du genre du texte}

Avant de comprendre le texte en détail, il faut identifier sa nature. Observons les indices matériels et linguistiques.

\begin{center}
\begin{tabularx}{\linewidth}{|X|X|X|}
\hline
\rowcolor{popblueL} \textbf{Critère} & \textbf{Observation} & \textbf{Conclusion} \\
\hline
Source probable & vocabulaire journalistique, ton engagé & un article de presse (rubrique opinion/tribune) \\
\hline
Genre & thèse + arguments + exemples & texte argumentatif (article d'opinion) \\
\hline
Destinataire & \esp{los ciudadanos}, le grand public & le lecteur ordinaire \\
\hline
Ton & sérieux, didactique, parfois solennel (\esp{un tesoro que hay que defender}) & ton engagé en faveur de la liberté \\
\hline
Intention & convaincre que la presse sert la paix & persuader et sensibiliser \\
\hline
\end{tabularx}
\end{center}

\begin{definition}[Cuarto poder]
\esp{El cuarto poder} (« le quatrième pouvoir ») est une expression attribuée à l'homme politique et écrivain britannique Edmund Burke (XVIII\textsuperscript{e} siècle), qui aurait désigné ainsi la galerie des journalistes au Parlement. Elle désigne la presse comme un quatrième pouvoir de contrôle dans une démocratie, à côté des trois pouvoirs classiques définis par Montesquieu : le pouvoir \esp{ejecutivo} (exécutif), \esp{legislativo} (législatif) et \esp{judicial} (judiciaire). La presse ne gouverne pas, mais elle surveille ceux qui gouvernent et informe l'opinion publique.
\end{definition}

\begin{exemplebox}[L'expression en contexte]
\begin{itemize}
\item \esp{La prensa es el cuarto poder porque controla a los que gobiernan.} — La presse est le quatrième pouvoir parce qu'elle contrôle ceux qui gouvernent.
\item \esp{Sin libertad de prensa, el cuarto poder desaparece y la democracia se debilita.} — Sans liberté de la presse, le quatrième pouvoir disparaît et la démocratie s'affaiblit.
\item \esp{Los periodistas de investigación son los guardianes del cuarto poder.} — Les journalistes d'investigation sont les gardiens du quatrième pouvoir.
\end{itemize}
\end{exemplebox}

\courssub{Les trois fonctions de la presse : carte mentale}

Le texte s'organise autour d'une idée centrale : la presse, en tant que \esp{cuarto poder}, remplit trois fonctions essentielles. La carte mentale ci-dessous résume la structure argumentative du texte.

\begin{popfigure}
\begin{center}\tcbox[colback=popblue,colframe=popblue,coltext=white,arc=9pt,boxsep=4pt,left=6pt,right=6pt]{\bfseries\small el cuarto poder}\end{center}
\vspace{-2pt}
\begin{tcbraster}[raster columns=2,raster equal height=rows,raster column skip=6pt,raster row skip=6pt,enhanced,blankest]
\begin{tcolorbox}[enhanced,colback=white,colframe=poporange,boxrule=0.9pt,arc=3pt,title={informar \emph{contar los hechos}},fonttitle=\bfseries\scriptsize,coltitle=white,colbacktitle=poporange,left=4pt,right=4pt,top=2pt,bottom=2pt,boxsep=1pt,toptitle=1.5pt,bottomtitle=1.5pt,halign=left]
\begin{itemize}[nosep,leftmargin=*,itemsep=1pt,font=\scriptsize]
\item veracidad
\item ciudadano informado
\end{itemize}
\end{tcolorbox}
\begin{tcolorbox}[enhanced,colback=white,colframe=poppurple,boxrule=0.9pt,arc=3pt,title={denunciar \emph{señalar injusticias}},fonttitle=\bfseries\scriptsize,coltitle=white,colbacktitle=poppurple,left=4pt,right=4pt,top=2pt,bottom=2pt,boxsep=1pt,toptitle=1.5pt,bottomtitle=1.5pt,halign=left]
\begin{itemize}[nosep,leftmargin=*,itemsep=1pt,font=\scriptsize]
\item corrupción
\item abusos de poder
\end{itemize}
\end{tcolorbox}
\begin{tcolorbox}[enhanced,colback=white,colframe=popgreen,boxrule=0.9pt,arc=3pt,title={educar para la paz},fonttitle=\bfseries\scriptsize,coltitle=white,colbacktitle=popgreen,left=4pt,right=4pt,top=2pt,bottom=2pt,boxsep=1pt,toptitle=1.5pt,bottomtitle=1.5pt,halign=left]
\begin{itemize}[nosep,leftmargin=*,itemsep=1pt,font=\scriptsize]
\item espíritu crítico
\item diálogo y reconciliación
\end{itemize}
\end{tcolorbox}
\end{tcbraster}

\legende{Carte mentale : les trois fonctions du cuarto poder selon le texte.}
\end{popfigure}

\begin{propriete}[Structure d'un texte argumentatif]
Un texte argumentatif comme celui-ci suit généralement quatre étapes repérables :
\begin{enumerate}
\item \textbf{La thèse} (\esp{la tesis}) : l'idée défendue. Ici : la presse, quatrième pouvoir, doit servir la paix.
\item \textbf{Les arguments} (\esp{los argumentos}) : les raisons qui soutiennent la thèse. Ici : les trois fonctions (informar, denunciar, educar) et le danger de la censure.
\item \textbf{Les exemples} (\esp{los ejemplos}) : les faits concrets qui illustrent les arguments. Ici : les radios du Guatemala, les médias colombiens.
\item \textbf{La conclusion} (\esp{la conclusión}) : la chute qui ouvre sur une responsabilité. Ici : la culture de la paix est une tâche partagée entre médias et citoyens.
\end{enumerate}
\end{propriete}

\courssub{Premiers éléments de compréhension}

\begin{exoresolu}[Questions de compréhension globale]
\textbf{Énoncé.} Réponds en espagnol, par des phrases complètes :
\begin{enumerate}
\item ¿Por qué se llama « cuarto poder » a la prensa?
\item ¿Cuáles son las tres funciones esenciales de la prensa según el texto?
\item ¿Qué ocurre cuando la censura silencia a los periodistas?
\item ¿Cómo puede el periodista convertirse en mediador?
\end{enumerate}
\tcblower
\textbf{Solution détaillée.}
\begin{enumerate}
\item \esp{Se llama « cuarto poder » a la prensa porque, después de los tres poderes clásicos (ejecutivo, legislativo y judicial), ejerce una influencia decisiva sobre la opinión pública.} — On appelle la presse « quatrième pouvoir » parce que, après les trois pouvoirs classiques, elle exerce une influence décisive sur l'opinion publique.
\item \esp{Las tres funciones esenciales de la prensa son informar, denunciar y educar.} — Les trois fonctions essentielles sont informer, dénoncer et éduquer.
\item \esp{Cuando la censura silencia a los periodistas, el pueblo queda ciego y sordo, y la opinión pública se convierte en juguete de los rumores y de la propaganda.} — Quand la censure réduit les journalistes au silence, le peuple devient aveugle et sourd, et l'opinion publique devient la proie des rumeurs et de la propagande.
\item \esp{El periodista puede convertirse en mediador dando la palabra a las dos partes enfrentadas, difundiendo los testimonios de las víctimas y rechazando los discursos de odio.} — Le journaliste peut devenir médiateur en donnant la parole aux deux parties, en diffusant les témoignages des victimes et en rejetant les discours de haine.
\end{enumerate}
\textbf{Méthode} : pour répondre, reprends la structure de la question (\esp{¿Por qué...?} → \esp{Se llama... porque...}) et cite les mots exacts du texte.
\end{exoresolu}

\courssub{Premier lexique : les mots-clés du débat}

Le texte introduit le vocabulaire essentiel du module. Voici les sept mots-clés de la fiche de progression, observés dans leur contexte ou déduits logiquement.

\begin{center}
\begin{tabularx}{\linewidth}{|X|X|X|}
\hline
\rowcolor{popblueL} \textbf{Mot espagnol} & \textbf{Français} & \textbf{Exemple / Contexte} \\
\hline
el cuarto poder & le quatrième pouvoir & \esp{Se llama cuarto poder a la prensa...} \\
\hline
la libertad de prensa & la liberté de la presse & \esp{...la libertad de prensa es un tesoro...} \\
\hline
la independencia & l'indépendance & (impliqué par \esp{prensa libre} ; \esp{prensa independiente} = presse indépendante) \\
\hline
la censura & la censure & \esp{...cuando la censura silencia a los periodistas...} \\
\hline
la opinión pública & l'opinion publique & \esp{...influencia decisiva sobre la opinión pública.} \\
\hline
el mediador / la mediadora & le médiateur / la médiatrice & \esp{...el periodista puede convertirse en mediador.} \\
\hline
la reconciliación & la réconciliation & \esp{...favorece el diálogo y la reconciliación.} \\
\hline
\end{tabularx}
\end{center}

\begin{attention}[Erreurs fréquentes des francophones]
\begin{itemize}
\item Genre : \esp{el poder} (masc.), \esp{la opinión}, \esp{la censura}, \esp{la libertad}, \esp{la reconciliación} (fém.). \cle{El mediador} / \cle{la mediadora}.
\item Accentuation : \esp{opinión}, \esp{reconciliación}, \esp{comunicación}, \esp{nación} portent l'accent sur le \esp{ó} final (oxytone) ; au pluriel, l'accent disparait : \esp{opiniones}, \esp{reconciliaciones}.
\item Ordre des mots : \esp{la opinión pública} (nom + adjectif), jamais *\esp{la pública opinión}.
\item \esp{la prensa} (ensemble des médias/journalistes) \neq \esp{la imprenta} (imprimerie/machine).
\item \esp{veraz} (adj. invariable) = véridique ; ne pas confondre avec \esp{verdad} (nom fém. = vérité).
\end{itemize}
\end{attention}

\begin{savaistu}[Les radios de la paix au Guatemala et en Guinée équatoriale]
Au Guatemala, après la guerre civile (1960--1996), les radios communautaires émettant en langues mayas (quiché, mam, kaqchikel...) ont joué un rôle clé dans la réconciliation. Dans les villages reculés, où beaucoup d'habitants ne parlent pas couramment l'espagnol, ces radios ont diffusé des programmes sur les droits des victimes, la mémoire du conflit et la coexistence pacifique. Elles illustrent parfaitement la troisième fonction de la presse : \esp{educar para la paz}. 

Plus près de nous, la Guinée équatoriale, seul pays hispanophone d'Afrique subsaharienne, connaît aussi des débats sur la liberté de la presse et l'indépendance des médias — preuve que le « quatrième pouvoir » est une question universelle, de Madrid à Malabo, de Bogotá à Yaoundé. La CRTV (Cameroon Radio Television) elle-même, service public, a pour mission d'informer avec impartialité et de promouvoir l'unité nationale, ce qui rejoint l'idéal du \esp{cuarto poder} au service de la paix.
\end{savaistu}

\begin{exercice}[À toi de jouer : repérage dans le texte]
\begin{enumerate}
\item Relève dans le texte trois mots de la famille de \esp{informar} ou liés au champ lexical de l'information.
\item Trouve dans le texte deux périphrases qui désignent la presse sans répéter le mot \esp{prensa}.
\item Relevez la métaphore filée du paragraphe 3 (\esp{ciego y sordo}, \esp{juguete}, \esp{tesoro}) et expliquez ce qu'elle suggère.
\item Souligne dans le dernier paragraphe les trois verbes qui expriment la responsabilité des citoyens.
\end{enumerate}
\end{exercice}

\begin{corrige}[Exercice : repérage dans le texte]
\begin{enumerate}
\item \esp{informar}, \esp{la información veraz} (ou \esp{veracidad}), \esp{un ciudadano bien informado} (on accepte aussi \esp{contar los hechos}, \esp{los hechos}).
\item \esp{el cuarto poder} et \esp{los medios de communication} (on accepte aussi \esp{la pluma del periodista} comme métonymie).
\item La métaphore filée oppose la lumière/la vue (\esp{salir a la luz}, \esp{ciego}) à l'obscurité de la censure : sans presse libre, le peuple est comme un aveugle et un sourd, manipulé comme un \esp{juguete} ; la liberté de presse est un \esp{tesoro}, c'est-à-dire un bien précieux et fragile qu'il faut protéger.
\item \esp{leer, escuchar y pensar con espíritu crítico} : les citoyens doivent lire, écouter et penser par eux-mêmes.
\end{enumerate}
\end{corrige}

\begin{aretenir}[Ce qu'il faut retenir de cette première étape]
\begin{itemize}
\item \esp{El cuarto poder} désigne la presse comme pouvoir de contrôle dans une démocratie (expression attribuée à Edmund Burke).
\item Le texte est un \textbf{article d'opinion} : un texte argumentatif avec une thèse, des arguments, des exemples et une conclusion.
\item Les trois fonctions de la presse : \esp{informar}, \esp{denunciar}, \esp{educar para la paz}.
\item La \esp{censura} détruit l'\esp{opinión pública} libre ; le journaliste peut devenir \esp{mediador} et servir la \esp{reconciliación}.
\item Vocabulaire-clé maîtrisé : \cle{libertad de prensa}, \cle{independencia}, \cle{censura}, \cle{opinión pública}, \cle{mediador}, \cle{reconciliación}.
\end{itemize}
\end{aretenir}

% ============================================================
% PARTIE 4 : MÉTHODE DU COMMENTAIRE DE TEXTE
% ============================================================

\coursec{Méthode du commentaire de texte argumentatif}

\courssub{Les quatre temps du commentaire}

Au baccalauréat (épreuve écrite), le commentaire de texte suit une structure rigoureuse en quatre parties. Chaque partie répond à une question précise.

\begin{methode}[Structure type du commentaire de texte (épreuve du Bac)]
\begin{enumerate}
\item \textbf{Introduction} (environ 10 lignes)
      \begin{itemize}
      \item \textbf{Accroche} : contextualisation large (actualité, citation, fait de société).
      \item \textbf{Présentation du texte} : auteur (si connu), source, date, titre, nature (article d'opinion), thème général.
      \item \textbf{Problématique} : question ouverte à laquelle le texte répond (souvent : \emph{Comment l'auteur montre-t-il que... ?} ou \emph{En quoi le texte défend-il l'idée que... ?}).
      \item \textbf{Annonce du plan} : en deux ou trois grandes parties équilibrées.
      \end{itemize}
\item \textbf{Développement} (2 ou 3 parties, chacune en 2--3 paragraphes)
      \begin{itemize}
      \item Chaque partie = une idée-force (argument principal du texte).
      \item Structure d'un paragraphe : \textbf{Idée directrice} (phrase d'intro) $\rightarrow$ \textbf{Citation/Repérage précis} $\rightarrow$ \textbf{Analyse/Explication} (effets de sens, procédés, vocabulaire, syntaxe) $\rightarrow$ \textbf{Mini-synthèse/Transition}.
      \item Équilibre entre analyse de fond (idées) et analyse de forme (langage).
      \end{itemize}
\item \textbf{Ouverture / Approfondissement} (1 paragraphe, 5--8 lignes)
      \begin{itemize}
      \item Élargir la réflexion : autre texte, autre contexte (Cameroun, Guinée équatoriale, histoire), expérience personnelle, actualité récente.
      \item Ne pas introduire de nouvelle analyse technique du texte support.
      \item Poser une question neuve ou une perspective d'avenir.
      \end{itemize}
\item \textbf{Conclusion} (3--4 lignes)
      \begin{itemize}
      \item Réponse synthétique à la problématique.
      \item Dernière phrase marquante (formule choc, citation, appel à l'action).
      \end{itemize}
\end{enumerate}
\end{methode}

\courssub{Outils d'analyse : repérer les procédés argumentatifs}

Pour analyser la \emph{forme}, il faut identifier les procédés utilisés par l'auteur pour convaincre.

\begin{center}
\begin{tabularx}{\linewidth}{|X|X|X|}
\hline
\rowcolor{popblueL} \textbf{Procédé} & \textbf{Définition / Effet} & \textbf{Exemple dans le texte} \\
\hline
\cle{La thèse explicite} & Phrase qui résume l'idée défendue & \esp{Promover una cultura de paz es, pues, una responsabilidad compartida} \\
\hline
\cle{L'énumération} & Liste d'éléments pour montrer l'ampleur & \esp{informar, denunciar y educar} ; \esp{la tolerancia, el diálogo y el respeto} \\
\hline
\cle{L'opposition / antithèse} & Contraste entre deux réalités & \esp{Sin embargo, esta libertad no existe en todas partes} (liberté vs censure) \\
\hline
\cle{La métaphore filée} & Image poursuivie sur plusieurs mots & \esp{ciego y sordo / juguete / tesoro} (aveugle/sourd, jouet, trésor) \\
\hline
\cle{L'exemple concret} & Fait réel qui ancre l'argument & \esp{radios en lenguas mayas... Guatemala} ; \esp{Colombia... proceso de paz} \\
\hline
\cle{L'adresse au lecteur} & \esp{usted / ustedes}, impératif, \esp{hay que} & \esp{hay que defender cada día} ; \esp{los ciudadanos deben leer...} \\
\hline
\cle{Les connecteurs logiques} & Articulent le raisonnement & \esp{porque, desde entonces, sin embargo, por último, por eso, pues, solo así} \\
\hline
\cle{L'hyperbole} & Exagération pour marquer les esprits & \esp{el pueblo queda ciego y sordo} (image forte de la privation) \\
\hline
\cle{Le champ lexical} & Mots d'un même domaine sémantique & Guerre/paix : \esp{guerra civil, heridas, curar, paz, reconciliación, discursos de odio} \\
\hline
\end{tabularx}
\end{center}

\courssub{Conseils de rédaction (français $\leftrightarrow$ espagnol)}

\begin{attention}[Pièges à éviter dans le commentaire]
\begin{itemize}
\item \textbf{Ne pas faire de paraphrase} : ne raconte pas le texte, \textbf{analyse}-le. « L'auteur dit que... » $\rightarrow$ « L'auteur \emph{dénonce} / \emph{montre} / \emph{met en opposition}... »
\item \textbf{Citer avec précision} : guillemets espagnols \esp{« ... »} ou italiques, référence de ligne si possible. Traduire la citation si elle est longue.
\item \textbf{Terminologie d'analyse en français} : thèse, argument, exemple, connecteur, métaphore, énumération, antithèse, champ lexical, ton, registre.
\item \textbf{Terminologie d'analyse en espagnol} (utile pour l'oral) : \esp{la tesis, el argumento, el ejemplo, el conector, la metáfora, la enumeración, la antítesis, el campo léxico, el tono, el registro}.
\item \textbf{Connecteurs pour structurer ton devoir} : \emph{En premier lieu / D'une part... D'autre part / En outre / Par conséquent / En définitive / Pour conclure}.
\end{itemize}
\end{attention}

% ============================================================
% PARTIE 5 : COMMENTAIRE MODÈLE
% ============================================================

\coursec{Commentaire modèle : « El cuarto poder al servicio de la paz »}

\courssub{Sujet}
\textbf{Commentez le texte suivant en montrant comment l'auteur défend l'idée que la presse, en tant que quatrième pouvoir, est un acteur indispensable de la culture de la paix.}

\courssub{Proposition de commentaire rédigé}

\textbf{Introduction}

[Accroche] Dans un monde saturé d'informations instantanées, où les « fake news » et les discours de haine prolifèrent sur les réseaux sociaux, la question du rôle éthique des médias n'a jamais été aussi brûlante. Au Cameroun comme en Colombie, au Guatemala comme en Guinée équatoriale, la radio communautaire ou le journal d'investigation peuvent faire basculer un conflit vers le dialogue ou vers la violence.

[Présentation] Le texte que nous commentons, intitulé « El cuarto poder al servicio de la paz », est un article d'opinion de facture journalistique, sans auteur identifié, probablement destiné à un public scolaire ou citoyen. Il aborde le thème classique du « quatrième pouvoir » pour en renouveler la portée : non plus seulement contrôle démocratique, mais service actif de la paix.

[Problématique] \emph{En quoi cet article démontre-t-il que la presse, par ses trois fonctions traditionnelles — informer, dénoncer, éduquer —, devient un médiateur indispensable pour la réconciliation des peuples ?}

[Annonce du plan] Nous verrons d'abord comment l'auteur fonde la légitimité de la presse sur ses trois fonctions historiques (I), puis comment il alerte sur le danger mortel que représente la censure pour la démocratie (II), pour enfin montrer comment le journaliste se mue en artisan de paix par la médiation et l'éducation critique (III).

\textbf{Développement}

\textbf{I. La presse, quatrième pouvoir : trois piliers pour une citoyenneté éclairée}

L'ouvre par une définition étymologique et historique : \esp{« Se llama cuarto poder a la prensa porque, después de los tres poderes clásicos..., los medios de comunicación ejercen una influencia decisiva sobre la opinión pública »}. La thèse est posée d'emblée : la presse n'est pas un accessoire, elle est un \emph{pouvoir} à part entière. L'auteur structure son argumentation par une \textbf{énumération ternaire} (\esp{informar, denunciar y educar}), rythme binaire/ternaire classique qui donne solennité et complétude.

La première fonction, \esp{informar}, est définie par l'exigence de \esp{veracidad} : \esp{« contar los hechos con veracidad, sin mentiras ni exageraciones »}. L'opposition binaire \esp{veracidad / mentiras} pose l'éthique du journaliste comme garde-fou. L'argument glisse vers le politique : \esp{« Un ciudadano bien informado puede votar, decidir y participar »} — l'information est condition \emph{sine qua non} de la démocratie.

La deuxième, \esp{denunciar}, utilise la métonymie \esp{« la pluma del periodista »} pour incarner l'action. Le champ lexical de la transgression (\esp{injusticias, corrupción, abusos de poder}) et le résultat concret (\esp{« muchos escándalos han salido a la luz »}) ancrent la fonction dans le réel. Le verbe \esp{salir a la luz} (métaphore lumineuse) prépare la métaphore filée du paragraphe suivant.

La troisième, \esp{educar}, élargit le rôle au-delà du factuel : \esp{« formar el espíritu crítico... enseñarles valores como la tolerancia, el diálogo y el respeto del otro »}. L'énumération de valeurs (tolérance, dialogue, respect) inscrit la presse dans une visée \emph{paideutique} (éducative), préfigurant la \esp{cultura de paz} du titre.

\textbf{II. La censure : l'aveuglement collectif et la mort de l'opinion publique}

Le connecteur \esp{« Sin embargo »} marque une rupture nette : l'antithèse. Le danger n'est pas théorique, il est présent (\esp{« no existe en todas partes »}). L'auteur déploie une \textbf{métaphore filée sensorielle} d'une grande force évocatrice : \esp{« el pueblo queda ciego y sordo »}. La privation des sens = privation de connaissance. L'opinion publique, concept clé du premier paragraphe, devient \esp{« juguete de los rumores y de la propaganda »} : de sujet souverain, elle devient objet manipulé. L'hyperbole \esp{« ciego y sordo »} frappe l'esprit.

La conclusion de cette partie est un appel solennel : \esp{« Por eso la libertad de prensa es un tesoro que hay que defender cada día »}. La métaphore du \esp{tesoro} (trésor) implique rareté, valeur, fragilité, et devoir de protection active. L'impersonnel \esp{hay que} + infinitif engage la responsabilité collective.

\textbf{III. Du journaliste au médiateur : l'engagement pour la réconciliation}

Le texte bascule du diagnostic à l'action concrète : \esp{« En los conflictos, además, el periodista puede convertirse en mediador »}. Le verbe \esp{convertirse} (réfléchi) souligne une transformation active, un choix éthique. Trois gérondifs (\esp{dando, difundiendo, rechazando}) détaillent la praxis du médiateur : donner la parole aux deux camps (équilibre), diffuser la voix des victimes (humanisation), rejeter la haine (éthique). La finalité : \esp{« favorece el diálogo y la reconciliación »}.

L'argumentation s'ancre dans deux exemples géographiques précis, symbole de deux continents hispanophones : \esp{Guatemala} (radios en langues mayas, guérison post-guerre civile) et \esp{Colombia} (programmes victimes, processus de paix). Le choix des langues autochtones au Guatemala illustre l'\esp{educar} inclusif ; l'accompagnement du processus colombien montre le rôle dans la \emph{justice transitionnelle}.

\textbf{Ouverture}

Cet article résonne avec l'actualité camerounaise. À Yaoundé, Douala, ou dans les zones anglophones en crise, les journalistes qui osent \esp{dar la palabra a las dos partes} — comme ce reporter de la CRTV imaginé en introduction — jouent exactement ce rôle de \esp{mediador}. De même, en Guinée équatoriale, la lutte pour une \esp{prensa independiente} conditionne la \esp{reconciliación} nationale. Mais au-delà des médias professionnels, chaque citoyen devient « journaliste » sur les réseaux sociaux : la \esp{cultura de paz} exige de nous tous \esp{leer, escuchar y pensar con espíritu crítico}, comme le conclut le texte. La paix n'est pas un don, c'est un \emph{devoir d'information}.

\textbf{Conclusion}

En définitive, ce texte démontre que le \esp{cuarto poder} n'est pas une métaphore désuète mais une nécessité vitale. Informer, dénoncer, éduquer : trois verbes qui conjugués au présent construisent l'avenir. La dernière phrase — \esp{« Solo así el cuarto poder será, de verdad, un poder al servicio de la vida y de la reconciliación entre los pueblos »} — scelle l'engagement : la presse libre est le premier rempart contre la barbarie. À nous, citoyens et futurs communicateurs, d'en être les gardiens.

% ============================================================
% PARTIE 6 : PRODUCTION GUIDÉE — DÉFENDRE LA CULTURE DE LA PAZ
% ============================================================

\coursec{Production guidée : défendre la culture de la paz}

\courssub{Expression écrite : article d'opinion / tribune libre}

\begin{methode}[Rédiger un article d'opinion (150--200 mots)]
\begin{enumerate}
\item \textbf{Analyser le sujet} : identifie la thèse à défendre, le public visé, le ton (engagé, mesuré, solennel).
\item \textbf{Trouver un titre accrocheur} : court, interrogatif ou affirmatif, avec vocabulaire fort (\esp{La voz de la paz}, \esp{Periodismo: arma contra el odio}, \esp{El micrófono de la reconciliación}).
\item \textbf{Structurer} :
      \begin{itemize}
      \item \textbf{Lead} (premier paragraphe) : accroche + thèse en une phrase.
      \item \textbf{Corps} : 2 arguments + 1 exemple concret (Cameroun, Guinée équatoriale, expérience personnelle).
      \item \textbf{Conclusion} : appel à l'action / phrase choc.
      \end{itemize}
\item \textbf{Soigner la langue} : connecteurs logiques (\esp{en primer lugar, por tanto, sin embargo, en definitiva}), champ lexical paix/guerre, impersonnels d'obligation (\esp{hay que, es necesario que + subjuntivo}), futur pour l'engagement.
\item \textbf{Relire} : accords, accents, \esp{ser/estar}, \esp{por/para}, subjonctif après expression de volonté/nécessité.
\end{enumerate}
\end{methode}

\begin{exemplebox}[Modèle d'article d'opinion (sujet : « La radio scolaire, outil de paix au lycée »)]
\textbf{La voz de los jóvenes: la radio escolar contra la violencia}

En los recreos de nuestro liceo en Yaoundé, los auriculares aíslan a los estudiantes; pero cuando suena la radio escolar, todos escuchan la misma voz. \textbf{La radio escolar no es un simple pasatiempo: es un laboratorio de ciudadanía y un antídoto contra la violencia.}

En primer lugar, informar desde dentro cambia la mirada. Cuando los alumnos periodistas cubren un conflicto entre clases —un balón disputado, un malentendido étnico— y dan la palabra a las dos partes, practican la \esp{veracidad} y la \esp{imparcialidad}. Aprenden que \esp{informar} no es repetir rumores, sino contrastar fuentes.

En segundo lugar, la radio educa en el diálogo. Nuestros programas en francés, en inglés y en lenguas nacionales (ewondo, bassa, fulfulde) valoran la diversidad lingüística de Camerún. Como las radios mayas en Guatemala, enseñan que \esp{la reconciliación} empieza por escuchar al otro en su propia lengua.

\esp{Por tanto}, hay que defender este espacio de libertad. \esp{Hay que} dar micrófonos a los estudiantes, formarles en ética periodística y proteger su \esp{independencia}. \esp{Solo así} el recreo dejará de ser zona de tensión para convertirse en \esp{agora} de paz.

\esp{La paz no cae del cielo; se construye en las ondas de una radio escolar.}
\end{exemplebox}

\courssub{Expression orale : débat structuré / jeu de rôle}

\begin{experience}[Débat : « La liberté de presse a-t-elle des limites ? »]
\textbf{Consigne.} En groupes de 4, préparez un débat de 6 minutes (3 min par camp) sur la motion : \esp{« La libertad de prensa debe ser absoluta, sin ninguna limitación. »}

\textbf{Rôles :}
\begin{itemize}
\item \textbf{Groupe A (Pour)} : journalistes, ONG de défense des droits de l'homme, blogueurs. Arguments : droit à l'info, contrôle du pouvoir, \esp{artículo 19 DUDH}, censure = dictature.
\item \textbf{Groupe B (Contre / Nuancé)} : gouvernants, juristes, parents d'élèves. Arguments : protection vie privée, sécurité nationale, discours de haine, \esp{por/para} la protección de menores, diffamation.
\item \textbf{Modérateur} : donne la parole, reformule, veille au temps, fait voter à la fin.
\item \textbf{Observateur} : note les arguments, le vocabulaire utilisé, les erreurs de langue, les connecteurs.
\end{itemize}

\textbf{Lexique utile pour le débat :}
\begin{center}
\begin{tabularx}{\linewidth}{|X|X|}
\hline
\rowcolor{popblueL} \textbf{Pour argumenter} & \textbf{Pour nuancer / contredire} \\
\hline
\esp{Estoy de acuerdo en que...} & \esp{No estoy de acuerdo, porque...} \\
\esp{Es innegable que...} & \esp{Es cierto, pero... / Sin embargo...} \\
\esp{El argumento principal es...} & \esp{Ese argumento no vale, ya que...} \\
\esp{Pongamos un ejemplo: ...} & \esp{En el caso de... vemos lo contrario.} \\
\esp{Por tanto / Por consiguiente,} & \esp{En cambio / Por el contrario,} \\
\hline
\end{tabularx}
\end{center}

\textbf{Évaluation (grille simplifiée) :}
\begin{itemize}
\item Contenu / argumentation (4 pts) : pertinence, exemples, structure.
\item Corrección lingüística (4 pts) : grammaire, vocabulaire, accents.
\item Fluidez / pronunciación (4 pts) : débit, intonation, \esp{v/b}, \esp{j/g}, \esp{ll/y}, \esp{r/rr}.
\item Interacción (4 pts) : écoute, réponses aux arguments adverses, politesse (\esp{permítame, disculpe, tiene razón en parte}).
\item Esprit critique / culture (4 pts) : références Cameroun / monde hispanophone.
\end{itemize}
\end{experience}

\courssub{Traduction : Thème (Français $\rightarrow$ Espagnol)}

\begin{exoresolu}[Thème : La presse et la paix]
\textbf{Énoncé.} Traduisez en espagnol le paragraphe suivant :

« Dans notre pays, la liberté de la presse est un trésor fragile. Quand les journalistes sont muselés, le peuple devient aveugle et sourd face à la réalité. Cependant, un reporter courageux peut devenir médiateur : en donnant la parole aux victimes et en refusant les discours de haine, il prépare la réconciliation. C'est pourquoi nous devons défendre chaque jour le droit à l'information. »

\tcblower
\textbf{Proposition de traduction.}

\esp{En nuestro país, la libertad de prensa es un tesoro frágil. Cuando los periodistas son silenciados (o: amordazados), el pueblo queda ciego y sordo ante la realidad. Sin embargo, un reportero valiente puede convertirse en mediador: dando la voz a las víctimas y rechazando los discursos de odio, prepara la reconciliación. Por eso debemos defender cada día el derecho a la información.}

\textbf{Points de vigilance (pièges pour francophones) :}
\begin{itemize}
\item \esp{un tesoro frágil} (accord adjectif, pas *fragile).
\item \esp{son silenciados / amordazados} (passif \esp{ser} + participe accordé) ; \esp{amordazar} = museler (mettre un bâillon), image forte pour journalistes.
\item \esp{queda ciego y sordo} (\esp{quedar} + adjectif = devenir/être ; accord au masculin singulier avec \esp{pueblo}).
\item \esp{ante la realidad} (devant / face à) ; pas *\esp{enfrente de}.
\item \esp{reportero} (anglicisme admis) ou \esp{periodista} / \esp{corresponsal}.
\item \esp{dando la voz a las víctimas} (gérondif de manière) ; \esp{rechazando los discursos de odio}.
\item \esp{prepara la reconciliación} (présent de l'indicatif, valeur de conséquence immédiate).
\item \esp{Por eso debemos defender...} (\esp{deber} + infinitif = obligation morale) ; \esp{el derecho a la información} (nom + \esp{a} + infinitif nominalisé).
\end{itemize}
\end{exoresolu}

\courssub{Traduction : Version (Espagnol $\rightarrow$ Français)}

\begin{exoresolu}[Version : Extrait du texte support]
\textbf{Énoncé.} Traduisez en français le passage suivant (lignes 13--18 du texte) :

\esp{« Sin embargo, esta libertad no existe en todas partes. Cuando la censura silencia a los periodistas, cuando los periódicos son cerrados y las radios intervenidas, el pueblo queda ciego y sordo: ya no sabe lo que ocurre en su propio país. La opinión pública, privada de información veraz, se convierte en juguete de los rumores y de la propaganda. »}

\tcblower
\textbf{Proposition de traduction.}

« Cependant, cette liberté n'existe pas partout. Quand la censure réduit les journalistes au silence, quand les journaux sont fermés et les radios mises sous tutelle (ou : saisies), le peuple reste aveugle et sourd : il ne sait plus ce qui se passe dans son propre pays. L'opinion publique, privée d'information véridique, devient le jouet (la proie) des rumeurs et de la propagande. »

\textbf{Notes de traduction :}
\begin{itemize}
\item \esp{Sin embargo} = Cependant / Toutefois (registre soutenu).
\item \esp{no existe en todas partes} = n'existe pas partout (pas *\esp{en todo lugar}).
\item \esp{son cerrados y las radios intervenidas} = passif \esp{ser} + participe (accord : \esp{cerrados} masc. pl., \esp{intervenidas} fém. pl.). \esp{Intervenir} une radio = la placer sous contrôle étatique / la saisir.
\item \esp{queda ciego y sordo} : \esp{quedar} + adjectif = « rester/devenir » avec changement d'état ; adjectifs au masculin singulier accordés à \esp{pueblo}.
\item \esp{ya no sabe} = ne sait plus / ne sait déjà plus (\esp{ya no} = plus).
\item \esp{privada de información veraz} : participe passé fém. sg. accordé à \esp{opinión pública} (gérondif/participe adjectival).
\item \esp{se convierte en juguete} : \esp{convertirse en} = devenir (changement de nature) ; \esp{juguete} = jouet, mais ici sens figuré « proie manipulée ».
\end{itemize}
\end{exoresolu}

\courssub{Exercices d'entraînement (Bac blanc)}

\begin{exercice}[Compréhension détaillée + Expression]
\begin{enumerate}
\item \textbf{Repérage} : Relevez dans le texte les connecteurs logiques et classez-les : cause, conséquence, opposition, addition, conclusion.
\item \textbf{Analyse} : « \esp{La pluma del periodista señala las injusticias} ». Identifiez la figure de style et expliquez son effet.
\item \textbf{Subjonctif / Indicatif} : Complétez avec le mode et le temps justifiés :
      \begin{enumerate}
      \item Es necesario que la prensa \_\_\_\_\_\_\_\_ (ser) libre.
      \item El gobierno quiere que los periodistas \_\_\_\_\_\_\_\_ (callar).
      \item No creo que la censura \_\_\_\_\_\_\_\_ (resolver) los problemas.
      \item Es verdad que los medios \_\_\_\_\_\_\_\_ (informar) con veracidad.
      \end{enumerate}
\item \textbf{Production} : Rédigez une tribune de 180 mots environ pour le journal de votre lycée sur le thème : \esp{« Los estudiantes también somos el cuarto poder: nuestra voz cuenta para la paz »}. Utilisez au moins 5 mots du lexique de la leçon et 3 connecteurs logiques.
\end{enumerate}
\end{exercice}

\begin{corrige}[Exercices d'entraînement]
\begin{enumerate}
\item \textbf{Connecteurs} : Cause : \esp{porque, ya que} (impliqué) ; Conséquence : \esp{por eso, por tanto, por consiguiente, solo así} ; Opposition : \esp{sin embargo} ; Addition : \esp{y, además, por último} ; Conclusion : \esp{pues, en definitiva}.
\item \textbf{Figure} : \esp{La pluma del periodista} = \textbf{métonymie} (l'instrument pour l'agent / l'action). Effet : concrétise l'abstrait, personnifie l'écriture, donne une image visuelle forte (la plume qui pointe, qui dénonce).
\item \textbf{Modes et temps} :
      \begin{enumerate}
      \item \esp{sea} (subj. présent) — \esp{es necesario que} + subjonctif (nécessité).
      \item \esp{callen} (subj. présent) — \esp{querer que} + subjonctif (volonté).
      \item \esp{resuelva} (subj. présent) — \esp{no creer que} + subjonctif (doute/négation).
      \item \esp{informan} (ind. présent) — \esp{es verdad que} + indicatif (certitude/réel).
      \end{enumerate}
\item \textbf{Production} : (Exemple de début) \esp{« Los estudiantes también somos el cuarto poder... »} Vérifier : 5 mots-clés (\esp{libertad de prensa, censura, opinión pública, mediador, reconciliación, independencia, cuarto poder}), 3 connecteurs (\esp{en primer lugar, sin embargo, por tanto}), subjonctifs après \esp{es necesario que / quiero que / exijo que}.
\end{enumerate}
\end{corrige}

\begin{aretenir}[Synthèse finale de la leçon E24]
\begin{itemize}
\item \textbf{Texte} : Article d'opinion argumentatif sur le \esp{cuarto poder} et la \esp{cultura de paz}.
\item \textbf{Structure} : Thèse (3 fonctions) — Antithèse (censure) — Synthèse/Action (médiateur, exemples Guatemala/Colombie) — Responsabilité partagée.
\item \textbf{Lexique clé} : \cle{cuarto poder, libertad de prensa, independencia, censura, opinión pública, mediador/mediadora, reconciliación, veracidad, imparcialidad, discurso de odio}.
\item \textbf{Grammaire} : Impersonnels (\esp{hay que, es necesario que} + subj.), passif \esp{ser} + part., gérondif de manière, connecteurs logiques, accord adjectifs (genre/nombre), \esp{ser/estar} (\esp{queda ciego} = état résultant).
\item \textbf{Compétences Bac} : Commentaire structuré (intro, développement 2--3 parties, ouverture, conclusion), thème/version, article d'opinion, débat oral.
\item \textbf{Contexte camerounais / hispanophone} : CRTV, radios communautaires, Guinée équatoriale, Guatemala, Colombie — la paix se construit par la parole libre et partagée.
\end{itemize}
\end{aretenir}

\coursec{Compréhension du texte}

Dans la section précédente, nous avons lu et découvert le texte support \emph{El cuarto poder al servicio de la paz}. Nous avons situé son thème : le rôle des médias dans la construction de la paix. Il s'agit maintenant de vérifier et d'approfondir notre compréhension, étape indispensable avant tout commentaire de texte. Au baccalauréat, les questions de compréhension (\esp{preguntas de comprensión}) portent toujours sur quatre niveaux : le repérage, l'idée générale, les détails et l'interprétation. Nous allons les travailler un par un, méthodiquement.

\courssub{Rappel du texte support}

Relisons d'abord le texte, car toute réponse doit s'appuyer sur lui.

\begin{textebox}[El cuarto poder al servicio de la paz]
Se llama « cuarto poder » a la prensa porque, después del poder ejecutivo, el legislativo y el judicial, los medios de comunicación influyen profundamente en la vida de una nación. Sin embargo, este poder solo es útil si se ejerce con libertad y responsabilidad.

Una prensa libre cumple tres funciones esenciales. Primero, informa : gracias a los periodistas, los ciudadanos conocen lo que ocurre en su país y en el mundo. Segundo, denuncia : una prensa libre denuncia las injusticias, la corrupción y los abusos de poder. Tercero, educa y media : los medios pueden enseñar el diálogo y servir de mediadores entre comunidades enfrentadas.

Pero ¿qué ocurre cuando hay censura ? Cuando los gobiernos callan a los periodistas, el pueblo queda ciego y sordo : ya no sabe lo que pasa, ya no puede decidir con conocimiento. La censura, pues, no solo ataca a los periodistas, sino también a todos los ciudadanos.

En conclusión, la libertad de prensa es una condición de la paz. Unos medios independientes y responsables promueven la reconciliación y la cultura de paz, porque un pueblo bien informado es un pueblo que dialoga en lugar de luchar.
\end{textebox}

\textbf{Glossaire :} \esp{el poder ejecutivo} : le pouvoir exécutif ; \esp{callar} : faire taire, réduire au silence ; \esp{enfrentadas} : opposées, en conflit ; \esp{queda ciego y sordo} : devient aveugle et sourd ; \esp{con conocimiento} : en connaissance de cause ; \esp{en lugar de luchar} : au lieu de se battre.

\courssub{La méthode pour répondre à une question de compréhension}

Avant de répondre aux questions, apprenons la bonne démarche. Beaucoup d'élèves perdent des points non pas parce qu'ils n'ont pas compris, mais parce qu'ils répondent mal.

\begin{methode}[Répondre à une question de compréhension en trois étapes]
\begin{enumerate}
\item \textbf{Reformuler} la question dans la réponse : à \esp{¿Cuál es la tesis del autor ?}, on commence par \esp{La tesis del autor es que...}. Ne répondez jamais par un simple mot isolé.
\item \textbf{Citer le texte} entre guillemets pour justifier : \esp{El autor escribe : « la libertad de prensa es una condición de la paz »}.
\item \textbf{Expliquer avec ses propres mots} ce que la citation signifie, sans la répéter mot à mot.
\end{enumerate}
\end{methode}

\begin{exemplebox}[Une réponse complète et correcte]
\textbf{Question :} \esp{¿Qué hace una prensa libre ?}

\textbf{Réponse modèle :} \esp{Según el texto, una prensa libre « denuncia las injusticias ». Esto significa que los periodistas revelan públicamente los actos injustos para que los ciudadanos los conozcan.}

Traduction : « Selon le texte, une presse libre "dénonce les injustices". Cela signifie que les journalistes révèlent publiquement les actes injustes afin que les citoyens les connaissent. »

Observez la construction : reformulation (\esp{según el texto}), citation entre guillemets, puis explication introduite par \esp{esto significa que}. Remarquez aussi le subjonctif \esp{conozcan} après \esp{para que} (but) : c'est une règle obligatoire.
\end{exemplebox}

\courssub{Questions de repérage : de quoi parle le texte ?}

Les questions de repérage sont les plus simples : elles demandent d'identifier le sujet, le locuteur, le destinataire, le lieu et le moment. Voici les réponses attendues pour notre texte.

\begin{itemize}
\item \esp{¿De qué habla el texto ?} (De quoi parle le texte ?) — Le texte parle du \cle{cuarto poder}, c'est-à-dire de la presse, et de son rôle dans la promotion de la paix.
\item \esp{¿Quién habla ?} (Qui parle ?) — C'est un \cle{articulista}, un journaliste ou un essayiste, qui s'adresse aux lecteurs. Le texte est impersonnel : il n'y a pas de « yo », mais un énonciateur qui expose des faits et défend une opinion.
\item \esp{¿A quién se dirige ?} (À qui s'adresse-t-il ?) — À l'opinion publique, aux citoyens, et en particulier aux jeunes lecteurs, aux élèves et étudiants.
\item \esp{¿Dónde y cuándo ?} (Où et quand ?) — Le texte ne précise ni lieu ni date : c'est un article d'opinion \emph{intemporel}, valable pour tout pays, y compris le Cameroun, où la question de la liberté de presse est d'actualité.
\end{itemize}

\begin{propriete}[Où trouver la thèse dans un texte argumentatif]
Dans un texte argumentatif, la \cle{tesis} (la thèse, l'idée principale défendue par l'auteur) se trouve le plus souvent dans l'\textbf{introduction} ou dans la \textbf{conclusion}. Les arguments, eux, sont développés dans le corps du texte et sont généralement introduits par des \cle{conectores argumentativos} : \esp{primero}, \esp{segundo}, \esp{además}, \esp{sin embargo}, \esp{por eso}, \esp{en conclusión}. Repérer ces connecteurs, c'est repérer le plan de l'auteur.
\end{propriete}

\courssub{L'idée générale : formuler la thèse en une phrase}

\textbf{Question :} \esp{¿Cuál es la idea general del texto ? ¿Cuál es la tesis del autor ?}

La thèse de l'auteur peut se formuler ainsi : \esp{La libertad de prensa es una condición de la paz porque una prensa libre informa, denuncia y educa, mientras que la censura deja al pueblo ciego y sordo.} (« La liberté de la presse est une condition de la paix, car une presse libre informe, dénonce et éduque, tandis que la censure laisse le peuple aveugle et sourd. »)

On remarque que cette thèse est annoncée presque mot pour mot dans la conclusion du texte : \esp{« la libertad de prensa es una condición de la paz »}. C'est la confirmation de la règle énoncée plus haut.

\begin{attention}[Ne recopiez pas tout le texte]
À la question \esp{¿cuál es la tesis del autor ?}, l'erreur classique des élèves francophones est de recopier un paragraphe entier, ou même tout le texte, en espérant que la bonne réponse s'y trouve. Le correcteur sanctionne ce procédé. La thèse doit être \textbf{résumée en une seule phrase}, avec vos propres mots ou une courte citation encadrée de guillemets. Une phrase, pas dix.
\end{attention}

\courssub{Questions de détail : les trois fonctions et la censure}

\textbf{Question 1 :} \esp{¿Cuáles son las tres funciones de la prensa citadas en el texto ?}

Le texte énumère explicitement trois fonctions, signalées par \esp{primero}, \esp{segundo}, \esp{tercero} :

\begin{enumerate}
\item \esp{Informar} : informer — \esp{« los ciudadanos conocen lo que ocurre en su país y en el mundo »} (les citoyens savent ce qui se passe dans leur pays et dans le monde).
\item \esp{Denunciar} : dénoncer — \esp{« una prensa libre denuncia las injusticias, la corrupción y los abusos de poder »} (une presse libre dénonce les injustices, la corruption et les abus de pouvoir).
\item \esp{Educar y mediar} : éduquer et servir de médiateur — \esp{« los medios pueden enseñar el diálogo y servir de mediadores »} (les médias peuvent enseigner le dialogue et servir de médiateurs).
\end{enumerate}

\textbf{Question 2 :} \esp{¿Qué ocurre cuando hay censura ?}

Réponse : \esp{Cuando hay censura, « el pueblo queda ciego y sordo » : los ciudadanos ya no saben lo que pasa y no pueden decidir con conocimiento.} (« Quand il y a censure, le peuple devient aveugle et sourd : les citoyens ne savent plus ce qui se passe et ne peuvent plus décider en connaissance de cause. »)

\begin{exemplebox}[Exemples traduits à retenir]
\esp{Una prensa libre denuncia las injusticias.} « Une presse libre dénonce les injustices. »

\esp{El pueblo queda ciego y sordo.} « Le peuple devient aveugle et sourd. »

\esp{Un pueblo bien informado es un pueblo que dialoga en lugar de luchar.} « Un peuple bien informé est un peuple qui dialogue au lieu de se battre. »
\end{exemplebox}

\courssub{Questions sur les procédés : les connecteurs argumentatifs}

\textbf{Question :} \esp{Señala los conectores argumentativos del texto y explica su función.}

Relevons les connecteurs et précisons leur rôle :

\begin{center}
\begin{tabularx}{\linewidth}{|l|X|X|}
\hline
\rowcolor{popblueL} \textbf{Conector} & \textbf{Función} & \textbf{Équivalent français} \\
\hline
\esp{sin embargo} & opposition, nuance & cependant, pourtant \\
\hline
\esp{primero / segundo / tercero} & énumération des arguments & premièrement, deuxièmement, troisièmement \\
\hline
\esp{pues} & explication, conséquence & donc, en effet, or \\
\hline
\esp{también} & addition & aussi, également \\
\hline
\esp{pero} & opposition forte & mais \\
\hline
\esp{en conclusión} & annonce de la conclusion & en conclusion \\
\hline
\esp{porque} & cause & parce que \\
\hline
\end{tabularx}
\end{center}

Analysons trois d'entre eux dans le texte :

\begin{itemize}
\item \esp{« Sin embargo, este poder solo es útil si se ejerce con libertad »} : \esp{sin embargo} introduit une \textbf{restriction} — le pouvoir de la presse a une condition.
\item \esp{« La censura, pues, no solo ataca a los periodistas »} : \esp{pues} marque ici la \textbf{conséquence logique} du raisonnement précédent (attention : en français on ne le traduit pas toujours ; il peut se rendre par « donc », « ainsi » ou ne pas se traduire).
\item \esp{« sino también a todos los ciudadanos »} : la structure \esp{no solo... sino también} (« non seulement... mais aussi ») ajoute un élément plus large et plus grave.
\end{itemize}

\begin{attention}[Faux amis et calques]
\esp{Sin embargo} ne signifie jamais « sans embargo » : c'est « cependant ». \esp{Actualmente} signifie « actuellement » au sens de « en ce moment », pas « effectivement ». Enfin, ne traduisez pas « en effet » mécaniquement : pour introduire une explication ou une conséquence, l'espagnol emploie le plus souvent \esp{pues} ou \esp{por eso} ; \esp{en efecto} existe, mais il appartient à un registre soutenu et sert à confirmer ce qui vient d'être dit.
\end{attention}

\courssub{Question d'interprétation : la métaphore du peuple aveugle et sourd}

\textbf{Question :} \esp{¿Por qué dice el autor que « el pueblo queda ciego y sordo » ? Explica la metáfora.}

C'est la question la plus difficile, car elle demande d'interpréter une image. Procédons par observation : quels organes permettent de s'informer ? Les yeux (lire les journaux, regarder les informations) et les oreilles (écouter la radio). Si la censure supprime les journaux, les radios et les chaînes libres, le peuple perd l'usage symbolique de ces organes : il ne \emph{voit} plus et n'\emph{entend} plus ce qui se passe réellement.

\begin{definition}[La metáfora]
Une \cle{metáfora} (métaphore) est une comparaison implicite, sans outil de comparaison (\esp{como} = « comme »). Ici, l'auteur ne dit pas que le peuple est \emph{comme} un aveugle : il dit qu'il \emph{devient} aveugle et sourd. La cécité symbolise l'impossibilité de voir la vérité ; la surdité, l'impossibilité d'entendre les informations. La métaphore exprime donc la \textbf{privation d'information}, qui rend le peuple incapable de juger et de décider.
\end{definition}

Réponse modèle complète : \esp{El autor utiliza una metáfora : los ojos representan la capacidad de ver la verdad y los oídos, la capacidad de recibir información. Con la censura, el pueblo pierde simbólicamente estos sentidos : ya no puede conocer la realidad ni decidir con conocimiento. La metáfora subraya que la censura es una forma de violencia contra todo el pueblo, no solo contra los periodistas.}

\courssub{Organigramme de la structure argumentative}

La figure suivante résume le cheminement de la pensée de l'auteur, de la thèse à la conclusion.

\begin{popfigure}
\begin{center}
\begin{tikzpicture}[
  node distance=0.55cm,
  caja/.style={draw=popblue, fill=popblueL, rounded corners=3pt, align=center, text width=6.2cm, inner sep=5pt, font=\small},
  tesis/.style={draw=poppurple, fill=poppurpleL, rounded corners=3pt, align=center, text width=6.2cm, inner sep=5pt, font=\small\bfseries},
  concl/.style={draw=popgreen, fill=popgreenL, rounded corners=3pt, align=center, text width=6.2cm, inner sep=5pt, font=\small\bfseries},
  flecha/.style={-{Stealth[length=2.5mm]}, thick, popdark}
]
\node[tesis, align=center] (t){TESIS\\La libertad de prensa es una condición de la paz};
\node[caja, below=of t, align=center] (a1){ARGUMENTO 1 : \emph{informar}\\Los ciudadanos conocen lo que ocurre};
\node[caja, below=of a1, align=center] (a2){ARGUMENTO 2 : \emph{denunciar}\\Una prensa libre denuncia las injusticias};
\node[caja, below=of a2, align=center] (a3){ARGUMENTO 3 : \emph{educar y mediar}\\Los medios enseñan el diálogo y median};
\node[caja, below=of a3, draw=poporange, fill=poporangeL, align=center] (cont){CONTRAARGUMENTO\\Con la censura, el pueblo queda ciego y sordo};
\node[concl, below=of cont, align=center] (c){CONCLUSIÓN\\Medios independientes = cultura de paz};
\draw[flecha] (t) -- (a1);
\draw[flecha] (a1) -- (a2);
\draw[flecha] (a2) -- (a3);
\draw[flecha] (a3) -- (cont);
\draw[flecha] (cont) -- (c);
\end{tikzpicture}
\end{center}
\legende{Structure argumentative du texte : de la thèse à la conclusion}
\end{popfigure}

\courssub{Question personnelle : donner son opinion}

\textbf{Question :} \esp{¿Crees que la prensa camerunesa es un mediador en los conflictos ? Justifica tu respuesta.}

Cette question n'a pas de réponse unique : le correcteur évalue votre capacité à \textbf{donner une opinion et à la justifier} en espagnol correct. Utilisez les formules d'opinion : \esp{creo que}, \esp{pienso que}, \esp{en mi opinión}, \esp{me parece que}, suivies de l'indicatif ; et \esp{no creo que} suivi du \textbf{subjonctif}.

\begin{exemplebox}[Deux réponses possibles]
\esp{En mi opinión, la prensa camerunesa puede ser un mediador en los conflictos porque algunos programas de radio invitan a las comunidades a dialogar. Por ejemplo, durante las elecciones, ciertos periodistas piden calma y respeto.}

\esp{No creo que la prensa sea siempre un mediador : a veces algunos medios difunden rumores que dividen a la población. Sin embargo, pienso que una prensa responsable puede promover la reconciliación.}
\end{exemplebox}

\begin{attention}[Subjonctif après « no creo que »]
Après \esp{creo que} et \esp{pienso que} (opinion affirmative), on emploie l'indicatif : \esp{Creo que la prensa \textbf{es} un mediador.} Mais après la négation \esp{no creo que}, l'espagnol exige le subjonctif : \esp{No creo que la prensa \textbf{sea} siempre un mediador.} Ne dites jamais \esp{no creo que la prensa es...}. Même règle avec \esp{no pienso que}, \esp{no me parece que}.
\end{attention}

\courssub{Entraînement type Bac : vrai ou faux justifié}

Au baccalauréat, l'exercice \esp{verdadero o falso} exige de cocher puis de \textbf{justifier par une citation exacte du texte}. Sans citation, la réponse ne vaut rien, même si elle est juste.

\begin{exoresolu}[Verdadero o falso : justifica con una cita del texto]
\textbf{Énoncé.} Indiquez si les affirmations suivantes sont \esp{verdadero} (V) ou \esp{falso} (F), en citant le texte.

1. \esp{La prensa se llama « cuarto poder » porque es el poder más importante del Estado.}

2. \esp{Según el texto, los medios de comunicación pueden servir de mediadores.}

3. \esp{La censura solo perjudica a los periodistas.}

4. \esp{Para el autor, un pueblo bien informado dialoga en lugar de luchar.}
\tcblower
\textbf{Solution détaillée.}

1. \textbf{Falso.} Le texte dit : \esp{« después del poder ejecutivo, el legislativo y el judicial, los medios de comunicación influyen profundamente »}. La presse vient \emph{après} les trois pouvoirs ; le texte ne dit jamais qu'elle est le pouvoir le plus important.

2. \textbf{Verdadero.} Citation : \esp{« los medios pueden enseñar el diálogo y servir de mediadores entre comunidades enfrentadas »}.

3. \textbf{Falso.} Citation : \esp{« La censura, pues, no solo ataca a los periodistas, sino también a todos los ciudadanos »}. La structure \esp{no solo... sino también} contredit exactement l'affirmation.

4. \textbf{Verdadero.} Citation finale du texte : \esp{« un pueblo bien informado es un pueblo que dialoga en lugar de luchar »}.

\textbf{Méthode :} pour chaque item, repérez dans le texte la phrase qui correspond à l'affirmation, comparez-la mot à mot (attention aux petits mots comme \esp{solo}, \esp{no}, \esp{más}), puis copiez la citation entre guillemets.
\end{exoresolu}

\begin{exercice}[À votre tour]
Répondez en espagnol, en phrases complètes, aux questions suivantes sur le texte \emph{El cuarto poder al servicio de la paz} :

1. \esp{¿Por qué se llama « cuarto poder » a la prensa ?}

2. \esp{¿Qué condición pone el autor para que el poder de la prensa sea útil ? Cita el texto.}

3. Verdadero o falso (justifica) : \esp{« La libertad de prensa no tiene ninguna relación con la paz. »}

4. \esp{En tu opinión, ¿qué función de la prensa es la más importante : informar, denunciar o educar ? Justifica tu respuesta.}
\end{exercice}

\begin{corrige}[Exercice « À votre tour »]
1. \esp{Se llama « cuarto poder » a la prensa porque, después de los tres poderes clásicos (ejecutivo, legislativo y judicial), los medios de comunicación influyen profundamente en la vida de una nación.}

2. \esp{El autor escribe : « este poder solo es útil si se ejerce con libertad y responsabilidad ». La condición es, pues, la libertad y la responsabilidad.}

3. \textbf{Falso.} Citation : \esp{« la libertad de prensa es una condición de la paz »}. Le texte affirme le contraire de l'affirmation.

4. Réponse personnelle possible : \esp{En mi opinión, la función más importante es informar, porque sin información los ciudadanos no pueden ni denunciar ni dialogar. Un pueblo informado es un pueblo libre.}
\end{corrige}

\begin{aretenir}[L'essentiel de la section]
\begin{itemize}
\item Une réponse de compréhension se construit en trois temps : reformuler, citer entre guillemets, expliquer avec ses mots.
\item La thèse se résume en \textbf{une seule phrase} ; elle se trouve souvent dans l'introduction ou la conclusion.
\item Les connecteurs (\esp{sin embargo}, \esp{pues}, \esp{también}, \esp{en conclusión}) balisent la structure argumentative : apprenez à les repérer et à dire leur fonction.
\item Au \esp{verdadero o falso}, toute réponse non justifiée par une citation est nulle.
\item Après \esp{no creo que}, employez le subjonctif ; après \esp{creo que}, l'indicatif.
\end{itemize}
\end{aretenir}

\coursec{Lexique thématique par champs}

Après avoir analysé le texte \emph{El cuarto poder al servicio de la paz} (section précédente), nous allons maintenant structurer le vocabulaire essentiel en quatre champs lexicaux cohérents. Cette organisation par thèmes permet de mémoriser les mots en réseau, de repérer les familles morphologiques et d'éviter les confusions fréquentes chez les francophones (faux amis, homonymes, calques). Chaque champ sera illustré par un tableau bilingue, des exemples contextualisés, des points de grammaire lexicale et un exercice résolu.

\courssub{Champ 1 — La presse et ses acteurs}

\begin{textebox}[Extrait journalistique simulé]
En la \cle{redacción} de \emph{El Diario de Yaoundé}, los \cle{periodistas} preparan el \cle{reportaje} sobre la cosecha de cacao. El \cle{titular} de la portada anuncia: « \esp{Los agricultores exigen precios justos.} » Dentro, un \cle{artículo} de opinión analiza la \cle{libertad de prensa} y una \cle{entrevista} da la palabra a un líder cooperativo.
\tcblower
\emph{Glossaire :} \cle{redacción} = rédaction (lieu et équipe) ; \cle{periodista} = journaliste ; \cle{reportaje} = reportage ; \cle{titular} = titre (une) ; \cle{artículo} = article ; \cle{entrevista} = interview ; \cle{agricultor} = cultivateur ; \cle{exigir} = exiger ; \cle{precio justo} = prix équitable.
\end{textebox}

\begin{center}
\begin{tabularx}{\linewidth}{|X|X|X|}
\hline
\rowcolor{popblueL} \textbf{Español} & \textbf{Français} & \textbf{Remarque} \\
\hline
la prensa & la presse (ensemble) & nom collectif, féminin \\
el periódico & le journal (quotidien) & synonyme : \emph{el diario} \\
el/la periodista & le/la journaliste & invariable en genre (el/la) \\
el artículo & l'article & attention : \emph{artículo} ≠ « article » (loi) → \emph{artículo de ley} \\
el titular & le titre (une) & aussi adjectif : \emph{socio titular} (membre titulaire) \\
la redacción & la rédaction (équipe/lieu) & verbe : \emph{redactar} \\
el reportaje & le reportage & \emph{reportaje de investigación} = reportage d'enquête \\
la entrevista & l'interview, l'entretien & \emph{entrevistar} (verbe) \\
\hline
\end{tabularx}
\end{center}

\begin{propriete}[Familles de mots : \emph{prensa / prensar}]
\emph{Prensa} (nom féminin) vient du verbe \emph{prensar} (presser, comprimer) → \emph{la imprenta} (l'imprimerie), \emph{el prensista} (l'imprimeur). Ne pas confondre avec \emph{prenda} (vêtement, gage) ni \emph{prenso} (participe de \emph{prensar}).
\end{propriete}

\begin{exemplebox}[Verbes clés du champ 1]
\begin{itemize}
\item \esp{El periodista \textbf{informa} con rigor.} → Le journaliste informe avec rigueur.
\item \esp{El diario \textbf{publica} la investigación.} → Le journal publie l'enquête.
\item \esp{Los medios \textbf{difunden} la cultura de paz.} → Les médias diffusent la culture de paix.
\item \esp{La entrevista \textbf{da la palabra} a las víctimas.} → L'interview donne la parole aux victimes.
\end{itemize}
\end{exemplebox}

\begin{exoresolu}[Exercice : Compléter avec le vocabulaire du champ 1]
\textbf{Énoncé :} Complétez les phrases suivantes en choisissant le mot adéquat parmi : \emph{redacción, titular, reportaje, entrevista, periodista, artículo, prensa, periódico}.

1. La \_\_\_\_\_\_ de ayer traía un \_\_\_\_\_\_ impactante sobre la corrupción. \\
2. El \_\_\_\_\_\_ entrevistó al ministro en una \_\_\_\_\_\_ exclusiva. \\
3. En la \_\_\_\_\_\_ trabajan veinte \_\_\_\_\_\_. \\
4. El \_\_\_\_\_\_ de opinión critica la \_\_\_\_\_\_ sensacionalista.

\tcblower
\textbf{Solution détaillée :}
1. \emph{periódico} (journal), \emph{reportaje} (reportage). \\
2. \emph{periodista} (journaliste), \emph{entrevista} (interview). \\
3. \emph{redacción} (rédaction), \emph{periodistas} (journalistes). \\
4. \emph{artículo} (article), \emph{prensa} (presse).
\end{exoresolu}

\courssub{Champ 2 — Le quatrième pouvoir et la liberté}

\begin{center}
\begin{tabularx}{\linewidth}{|X|X|X|}
\hline
\rowcolor{poporangeL} \textbf{Español} & \textbf{Français} & \textbf{Remarque} \\
\hline
el cuarto poder & le quatrième pouvoir & métaphore politique (après exécutif, législatif, judiciaire) \\
la libertad de prensa & la liberté de la presse & \emph{libertad de expresión} = liberté d'expression \\
la independencia & l'indépendance & \emph{independiente} (adj.) ; \emph{independizarse} (s'émanciper) \\
la imparcialidad & l'impartialité & \emph{imparcial} (adj.) ; antonyme : \emph{parcial} \\
la opinión pública & l'opinion publique & voir définition ci-dessous \\
la influencia & l'influence & \emph{ejercer influencia sobre} = exercer une influence sur \\
\hline
\end{tabularx}
\end{center}

\begin{definition}[opinión pública]
Ensemble des opinions, attitudes et croyances d'une population sur un sujet d'intérêt général, mesuré ou non par des sondages. La presse libre est le principal vecteur de formation de l'\emph{opinión pública}.
\end{definition}

\begin{attention}[Confusion \emph{mediador} vs \emph{medios}]
\begin{itemize}
\item \cle{el mediador} / \cle{la mediadora} = le médiateur / la médiatrice (personne qui facilite le dialogue).
\item \cle{los medios (de comunicación)} = les médias (presse, radio, TV, internet).
\item \cle{el medio} = le moyen, le milieu, la manière.
\end{itemize}
\emph{Exemple :} \esp{Los medios de comunicación entrevistaron al mediador del conflicto.} (Les médias ont interviewé le médiateur du conflit.)
\end{attention}

\begin{exemplebox}[Structures avec \emph{influencia} y \emph{opinión pública}]
\begin{itemize}
\item \esp{La prensa libre \textbf{forma la opinión pública}.} → La presse libre forme l'opinion publique.
\item \esp{Los grandes grupos \textbf{ejercen gran influencia sobre} la agenda política.} → Les grands groupes exercent une grande influence sur l'agenda politique.
\item \esp{La \textbf{independencia} editorial garantiza la \textbf{imparcialidad}.} → L'indépendance éditoriale garantit l'impartialité.
\end{itemize}
\end{exemplebox}

\courssub{Champ 3 — Les menaces : censure, désinformation, manipulation}

\begin{center}
\begin{tabularx}{\linewidth}{|X|X|X|}
\hline
\rowcolor{poppurpleL} \textbf{Español} & \textbf{Français} & \textbf{Remarque} \\
\hline
la censura & la censure & \emph{censurar} (verbe), \emph{censurado} (participe/adjectif) \\
la desinformación & la désinformation & \emph{desinformar} (verbe), \emph{desinformado} (mal informé) \\
las noticias falsas / \emph{fake news} & les fausses nouvelles & \emph{noticia falsa} (sing.) ; attention faux ami \emph{noticia} \\
la manipulación & la manipulation & \emph{manipular} (verbe), \emph{manipulador/a} (adj./nom) \\
la propaganda & la propagande & \emph{propagandístico} (adj.) \\
silenciar & réduire au silence & \emph{silenciar a los periodistas} = museler les journalistes \\
\hline
\end{tabularx}
\end{center}

\begin{propriete}[Famille morphologique : \emph{censura / censurar / censurado}]
\begin{itemize}
\item \emph{la censura} (nom f.) : l'acte ou l'institution de censure.
\item \emph{censurar} (verbe régulier -ar) : censurer, interdire.
\item \emph{censurado, -a} (participe/adjectif) : censuré, interdit.
\item \emph{el censor} (nom m.) : le censeur.
\end{itemize}
\emph{Exemple :} \esp{El gobierno \textbf{censuró} el artículo; quedó \textbf{censurado} por la \textbf{censura} militar.}
\end{propriete}

\begin{attention}[Faux amis fréquents]
\begin{itemize}
\item \cle{una noticia} = une nouvelle, une information journalistique (≠ « une notice » = \emph{un aviso}, \emph{una nota}, \emph{un prospecto}).
\item \cle{actualmente} = actuellement, de nos jours (≠ « en fait, en réalité » = \emph{en realidad}, \emph{de hecho}).
\item \cle{la propaganda} = la propagande (souvent politique) ; \emph{publicidad} = la publicité (commerciale).
\item \cle{un reportaje} = un reportage ; \emph{un informe} = un rapport.
\end{itemize}
\end{attention}

\begin{exemplebox}[Verbes du champ 3]
\begin{itemize}
\item \esp{El régimen \textbf{censura} las noticias incómodas.} → Le régime censure les nouvelles gênantes.
\item \esp{Las redes \textbf{difunden desinformación} sin control.} → Les réseaux diffusent de la désinformation sans contrôle.
\item \esp{El poder \textbf{manipula} la opinión pública con \textbf{propaganda}.} → Le pouvoir manipule l'opinion publique avec de la propagande.
\item \esp{Intentan \textbf{silenciar} a los periodistas críticos.} → Ils tentent de réduire au silence les journalistes critiques.
\end{itemize}
\end{exemplebox}

\begin{exoresolu}[Exercice : Choisir le bon mot (faux amis)]
\textbf{Énoncé :} Remplacez l'anglicisme ou le mot incorrect par le terme espagnol exact.

1. El periódico publicó una \emph{notice} sobre el accidente. \\
2. \emph{Actualmente}, yo creo que la prensa es libre. (sens : « En fait, je crois… ») \\
3. El gobierno usa \emph{publicidad} política para manipular. \\
4. Los \emph{fake news} son un problema grave.

\tcblower
\textbf{Solution :}
1. \emph{notice} → \textbf{noticia} (information). \\
2. \emph{Actualmente} (actuellement) → \textbf{En realidad} / \textbf{De hecho} (en fait). \\
3. \emph{publicidad} (publicité commerciale) → \textbf{propaganda} (propagande politique). \\
4. \emph{fake news} → \textbf{noticias falsas} / \textbf{desinformación}.
\end{exoresolu}

\courssub{Champ 4 — La paix, la médiation et la réconciliation}

\begin{center}
\begin{tabularx}{\linewidth}{|X|X|X|}
\hline
\rowcolor{popgreenL} \textbf{Español} & \textbf{Français} & \textbf{Remarque} \\
\hline
el mediador / la mediadora & le médiateur / la médiatrice & agent de \emph{mediación} \\
la mediación & la médiation & processus de résolution pacifique \\
el diálogo & le dialogue & \emph{dialogar} (verbe) \\
la reconciliación & la réconciliation & \emph{reconciliar(se)} (verbe pronominal) \\
la cultura de paz & la culture de paix & concept UNESCO \\
la convivencia & la coexistence, la vie commune & \emph{convivir} (vivre ensemble) \\
el conflicto & le conflit & \emph{conflicto armado}, \emph{resolución de conflictos} \\
la víctima & la victime & \emph{dar voz a las víctimas} = donner la parole aux victimes \\
\hline
\end{tabularx}
\end{center}

\begin{propriete}[Familles de mots : \emph{paz} et \emph{medio}]
\begin{itemize}
\item \textbf{paz} (f.) → \emph{pacífico/a} (adj. : pacifique), \emph{pacifista} (pacifiste), \emph{apaciguar} (apaiser), \emph{pacificación} (pacification).
\item \textbf{medio} (moyen/milieu) → \emph{los medios de comunicación} (médias), \emph{el mediador} (médiateur), \emph{la mediación} (médiation), \emph{mediar} (médier).
\end{itemize}
\end{propriete}

\begin{exemplebox}[Verbes et expressions du champ 4]
\begin{itemize}
\item \esp{El mediador \textbf{media} entre las partes en \textbf{conflicto}.} → Le médiateur médie entre les parties en conflit.
\item \esp{La prensa \textbf{promueve} la \textbf{cultura de paz} y la \textbf{convivencia}.} → La presse promeut la culture de paix et la coexistence.
\item \esp{Es necesario \textbf{reconciliar} a las comunidades tras la guerra.} → Il faut réconcilier les communautés après la guerre.
\item \esp{El periodismo de paz \textbf{da voz a las víctimas} y \textbf{educa} para la tolerancia.} → Le journalisme de paix donne la parole aux victimes et éduque à la tolérance.
\end{itemize}
\end{exemplebox}

\begin{textebox}[Texte original : « El periodismo de paz en Camerún »]
En \cle{Yaoundé}, una joven \cle{periodista} llamada Amina trabaja en una \cle{redacción} comunitaria. Su \cle{reportaje} semanal no busca el \cle{titular} sensacionalista, sino \cle{dar voz a las víctimas} de los conflictos por la tierra entre \cle{agricultores} y \cle{ganaderos}. Amina sabe que la \cle{desinformación} en \cle{WhatsApp} alimenta la \cle{manipulación} y el \cle{odio}. Por eso, cada programa de radio organiza un \cle{diálogo} entre líderes locales; ella actúa como \cle{mediadora} y \cle{promueve} la \cle{reconciliación}. « La \cle{prensa} libre —dice— no es solo \cle{informar}, es \cle{educar} para la \cle{convivencia} y la \cle{cultura de paz}. » Su \cle{independencia} editorial le permite \cle{denunciar} la \cle{censura} y la \cle{propaganda} de los poderosos. Así, el \cle{cuarto poder} se convierte en \cle{herramienta} de \cle{mediación} y no en arma de \cle{conflicto}.
\tcblower
\emph{Glossaire sélectif :} \cle{ganadero} = éleveur ; \cle{alimenta} (verbe \emph{alimentar}) = alimente ; \cle{odio} = haine ; \cle{poderoso} = puissant ; \cle{herramienta} = outil ; \cle{denunciar} = dénoncer.
\end{textebox}

\begin{exercice}[Questions sur le texte « El periodismo de paz en Camerún »]
\begin{enumerate}
\item Identifiez cinq noms du champ lexical « presse » et cinq du champ « paix/médiation ».
\item Trouvez trois verbes conjugués au présent de l'indicatif et donnez leur infinitif.
\item Expliquez en français l'expression \esp{« dar voz a las víctimas »} et son rôle dans le journalisme de paix.
\item Traduisez en espagnol : « La désinformation sur WhatsApp alimente la haine. »
\item Selon le texte, en quoi l'indépendance éditoriale protège-t-elle la culture de paix ?
\end{enumerate}
\end{exercice}

\begin{savaistu}[Du « cuarto poder » au « quinto poder »]
L'expression \emph{cuarto poder} (quatrième pouvoir) date du \textsc{xviii}\textsuperscript{e} siècle (Edmund Burke, 1787) pour désigner la presse comme contre-pouvoir face aux trois pouvoirs classiques. Aujourd'hui, les théoriciens de la communication parlent de \emph{quinto poder} (cinquième pouvoir) pour qualifier les \emph{redes sociales} et le \emph{periodismo ciudadano} (journalisme citoyen) : les citoyens deviennent producteurs d'information, ce qui renouvelle les enjeux de \emph{veracidad} (véracité), \emph{ética} (éthique) et \emph{responsabilidad} (responsabilité). Ce concept est un atout pour l'oral du Bac : il montre votre culture médiatique actualisée.
\end{savaistu}

\begin{popfigure}
\begin{center}\tcbox[colback=popink,colframe=popink,coltext=white,arc=9pt,boxsep=4pt,left=6pt,right=6pt]{\bfseries\small medios y paz}\end{center}
\vspace{-2pt}
\begin{tcbraster}[raster columns=2,raster equal height=rows,raster column skip=6pt,raster row skip=6pt,enhanced,blankest]
\begin{tcolorbox}[enhanced,colback=white,colframe=poporange,boxrule=0.9pt,arc=3pt,title={prensa},fonttitle=\bfseries\scriptsize,coltitle=white,colbacktitle=poporange,left=4pt,right=4pt,top=2pt,bottom=2pt,boxsep=1pt,toptitle=1.5pt,bottomtitle=1.5pt,halign=left]
\begin{itemize}[nosep,leftmargin=*,itemsep=1pt,font=\scriptsize]
\item periódico
\item periodista
\item reportaje
\item entrevista
\item redacción
\end{itemize}
\end{tcolorbox}
\begin{tcolorbox}[enhanced,colback=white,colframe=popgreen,boxrule=0.9pt,arc=3pt,title={libertad},fonttitle=\bfseries\scriptsize,coltitle=white,colbacktitle=popgreen,left=4pt,right=4pt,top=2pt,bottom=2pt,boxsep=1pt,toptitle=1.5pt,bottomtitle=1.5pt,halign=left]
\begin{itemize}[nosep,leftmargin=*,itemsep=1pt,font=\scriptsize]
\item cuarto poder
\item independencia
\item imparcialidad
\item opinión pública
\item influencia
\end{itemize}
\end{tcolorbox}
\begin{tcolorbox}[enhanced,colback=white,colframe=poppurple,boxrule=0.9pt,arc=3pt,title={amenazas},fonttitle=\bfseries\scriptsize,coltitle=white,colbacktitle=poppurple,left=4pt,right=4pt,top=2pt,bottom=2pt,boxsep=1pt,toptitle=1.5pt,bottomtitle=1.5pt,halign=left]
\begin{itemize}[nosep,leftmargin=*,itemsep=1pt,font=\scriptsize]
\item censura
\item desinformación
\item noticias falsas
\item manipulación
\item propaganda
\end{itemize}
\end{tcolorbox}
\begin{tcolorbox}[enhanced,colback=white,colframe=poppink,boxrule=0.9pt,arc=3pt,title={mediación},fonttitle=\bfseries\scriptsize,coltitle=white,colbacktitle=poppink,left=4pt,right=4pt,top=2pt,bottom=2pt,boxsep=1pt,toptitle=1.5pt,bottomtitle=1.5pt,halign=left]
\begin{itemize}[nosep,leftmargin=*,itemsep=1pt,font=\scriptsize]
\item mediador
\item diálogo
\item reconciliación
\item cultura de paz
\item convivencia
\end{itemize}
\end{tcolorbox}
\end{tcbraster}

\legende{Carte mentale : les quatre champs lexicaux « medios y paz » (branches colorées).}
\end{popfigure}

\begin{aretenir}[Synthèse lexicale — Mots-clés à maîtriser pour le Bac]
\begin{itemize}
\item \cle{prensa, periódico, periodista, reportaje, entrevista, redacción, titular, artículo}
\item \cle{cuarto poder, libertad de prensa, independencia, imparcialidad, opinión pública, influencia}
\item \cle{censura, desinformación, noticias falsas, manipulación, propaganda, silenciar}
\item \cle{mediador, mediación, diálogo, reconciliación, cultura de paz, convivencia, conflicto, víctima}
\item Verbes : \cle{informar, denunciar, publicar, difundir, censurar, mediar, reconciliar(se), promover, educar}
\item Expressions : \cle{dar la palabra a alguien, ejercer influencia sobre, silenciar a los periodistas, dar voz a las víctimas}
\item Faux amis : \cle{noticia ≠ notice}, \cle{actualmente ≠ en fait}, \cle{propaganda ≠ publicidad}
\item Familles : \cle{prensa/prensar}, \cle{censura/censurar/censurado}, \cle{paz/pacífico/apaciguar}, \cle{medio/medios/mediador/mediación}
\end{itemize}
\end{aretenir}

\coursec{Méthode du commentaire de texto}

Après avoir constitué, dans la section précédente, une solide boîte à outils lexicale (\esp{cuarto poder}, \esp{libertad de prensa}, \esp{censura}, \esp{opinión pública}, \esp{mediador}, \esp{reconciliación}...), il faut maintenant apprendre à \cle{mobiliser ce vocabulaire dans un exercice rédigé}. Au baccalauréat, l'épreuve phare est le \cle{comentario de texto} : un texte argumentatif ou explicatif vous est proposé, et vous devez en rendre compte de façon organisée, entièrement en espagnol. Cette section vous donne une méthode en sept étapes, des phrases toutes faites, un modèle d'introduction et les pièges à éviter.

\courssub{Qu'est-ce qu'un commentaire de texte ?}

\begin{definition}[Le comentario de texto]
Le \cle{commentaire de texte} (\esp{comentario de texto}) est un exercice écrit qui consiste à : (1) présenter un texte (nature, auteur, source, thème), (2) dégager son idée générale et la thèse de l'auteur, (3) analyser ses idées principales et ses procédés argumentatifs, (4) donner son avis personnel argumenté, (5) conclure par un bilan et une ouverture. Il se rédige \textbf{entièrement en espagnol}, en paragraphes rédigés, sans plan apparent avec des chiffres.
\end{definition}

\begin{attention}[Règle d'or du Bac]
Au baccalauréat, le commentaire se rédige \textbf{entièrement en espagnol}, de la première à la dernière ligne. Toute phrase en français est pénalisée. Il faut aussi éviter les calques du français : on ne dit pas \esp{el texto habla sobre la paz} (calque de « le texte parle de »), mais \esp{el texto trata de la paz} ou \esp{el texto aborda el tema de la paz}. De même, évitez \esp{según yo} : dites \esp{en mi opinión} ou \esp{a mi juicio}.
\end{attention}

\courssub{Les sept étapes du commentaire}

Voici d'abord une vue d'ensemble de la démarche, sous forme d'escalier : chaque marche vous rapproche de la rédaction finale.

\begin{popfigure}
\begin{tikzpicture}[font=\small, node distance=0.5cm and 1cm]
  \node[draw=popblue, fill=popblueL, rounded corners=2pt, align=center, minimum width=3.1cm, minimum height=1.1cm] (s1) {1. Leer dos veces\\y subrayar};
  \node[draw=popblue, fill=popblueL, rounded corners=2pt, align=center, minimum width=3.1cm, minimum height=1.1cm] (s2) [below right=of s1] {2. Introducción\\y problemática};
  \node[draw=popblue, fill=popblueL, rounded corners=2pt, align=center, minimum width=3.1cm, minimum height=1.1cm] (s3) [below right=of s2] {3. Idea general\\y tesis};
  \node[draw=popblue, fill=popblueL, rounded corners=2pt, align=center, minimum width=3.1cm, minimum height=1.1cm] (s4) [below right=of s3] {4. Análisis de las\\y ideas principales};
  \node[draw=popblue, fill=popblueL, rounded corners=2pt, align=center, minimum width=3.1cm, minimum height=1.1cm] (s5) [below right=of s4] {5. Procedimientos\\argumentativos};
  \node[draw=popblue, fill=popblueL, rounded corners=2pt, align=center, minimum width=3.1cm, minimum height=1.1cm] (s6) [below right=of s5] {6. Opinión personal\\argumentada};
  \node[draw=popblue, fill=popblueL, rounded corners=2pt, align=center, minimum width=3.1cm, minimum height=1.1cm] (s7) [below right=of s6] {7. Conclusión\\y apertura};
  \draw[-{Stealth}, thick, poporange] (s1) -- (s2);
  \draw[-{Stealth}, thick, poporange] (s2) -- (s3);
  \draw[-{Stealth}, thick, poporange] (s3) -- (s4);
  \draw[-{Stealth}, thick, poporange] (s4) -- (s5);
  \draw[-{Stealth}, thick, poporange] (s5) -- (s6);
  \draw[-{Stealth}, thick, poporange] (s6) -- (s7);
  \node[align=center, popdark, below=1cm of s4, xshift=-3cm] {De la lectura a la conclusión : siete peldaños};
\end{tikzpicture}
\legende{Schéma en escalier : les sept étapes du commentaire de texte.}
\end{popfigure}

\textbf{Étape 1 — Lire deux fois le texte et souligner les mots-clés.}\par
La première lecture se fait sans stylo, pour saisir le sens global. La deuxième lecture est active : soulignez les \cle{mots-clés} (noms, verbes d'opinion, connecteurs), entourez les mots inconnus et essayez de les deviner par le contexte. Repérez aussi les marques du genre : titre, signature, source, date.

\textbf{Étape 2 — Rédiger l'introduction.}\par
L'introduction présente le texte en deux ou trois phrases : sa \cle{nature} (article d'opinion, éditorial, reportage, discours...), son \cle{auteur} s'il est connu, sa \cle{source} (journal, revue, site), son \cle{thème}, puis elle annonce la \cle{problématique}, c'est-à-dire la question à laquelle le texte répond.

\begin{methode}[Rédiger l'introduction type]
\begin{enumerate}
\item Identifier la nature du document : \esp{Es un artículo de opinión / un editorial / un discurso...}
\item Nommer l'auteur et la source : \esp{publicado en el periódico... / escrito por...}
\item Présenter le thème : \esp{que trata del papel de la prensa en la promoción de la paz.}
\item Formuler la problématique : \esp{La cuestión que plantea el texto es la siguiente : ¿puede la prensa contribuir a la reconciliación de un país?}
\end{enumerate}
\textbf{Modèle complet à adapter :} \esp{El texto que vamos a comentar es un artículo de opinión titulado «El cuarto poder al servicio de la paz», publicado en un periódico nacional, que trata del papel de la prensa en la promoción de la paz. El autor se pregunta si los medios de comunicación pueden ser mediadores en los conflictos sociales.} « Le texte que nous allons commenter est un article d'opinion intitulé "Le quatrième pouvoir au service de la paix", publié dans un journal national, qui traite du rôle de la presse dans la promotion de la paix. L'auteur se demande si les médias peuvent être des médiateurs dans les conflits sociaux. »
\end{methode}

\textbf{Étape 3 — Dégager l'idée générale et la thèse de l'auteur.}\par
L'\cle{idée générale} est le sujet du texte en une phrase ; la \cle{thèse} est la position que défend l'auteur. Ne les confondez pas : le sujet est « la presse et la paix » ; la thèse est « la presse doit éduquer les citoyens à la paix au lieu d'alimenter les conflits ».

\begin{exemplebox}[Formules pour présenter la thèse]
\esp{El autor defiende la idea de que la prensa libre es un pilar de la democracia.} « L'auteur défend l'idée que la presse libre est un pilier de la démocratie. »\par
\esp{El autor denuncia la censura.} « L'auteur dénonce la censure. »\par
\esp{Según el texto, la prensa educa a los ciudadanos.} « Selon le texte, la presse éduque les citoyens. »\par
\esp{Para el autor, los periodistas deben ser independientes.} « Pour l'auteur, les journalistes doivent être indépendants. »
\end{exemplebox}

\textbf{Étape 4 — Analyser les idées principales.}\par
Construisez un plan en \cle{deux ou trois parties} (thèse / antithèse, ou causes / conséquences / solutions). Chaque partie doit être \textbf{justifiée par une citation courte} du texte, introduite par une formule de renvoi : \esp{como dice el autor}, \esp{según el texto}, \esp{el autor afirma que...}. Une citation sans analyse ne vaut rien : commentez toujours ce que vous citez.

\textbf{Étape 5 — Étudier les procédés argumentatifs.}\par
Un texte argumentatif convainc grâce à des procédés qu'il faut savoir nommer :

\begin{center}
\begin{tabularx}{\linewidth}{|X|X|X|}
\hline
\rowcolor{popblueL} \textbf{Procedimiento} & \textbf{Ejemplo español} & \textbf{Effet / traduction} \\
\hline
Ejemplo (exemple) & \esp{Por ejemplo, en muchos países...} & Illustre la thèse : « par exemple » \\
\hline
Comparación & \esp{La prensa es como un puente entre el pueblo y el poder.} & Rend l'idée concrète : « la presse est comme un pont » \\
\hline
Pregunta retórica & \esp{¿Quién vigilará a los poderosos si no hay periodistas?} & Fait réfléchir le lecteur : question rhétorique \\
\hline
Léxico valorativo & \esp{prensa libre / prensa manipuladora} & Lexique mélioratif / péjoratif \\
\hline
Repetición & \esp{sin prensa libre, sin democracia} & Martèle l'idée : « sans presse libre, pas de démocratie » \\
\hline
Conectores & \esp{sin embargo, por lo tanto, además} & Structurent l'argumentation : « cependant, donc, de plus » \\
\hline
\end{tabularx}
\end{center}

\textbf{Étape 6 — Donner son avis personnel argumenté.}\par
C'est le moment de vous engager : êtes-vous d'accord avec l'auteur ? Votre avis doit être \cle{argumenté}, jamais gratuit. Utilisez \esp{Estoy de acuerdo con el autor porque...} ou \esp{No estoy de acuerdo porque...}, et appuyez-vous sur votre expérience (la presse camerounaise, les radios locales, les réseaux sociaux...).

\begin{propriete}[Exprimer un désaccord nuancé : no creo que + subjonctif]
Après \esp{no creo que}, \esp{no pienso que}, \esp{no me parece que}, le verbe de la subordonnée se met au \textbf{subjonctif}, car on nie la réalité du fait : \esp{No creo que la prensa \textbf{sea} siempre imparcial.} « Je ne crois pas que la presse soit toujours impartiale. »\par
En revanche, à la forme affirmative, on utilise l'\textbf{indicatif} : \esp{Creo que la prensa \textbf{es} necesaria.} « Je crois que la presse est nécessaire. »\par
\textbf{Attention :} le français utilise aussi le subjonctif après « je ne crois pas que », donc la difficulté n'est pas le choix du mode mais la \textbf{forme} du subjonctif espagnol : \esp{sea} (ser), \esp{tenga} (tener), \esp{pueda} (poder), \esp{diga} (decir).
\end{propriete}

\textbf{Étape 7 — Conclure : bilan et ouverture.}\par
La conclusion résume l'essentiel en une ou deux phrases (\esp{En conclusión, este texto muestra que...}) puis \cle{ouvre} sur un sujet voisin : le rôle des réseaux sociaux, la formation des journalistes, la presse en Guinée équatoriale ou au Cameroun. Formule utile : \esp{Esta reflexión nos lleva a otra pregunta : ¿qué papel juegan las redes sociales en la cultura de paz?}

\courssub{Phrases toutes faites pour le commentaire}

\begin{aretenir}[Boîte à phrases du comentario]
\begin{itemize}
\item Présenter : \esp{El texto trata de... / El texto aborda el tema de...}
\item Citer l'auteur : \esp{El autor defiende la idea de que... / Según el autor... / Como afirma el texto...}
\item Analyser : \esp{Este ejemplo ilustra que... / Con esta comparación, el autor quiere decir que...}
\item Donner son avis : \esp{En mi opinión... / A mi juicio... / Estoy de acuerdo porque... / No creo que... + subjuntivo}
\item Conclure : \esp{En conclusión... / Para terminar... / Esta reflexión nos lleva a preguntarnos...}
\end{itemize}
\end{aretenir}

\courssub{Un mini-texte pour s'entraîner}

\begin{textebox}[La voz que calma (texto original)]
En los días difíciles que atraviesa nuestra sociedad, la prensa tiene una responsabilidad enorme. Un periodista no es un simple testigo: es un mediador entre los hechos y los ciudadanos. Cuando una radio local difunde rumores sin verificarlos, enciende la mecha del conflicto; cuando, por el contrario, informa con rigor, apaga el fuego antes de que arda.

¿Quién no recuerda emisiones que han reconciliado a vecinos enfrentados? La palabra, bien usada, es un puente; mal usada, es un arma. Por eso la libertad de prensa no significa libertad para mentir: significa libertad para buscar la verdad y decirla sin miedo.

Sin embargo, esta misión exige independencia. Un periodista comprado por el poder ya no es un vigilante, sino un sirviente. La censura mata la verdad, pero la corrupción la envenena lentamente.

En conclusión, formar periodistas honestos es construir la paz. Cada artículo riguroso es un ladrillo en el edificio de la reconciliación nacional.
\end{textebox}

\textbf{Glossaire :} \esp{testigo} : témoin ; \esp{difundir rumores} : répandre des rumeurs ; \esp{encender la mecha} : allumer la mèche ; \esp{apagar el fuego} : éteindre le feu ; \esp{comprado por el poder} : acheté par le pouvoir ; \esp{vigilante} : sentinelle, gardien ; \esp{sirviente} : serviteur ; \esp{envenenar} : empoisonner ; \esp{ladrillo} : brique ; \esp{edificio} : édifice.

\textbf{Questions de préparation au commentaire :}
\begin{enumerate}
\item ¿Cuál es la naturaleza y el tema de este texto?
\item ¿Qué tesis defiende el autor sobre el papel del periodista?
\item Señala dos comparaciones y explica su efecto.
\item ¿Estás de acuerdo con la frase «la palabra, mal usada, es un arma»? Justifica tu respuesta.
\end{enumerate}

\courssub{Exercice résolu et pièges à éviter}

\begin{exoresolu}[Rédiger une introduction et une opinion]
\textbf{Énoncé.} À partir du texte \esp{La voz que calma} ci-dessus : 1) rédigez une introduction complète (4 phrases) ; 2) exprimez un désaccord nuancé avec \esp{no creo que} + subjonctif.
\tcblower
\textbf{Solution détaillée.}
1) \esp{El texto que vamos a comentar es un artículo de opinión titulado «La voz que calma». Trata del papel de la prensa en la promoción de la cultura de paz. El autor defiende la idea de que el periodista es un mediador que puede encender o apagar los conflictos. La cuestión que plantea es la siguiente: ¿cómo puede la prensa servir a la reconciliación nacional?}\par
2) \esp{En mi opinión, el autor tiene razón en general; sin embargo, no creo que todos los periodistas sean mediadores honestos, porque algunos difunden rumores por interés. Tampoco creo que la libertad de prensa sea suficiente: hace falta también la formación de los ciudadanos.}\par
Remarquez : après \esp{no creo que}, les verbes sont au subjonctif (\esp{sean}, \esp{sea}) ; la nuance est marquée par \esp{sin embargo} et \esp{tampoco}.
\end{exoresolu}

\begin{attention}[Les erreurs fréquentes des francophones]
\begin{itemize}
\item \esp{El texto habla sobre la paz.} (calque) → dites \esp{El texto trata de la paz.}
\item \esp{Estoy de acuerdo con que...} → dites \esp{Estoy de acuerdo en que...} ou simplement \esp{Estoy de acuerdo porque...}
\item \esp{No creo que la prensa \emph{es} imparcial.} → subjonctif obligatoire : \esp{no creo que \emph{sea} imparcial.}
\item \esp{Según yo...} → \esp{En mi opinión...} ou \esp{A mi juicio...}
\item Oublier de citer le texte : un commentaire sans citations est un commentaire vide.
\item Rédiger en français ou mélanger les deux langues : tout doit être en espagnol.
\end{itemize}
\end{attention}

\begin{exercice}[À vous de jouer]
Reprenez le texte \esp{La voz que calma} et rédigez, en espagnol, le début d'un commentaire : introduction complète, idée générale, thèse de l'auteur, puis une première partie d'analyse avec deux citations commentées. Terminez par une phrase d'opinion personnelle utilisant \esp{no creo que} + subjonctif.
\end{exercice}

Vous disposez maintenant de la méthode complète : dans la section suivante, nous l'appliquerons pas à pas au texte support \esp{El cuarto poder al servicio de la paz} pour produire un commentaire modèle.

\coursec{Commentaire modèle}

Dans la section précédente, nous avons découvert la \cle{méthode du commentaire de texte} : introduction (présentation, thème, problématique), développement en deux ou trois parties avec citations, avis personnel argumenté et conclusion avec ouverture. Il est temps de passer de la théorie à la pratique. Nous allons lire ensemble un \cle{commentaire modèle} complet, rédigé à partir du texte support \esp{El cuarto poder al servicio de la paz}, puis nous l'analyserons pas à pas pour que vous puissiez reproduire cette démarche le jour de l'examen.

\courssub{Le commentaire modèle intégral}

Lisez d'abord le commentaire en entier, sans dictionnaire, pour saisir sa logique générale. Repérez ensuite les connecteurs et les citations.

\begin{textebox}[Comentario modelo : \esp{¿Puede la prensa contribuir a la cultura de paz?}]
\textbf{Introducción.} El texto que vamos a comentar es un artículo de opinión titulado \emph{El cuarto poder al servicio de la paz}. El autor defiende la idea de que los medios de comunicación, llamados el cuarto poder, pueden contribuir a la promoción de la cultura de paz. El tema es, por tanto, el papel de la prensa en la sociedad. La problemática que plantea el texto es la siguiente: ¿puede la prensa contribuir a la cultura de paz?

\textbf{Primera parte.} En primer lugar, el autor presenta la prensa como un poder de control: \emph{« informa y denuncia las injusticias »}. Esta cita muestra que el periodista tiene dos funciones esenciales: dar información objetiva a los ciudadanos y señalar los abusos de los poderosos. Gracias a la libertad de prensa, los medios vigilan al gobierno y protegen la democracia. Sin embargo, el autor recuerda que este poder exige independencia: sin ella, la censura destruye la credibilidad de los periodistas.

\textbf{Segunda parte.} En segundo lugar, afirma que el periodista puede ser un mediador que favorece la reconciliación. Según el texto, los medios \emph{« dan la palabra a todas las partes del conflicto »}. Esta idea es muy importante: al escuchar a las víctimas y a los adversarios, la prensa crea un espacio de diálogo y reduce el odio. Así, el cuarto poder no solo informa, sino que también construye la paz.

\textbf{Opinión personal.} Estoy de acuerdo con esta tesis, aunque pienso que hay que matizarla. En mi país, \esp{Camerún}, los medios han ayudado a calmar los conflictos dando la palabra a las víctimas y difundiendo mensajes de unidad. Sin embargo, algunos medios sensacionalistas pueden también encender las tensiones. En Colombia, por ejemplo, la prensa acompañó el proceso de paz explicando los acuerdos a la población, lo que demuestra que los medios \textbf{pueden} ser verdaderos actores de reconciliación.

\textbf{Conclusión.} En conclusión, el texto demuestra que la prensa es a la vez un poder de control y una mediadora de paz. No obstante, esta reflexión debe ampliarse hoy a las redes sociales: frente a las noticias falsas, la educación de los ciudadanos es más necesaria que nunca.
\end{textebox}

\textbf{Glossaire.} \esp{denunciar} : dénoncer ; \esp{los abusos} : les abus ; \esp{vigilar} : surveiller ; \esp{matizar} : nuancer ; \esp{encender las tensiones} : attiser les tensions ; \esp{los acuerdos} : les accords ; \esp{las noticias falsas} : les fausses informations (\emph{fake news}) ; \esp{ampliarse} : s'étendre, être élargie.

\begin{aretenir}[Questions de vérification]
\begin{enumerate}
\item ¿Cuál es la problemática del comentario?
\item ¿Qué dos funciones de la prensa se presentan en las dos partes del desarrollo?
\item ¿Qué ejemplos utiliza el alumno para expresar su opinión personal?
\item ¿Sobre qué tema se abre la conclusión?
\end{enumerate}
\end{aretenir}

\courssub{Analyse guidée du modèle : l'introduction}

Relisons la première phrase : \esp{El texto que vamos a comentar es un artículo de opinión titulado...}. Observez la structure : on présente d'abord la \cle{nature du texte} (article d'opinion), puis son \cle{titre}, puis la \cle{thèse de l'auteur}, puis le \cle{thème}, et enfin la \cle{problématique} sous forme de question directe avec les signes ¿ ... ?.

\begin{propriete}[Structure de l'introduction de commentaire]
L'introduction suit toujours quatre étapes, dans cet ordre :
\begin{enumerate}
\item \textbf{Présentation} : \esp{El texto que vamos a comentar es...} (nature, titre, auteur si connu).
\item \textbf{Thème} : \esp{El tema es...} / \esp{El texto trata de...}.
\item \textbf{Idée principale} : \esp{El autor defiende la idea de que...} (attention : \esp{la idea de que} + indicatif).
\item \textbf{Problématique} : une question : \esp{¿Puede la prensa contribuir a la cultura de paz?}
\end{enumerate}
Exception : si le texte est un dialogue ou un récit, remplacez \esp{artículo de opinión} par \esp{un diálogo}, \esp{un relato}, \esp{una noticia}.
\end{propriete}

\begin{attention}[Piège des francophones]
Ne traduisez pas mot à mot « le texte parle de » par \esp{el texto habla de} : c'est un calque acceptable mais faible. Préférez \esp{el texto trata de} ou \esp{el texto aborda el tema de}. De même, « selon l'auteur » se dit \esp{según el autor} (attention : n'écrivez jamais \emph{selón}, calque phonétique du français ; l'espagnol s'écrit \textbf{según} avec accent sur le \emph{u} et \emph{g}).
\end{attention}

\courssub{Analyse guidée du modèle : le développement et les citations}

Le développement comporte deux parties, annoncées par les connecteurs \esp{en primer lugar} et \esp{en segundo lugar}. Chaque partie suit le même schéma : \textbf{affirmation} puis \textbf{citation entre guillemets} puis \textbf{analyse personnelle de la citation}.

\begin{methode}[Comment citer le texte en trois temps]
\begin{enumerate}
\item \textbf{Introduire la citation} avec un verbe de parole : \esp{el autor afirma que...}, \esp{según el texto...}, \esp{como dice el autor...}.
\item \textbf{Citer entre guillemets} les mots exacts du texte, et si possible indiquer la ligne : \esp{« informa y denuncia las injusticias » (línea 5)}.
\item \textbf{Commenter la citation} : expliquez ce qu'elle prouve, avec \esp{esta cita muestra que...}, \esp{esta idea significa que...}, \esp{con estas palabras, el autor subraya...}.
\end{enumerate}
\end{methode}

\begin{attention}[Le commentaire n'est pas un résumé]
L'erreur la plus grave au Bac est de transformer le commentaire en simple résumé du texte. Règle d'or : \textbf{chaque citation doit être suivie d'une analyse}. Une citation sans commentaire ne rapporte aucun point ; une paraphrase sans citation non plus. Dans le modèle, après la citation \esp{« informa y denuncia las injusticias »}, l'élève ajoute : \esp{Esta cita muestra que el periodista tiene dos funciones esenciales...} : c'est cela, l'analyse.
\end{attention}

\begin{exemplebox}[Citations analysées / citations non analysées]
\begin{itemize}
\item Mauvais (résumé) : \esp{El autor dice que la prensa informa y denuncia. Después habla del diálogo.} « L'auteur dit que la presse informe et dénonce. Ensuite il parle du dialogue. »
\item Bon (analyse) : \esp{El autor afirma que la prensa « informa y denuncia las injusticias »: esta cita revela la doble misión del periodista, vigilante y defensor de los ciudadanos.} « L'auteur affirme que la presse "informe et dénonce les injustices" : cette citation révèle la double mission du journaliste, vigilant et défenseur des citoyens. »
\end{itemize}
\end{exemplebox}

\courssub{Analyse guidée du modèle : l'avis personnel et la conclusion}

L'avis personnel commence par une formule claire : \esp{Estoy de acuerdo con esta tesis, aunque pienso que hay que matizarla.} Remarquez la construction : \esp{estar de acuerdo con} + nom, et la nuance introduite par \esp{aunque} + indicatif (fait réel : \esp{pienso que...}). L'élève appuie ensuite son avis sur deux exemples concrets : un exemple \cle{camerounais} (le rôle des médias dans l'apaisement des conflits) et un exemple \cle{hispanophone} (l'accompagnement du processus de paix en Colombie). C'est exactement ce que l'examinateur attend : un avis \textbf{nuancé, illustré et personnel}.

La conclusion, introduite par \esp{en conclusión}, fait le bilan des deux parties (\esp{un poder de control y una mediadora de paz}) puis s'ouvre sur un sujet d'actualité : les réseaux sociaux et les \esp{noticias falsas}.

\begin{center}
\begin{tabularx}{\linewidth}{|X|X|X|}
\hline
\rowcolor{popblueL} Étape & \esp{Español (modèle)} & Français \\
\hline
Annoncer la 1\textsuperscript{re} partie & \esp{En primer lugar} & En premier lieu \\
\hline
Annoncer la 2\textsuperscript{e} partie & \esp{En segundo lugar} & En second lieu \\
\hline
Citer & \esp{según el texto / el autor afirma que} & selon le texte / l'auteur affirme que \\
\hline
Analyser & \esp{esta cita muestra que} & cette citation montre que \\
\hline
Ajouter & \esp{además / también} & de plus / aussi \\
\hline
Opposer & \esp{sin embargo / no obstante} & cependant / néanmoins \\
\hline
Nuancer & \esp{aunque pienso que hay que matizarla} & bien que je pense qu'il faut la nuancer \\
\hline
Donner son avis & \esp{estoy de acuerdo con / en mi opinión} & je suis d'accord avec / à mon avis \\
\hline
Conclure & \esp{en conclusión / para concluir} & en conclusion / pour conclure \\
\hline
Ouvrir & \esp{esta reflexión debe ampliarse a...} & cette réflexion doit s'étendre à... \\
\hline
\end{tabularx}
\end{center}

\begin{popfigure}
\begin{tikzpicture}[
  col/.style={rectangle, rounded corners=2pt, draw=popline, fill=popcream, align=left, text width=5.6cm, inner sep=4pt, font=\small},
  head/.style={rectangle, rounded corners=2pt, fill=popblue, text=white, align=center, text width=5.6cm, inner sep=4pt, font=\small\bfseries}
]
\node[head] (h1) at (0,0) {Estructura del comentario};
\node[head, right=0.6cm of h1] (h2) {Frases del modelo};
\node[col, below=0.25cm of h1] (a1) {Introducción: naturaleza, tema, problemática};
\node[col, below=0.25cm of h2] (b1) {\esp{El texto que vamos a comentar es... ¿Puede la prensa contribuir a la cultura de paz?}};
\node[col, below=0.25cm of a1] (a2) {Parte 1: la prensa, poder de control};
\node[col, below=0.25cm of b1] (b2) {\esp{En primer lugar... « informa y denuncia las injusticias »}};
\node[col, below=0.25cm of a2] (a3) {Parte 2: la prensa, mediadora de paz};
\node[col, below=0.25cm of b2] (b3) {\esp{En segundo lugar... « dan la palabra a todas las partes »}};
\node[col, below=0.25cm of a3] (a4) {Opinión personal con ejemplos};
\node[col, below=0.25cm of b3] (b4) {\esp{Estoy de acuerdo..., aunque hay que matizarla. En Camerún... En Colombia...}};
\node[col, below=0.25cm of a4] (a5) {Conclusión y apertura};
\node[col, below=0.25cm of b4] (b5) {\esp{En conclusión... frente a las noticias falsas...}};
\draw[-{Stealth}, poporange, thick] (a1.east) -- (b1.west);
\draw[-{Stealth}, poporange, thick] (a2.east) -- (b2.west);
\draw[-{Stealth}, poporange, thick] (a3.east) -- (b3.west);
\draw[-{Stealth}, poporange, thick] (a4.east) -- (b4.west);
\draw[-{Stealth}, poporange, thick] (a5.east) -- (b5.west);
\end{tikzpicture}
\legende{Correspondance entre la structure du commentaire et les phrases du modèle.}
\end{popfigure}

\courssub{Exercice résolu et entraînement}

\begin{exoresolu}[Repérer les éléments du modèle]
Énoncé : dans le commentaire modèle, relevez (a) deux connecteurs d'ordre, (b) deux verbes pour introduire une citation, (c) la formule d'avis personnel, (d) la formule d'ouverture.
\tcblower
Solution : (a) \esp{en primer lugar}, \esp{en segundo lugar} ; (b) \esp{el autor presenta}, \esp{el autor afirma que} (ainsi que \esp{según el texto}) ; (c) \esp{Estoy de acuerdo con esta tesis, aunque pienso que hay que matizarla} ; (d) \esp{esta reflexión debe ampliarse hoy a las redes sociales}. Chacune de ces formules est réutilisable dans n'importe quel commentaire : apprenez-les par cœur.
\end{exoresolu}

\begin{exercice}[À votre tour]
À partir du texte \esp{El cuarto poder al servicio de la paz}, rédigez en dix lignes un développement en deux parties : la première sur la \esp{libertad de prensa}, la seconde sur la \esp{reconciliación}. Obligations : deux citations entre guillemets, chacune suivie d'une analyse commençant par \esp{esta cita muestra que...}, et les connecteurs \esp{en primer lugar} / \esp{en segundo lugar}.
\end{exercice}

\begin{corrige}[Exercice]
Exemple de production attendue : \esp{En primer lugar, el autor subraya la importancia de la libertad de prensa: « sin libertad, no hay democracia ». Esta cita muestra que la independencia de los periodistas es la condición de su credibilidad frente a la censura. En segundo lugar, el texto afirma que los medios « construyen puentes entre los enemigos ». Esta cita significa que la prensa puede favorecer el diálogo y la reconciliación dando la palabra a todas las partes.}
\end{corrige}

\begin{savaistu}[Le commentaire au Bac espagnol Tle A]
Au baccalauréat camerounais, le commentaire de texte en espagnol est évalué selon quatre critères : la \cle{compréhension} du texte (avez-vous bien identifié le thème et la thèse ?), la \cle{richesse du lexique} (employez les mots du champ du quatrième pouvoir : \esp{censura}, \esp{opinión pública}, \esp{mediador}...), la \cle{correction grammaticale} (accords, conjugaisons, accents) et la \cle{pertinence de l'avis personnel}. Un avis illustré par un exemple camerounais ou hispanophone, comme dans notre modèle, fait toujours la différence.
\end{savaistu}

En résumé, le commentaire modèle vous offre une \textbf{charpente} : introduction en quatre temps, deux parties articulées autour de citations analysées, avis nuancé et illustré, conclusion ouverte sur l'actualité. La section suivante vous invitera à mobiliser cette charpente dans une production guidée au service de la culture de la paix.

\coursec{Production guidée : défendre la culture de la paix}

Après avoir lu et commenté le texte \esp{El cuarto poder al servicio de la paz}, nous savons désormais analyser un texte argumentatif sur la presse. Il reste l'étape la plus importante : passer du rôle de lecteur à celui d'\emph{auteur}. Dans cette dernière partie, vous allez produire vous-mêmes, à l'écrit et à l'oral, des textes qui défendent la culture de la paix, en réutilisant le lexique du quatrième pouvoir (\esp{cuarto poder, libertad de prensa, censura, opinión pública, mediador, reconciliación}) et les connecteurs argumentatifs étudiés. L'objectif est double : maîtriser la langue et devenir, par la langue, un acteur de paix dans votre lycée et votre quartier.

\courssub{Observer avant de produire : un modèle de paragraphe argumentatif}

Avant d'écrire, observons comment un élève de Terminale peut organiser un paragraphe d'opinion. Lisez ce texte original, écrit dans le style d'une rédaction scolaire camerounaise.

\begin{textebox}[Un paragraphe modèle : ¿Deben los medios promover la paz?]
En mi opinión, los medios de comunicación cameruneses deben promover activamente la cultura de paz. En primer lugar, la prensa forma la opinión pública: cuando un periódico de Yaoundé publica artículos sobre la reconciliación entre comunidades, los lectores aprenden a dialogar en lugar de pelear. Por ejemplo, durante las elecciones, algunos periodistas invitan a los ciudadanos a respetar los resultados y a evitar la violencia. Además, los medios pueden actuar como mediadores en los conflictos: una radio comunitaria de Douala puede dar la palabra a los dos bandos de una disputa y ayudarlos a entenderse. Sin embargo, esta misión exige una verdadera independencia: un periodista que depende del dinero de un político no puede ser un mediador honesto. Por eso, hay que luchar contra la censura y contra las noticias falsas que dividen a la población. En conclusión, el cuarto poder no debe limitarse a informar: debe construir la paz, porque sin paz no hay escuela, ni mercado, ni futuro para los jóvenes cameruneses.
\end{textebox}

\textbf{Glossaire :} \esp{en mi opinión} = à mon avis ; \esp{en primer lugar} = en premier lieu ; \esp{dar la palabra} = donner la parole ; \esp{los dos bandos} = les deux camps ; \esp{sin embargo} = cependant ; \esp{por eso} = c'est pourquoi ; \esp{noticias falsas} = fausses informations ; \esp{ayudarlos} = les aider (COD masculin pluriel).

Observez la structure : le paragraphe commence par une \cle{opinión}, donne deux \cle{argumentos} chacun suivi d'un \cle{ejemplo}, et se termine par une \cle{conclusión}. C'est exactement le schéma que vous devrez reproduire.

\begin{popfigure}
\begin{tikzpicture}[node distance=0.55cm, every node/.style={font=\small}]
\node[draw=popblue, fill=popblueL, rounded corners, align=center, text width=8.5cm] (op) {\textbf{1. Opinión}\\ \esp{En mi opinión, los medios deben promover la paz.}};
\node[draw=popgreen, fill=popgreenL, rounded corners, align=center, text width=8.5cm, below=of op] (a1) {\textbf{2. Argumento 1 + ejemplo}\\ \esp{En primer lugar, la prensa forma la opinión pública... Por ejemplo...}};
\node[draw=poporange, fill=poporangeL, rounded corners, align=center, text width=8.5cm, below=of a1] (a2) {\textbf{3. Argumento 2 + ejemplo}\\ \esp{Además, los medios pueden actuar como mediadores...}};
\node[draw=poppurple, fill=poppurpleL, rounded corners, align=center, text width=8.5cm, below=of a2] (ccl) {\textbf{4. Conclusión}\\ \esp{En conclusión, el cuarto poder debe construir la paz.}};
\draw[-{Stealth}, thick, popdark] (op) -- (a1);
\draw[-{Stealth}, thick, popdark] (a1) -- (a2);
\draw[-{Stealth}, thick, popdark] (a2) -- (ccl);
\end{tikzpicture}
\legende{Organigramme du paragraphe argumentatif : opinión, argumentos con ejemplos, conclusión.}
\end{popfigure}

\courssub{Les outils linguistiques de l'argumentation}

Pour exprimer votre opinion et réagir à celle des autres, voici les formules indispensables, à comparer avec le français.

\begin{center}
\begin{tabularx}{\linewidth}{|X|X|X|}
\hline
\rowcolor{popblueL} Français & Español & Emploi \\
\hline
À mon avis / Selon moi & \esp{En mi opinión / Desde mi punto de vista} & Introduire son opinion \\
\hline
Je pense que + indicatif & \esp{Pienso que / Creo que + indicativo} & Affirmer une conviction \\
\hline
Il faut + infinitif & \esp{Hay que + infinitivo} & Obligation générale \\
\hline
Je suis d'accord avec toi, mais... & \esp{Estoy de acuerdo contigo, pero...} & Nuancer \\
\hline
Je ne partage pas ton opinion parce que... & \esp{No comparto tu opinión porque...} & Contredire poliment \\
\hline
De plus / Cependant & \esp{Además / Sin embargo} & Ajouter / Opposer \\
\hline
En conclusion & \esp{En conclusión / En definitiva} & Conclure \\
\hline
\end{tabularx}
\end{center}

\begin{propriete}[Conclure un écrit argumentatif]
Pour terminer un paragraphe ou un essai argumentatif en espagnol, on utilise les connecteurs : \esp{En conclusión} (en conclusion), \esp{Para terminar} (pour terminer), \esp{En definitiva} (en définitive, tout compte fait). Ils sont toujours suivis d'une virgule et introduisent la réaffirmation de l'opinion initiale, jamais un argument nouveau. Exemple : \esp{En definitiva, la libertad de prensa es esencial para la democracia.} « En définitive, la liberté de presse est essentielle pour la démocratie. »
\end{propriete}

\begin{exemplebox}[Phrases modèles traduites]
\esp{Los medios deben promover la reconciliación.} « Les médias doivent promouvoir la réconciliation. »\\
\esp{Hay que luchar contra las noticias falsas.} « Il faut lutter contre les fausses informations. »\\
\esp{Los periodistas deben respetar la verdad.} « Les journalistes doivent respecter la vérité. »\\
\esp{Hay que proteger la libertad de prensa.} « Il faut protéger la liberté de la presse. »
\end{exemplebox}

\begin{attention}[\esp{Hay que} et non \esp{tener que} pour l'obligation générale]
\esp{Hay que + infinitivo} exprime une obligation \emph{impersonnelle} et générale : « il faut ». \esp{Hay que verificar las fuentes.} « Il faut vérifier les sources. » Ne dites jamais \esp{tener que} pour une obligation impersonnelle : \esp{tenemos que} signifie « nous devons » (obligation personnelle, sujet précis). Comparez : \esp{Hay que respetar la paz.} « Il faut respecter la paix » (tout le monde) ; \esp{Los alumnos tienen que respetar la paz.} « Les élèves doivent respecter la paix » (sujet précis).
\end{attention}

\courssub{Tâche 1 (écrite) : le paragraphe argumentatif}

\begin{exercice}[Tâche 1 : ¿Deben los medios cameruneses promover la cultura de paz?]
Rédigez un paragraphe argumentatif de 8 à 10 lignes répondant à la question : \esp{¿Deben los medios de comunicación cameruneses promover la cultura de paz?} Contraintes : (1) suivez l'organigramme opinión / argumentos + ejemplos / conclusión ; (2) employez au moins 5 mots du lexique étudié : \esp{cuarto poder, libertad de prensa, independencia, censura, opinión pública, mediador, reconciliación} ; (3) utilisez au moins une fois \esp{hay que + infinitivo} ; (4) terminez par \esp{En conclusión} ou \esp{En definitiva}.
\end{exercice}

\begin{exoresolu}[Corrigé guidé de la tâche 1]
Énoncé : produire le paragraphe demandé avec les contraintes ci-dessus.
\tcblower
Solution détaillée. Étape 1 : je choisis mon opinion (oui, les médias doivent promouvoir la paix). Étape 2 : je sélectionne mes mots du lexique : \esp{opinión pública, mediador, reconciliación, censura, cuarto poder}. Étape 3 : je rédige en suivant l'organigramme.\\
Modèle de réponse : \esp{Desde mi punto de vista, los medios cameruneses deben promover la cultura de paz. En primer lugar, la prensa influye en la opinión pública: por ejemplo, un reportaje sobre la reconciliación en una aldea puede cambiar las actitudes de miles de lectores. Además, el periodista puede ser un mediador entre las comunidades en conflicto, como ocurre cuando una radio da la palabra a todos los bandos. Sin embargo, para cumplir esta misión, hay que rechazar la censura y defender la independencia de la prensa. En conclusión, el cuarto poder es un instrumento de paz: sin medios libres y responsables, no hay reconciliación posible.}\\
Vérification : 5 mots du lexique soulignés, un \esp{hay que}, une conclusion introduite par \esp{En conclusión}, environ 9 lignes. Toutes les contraintes sont respectées.
\end{exoresolu}

\courssub{Tâche 2 (orale) : le débat en binômes}

\begin{experience}[Débat : libertad de prensa contra regulación]
Travaillez en binômes. L'élève A défend la \emph{liberté de presse totale} ; l'élève B plaide pour une \emph{régulation contre les fake news}. Chacun parle une minute, puis réagit aux arguments de l'autre pendant deux minutes. La classe vote ensuite pour la position la plus convaincante.
\end{experience}

\begin{methode}[Réussir le débat oral]
\begin{enumerate}
\item \textbf{Commencez} par annoncer votre position : \esp{Desde mi punto de vista, la libertad de prensa no debe tener límites...} ou \esp{En mi opinión, hay que regular los medios contra las noticias falsas...}.
\item \textbf{Argumentez} avec des connecteurs : \esp{En primer lugar... Además... Por ejemplo...}.
\item \textbf{Réagissez} à votre partenaire avec politesse : \esp{Estoy de acuerdo contigo, pero...} (je suis d'accord avec toi, mais...) ou \esp{No comparto tu opinión porque...} (je ne partage pas ton opinion parce que...).
\item \textbf{Concluez} : \esp{En definitiva, pienso que...}.
\end{enumerate}
\end{methode}

\begin{exemplebox}[Répliques possibles pour le débat]
Élève A : \esp{La censura es peor que las noticias falsas: sin libertad de prensa no hay democracia.} « La censure est pire que les fausses informations : sans liberté de presse, il n'y a pas de démocratie. »\\
Élève B : \esp{No comparto tu opinión porque las noticias falsas pueden provocar violencia en las elecciones.} « Je ne partage pas ton opinion parce que les fausses informations peuvent provoquer des violences lors des élections. »\\
Élève A : \esp{Estoy de acuerdo contigo, pero hay que educar a los lectores en lugar de censurar.} « Je suis d'accord avec toi, mais il faut éduquer les lecteurs au lieu de censurer. »
\end{exemplebox}

\courssub{Tâche 3 : titular y chapô d'un article de paix}

Le \esp{titular} (titre) est court, frappant, souvent sans article ; le \esp{chapó} (chapeau) est le petit paragraphe en gras sous le titre qui résume le qui, le quoi, le où et le quand.

\begin{exemplebox}[Modèle : un club de paix au lycée]
\textbf{Titular} : \esp{El Club de Paz del Liceo de Yaoundé siembra diálogo entre los alumnos}\\
\textbf{Chapó} : \esp{Creado por veinte estudiantes de Terminale, el Club de Paz organiza cada viernes sesiones de mediación escolar para resolver los conflictos sin violencia y promover la reconciliación en el establecimiento.} « Créé par vingt élèves de Terminale, le Club de Paix organise chaque vendredi des séances de médiation scolaire pour résoudre les conflits sans violence et promouvoir la réconciliation dans l'établissement. »
\end{exemplebox}

\begin{exercice}[Tâche 3 : à vous d'écrire]
Imaginez une initiative de paix dans votre établissement (club de paix, médiation scolaire, match de football pour la réconciliation, journée du vivre-ensemble). Rédigez le \esp{titular} (maximum 10 mots) et le \esp{chapó} (2 à 3 phrases) de l'article qui la présente. Employez au moins deux mots du lexique : \esp{mediador, reconciliación, cultura de paz, opinión pública}.
\end{exercice}

\courssub{Reformulation : jouer avec les synonymes}

Reformuler, c'est dire la même chose avec d'autres mots : compétence essentielle pour le commentaire et la traduction, qui prouve que vous possédez le lexique au lieu de le répéter.

\begin{center}
\begin{tabularx}{\linewidth}{|X|X|}
\hline
\rowcolor{popgreenL} Mot du texte & Synonymes possibles \\
\hline
\esp{la prensa} & \esp{los medios de comunicación, los periodistas, el cuarto poder} \\
\hline
\esp{la paz} & \esp{la reconciliación, la convivencia, la armonía} \\
\hline
\esp{informar} & \esp{comunicar, dar a conocer, transmitir} \\
\hline
\esp{las noticias falsas} & \esp{la desinformación, los rumores, las mentiras} \\
\hline
\esp{debe} & \esp{tiene que, hay que (+ inf.), es necesario que (+ subj.)} \\
\hline
\end{tabularx}
\end{center}

\begin{exoresolu}[Reformulation de 5 phrases du texte]
Énoncé : transformez ces phrases en utilisant des synonymes du lexique étudié, sans changer le sens.\\
1. \esp{La prensa debe informar con verdad.} 2. \esp{Los medios promueven la paz.} 3. \esp{El periodista actúa como mediador.} 4. \esp{Hay que combatir las noticias falsas.} 5. \esp{La censura limita la libertad de prensa.}
\tcblower
Solutions possibles :\\
1. \esp{El cuarto poder tiene que comunicar con honestidad.} 2. \esp{Los medios de comunicación fomentan la reconciliación y la convivencia.} 3. \esp{El periodista sirve de árbitro entre las partes en conflicto.} 4. \esp{Es necesario luchar contra la desinformación.} 5. \esp{El control del Estado reduce la independencia de los medios.}\\
Remarquez que chaque reformulation garde le sens mais change au moins deux éléments lexicaux : c'est la marque d'un bon commentaire.
\end{exoresolu}

\courssub{Grille d'auto-évaluation}

Après chaque production, évaluez-vous honnêtement avec cette grille (cochez 0 = absent, 1 = partiel, 2 = réussi).

\begin{center}
\begin{tabularx}{\linewidth}{|X|c|c|}
\hline
\rowcolor{popgoldL} Critère & Tâche 1 & Tâche 3 \\
\hline
J'ai employé au moins 5 mots du lexique du cuarto poder & /2 & /2 \\
\hline
J'ai utilisé des connecteurs (\esp{en primer lugar, además, sin embargo}) & /2 & /2 \\
\hline
J'ai exprimé clairement mon opinion personnelle & /2 & /2 \\
\hline
J'ai conclu avec \esp{En conclusión} ou \esp{En definitiva} & /2 & /2 \\
\hline
Correction grammaticale (accords, conjugaisons, accents) & /2 & /2 \\
\hline
\textbf{Total} & /10 & /10 \\
\hline
\end{tabularx}
\end{center}

\begin{aretenir}[L'essentiel de la production guidée]
Un bon texte argumentatif en espagnol suit toujours le même chemin : \cle{opinión} (\esp{En mi opinión...}), \cle{argumentos + ejemplos} (\esp{En primer lugar... Además... Por ejemplo...}), \cle{conclusión} (\esp{En conclusión...}). Pour l'obligation générale, dites \esp{hay que + infinitivo}, jamais \esp{tener que} sans sujet précis. Au débat, réagissez avec courtoisie : \esp{Estoy de acuerdo contigo, pero...} / \esp{No comparto tu opinión porque...}. Enfin, reformuler avec des synonymes du lexique (\esp{prensa = cuarto poder = medios de comunicación}) est la preuve que le vocabulaire est vraiment acquis.
\end{aretenir}

\begin{savaistu}[La paix, une valeur partagée]
Le Cameroun célèbre chaque année la Journée internationale de la paix avec des marches, des matchs de football et des émissions de radio communautaires. En Guinée équatoriale, voisin hispanophone, les radios nationales diffusent aussi des programmes en espagnol consacrés à la \esp{convivencia} (le vivre-ensemble). En produisant vos propres textes pour la paix en espagnol, vous rejoignez une grande famille de jeunes francophones et hispanophones qui utilisent les mots comme des ponts, jamais comme des armes.
\end{savaistu}

\newpage

\coursec{Activité d'intégration}

\courssub{Objectifs de l'activité}

Cette activité d'intégration clôt la leçon E24 en faisant travailler ensemble toutes les compétences acquises. À la fin de la séance, l'élève doit être capable de :

\begin{itemize}
\item mobiliser le \cle{lexique du quatrième pouvoir} (\esp{cuarto poder, libertad de prensa, censura, opinión pública, mediador, reconciliación, diálogo}) dans un texte d'opinion cohérent ;
\item appliquer la \cle{méthode du commentaire de texte} : repérer la thèse, les arguments, les procédés de persuasion, et s'en inspirer pour sa propre production ;
\item construire une \cle{argumentation écrite} en espagnol : présenter un problème, prendre position, proposer des solutions, conclure ;
\item employer les \cle{connecteurs argumentatifs} et les temps verbaux adaptés au texte d'opinion (présent de vérité générale, subjonctif après les verbes d'appel et de souhait) ;
\item s'exprimer à l'oral : lire un texte à voix haute, justifier un vote, défendre un point de vue en espagnol.
\end{itemize}

\courssub{Situation}

\begin{textebox}[La situación : tensión en el Liceo Bilingüe de Yaoundé]
Desde hace dos semanas, el clima escolar del Liceo Bilingüe de Yaoundé se ha deteriorado. Tras la derrota del equipo de fútbol de la clase de Terminale A2 frente al de Terminale C durante el torneo interclases, circularon rumores falsos en los grupos de WhatsApp: se acusó al árbitro, alumno del liceo, de haber favorecido a los rivales. Los insultos se multiplicaron en las redes sociales y ayer, durante el recreo, dos alumnos estuvieron a punto de pelearse delante del comedor. El director de estudios ha pedido calma, pero la tensión sigue presente en los pasillos.

La redacción del periódico escolar bilingüe \emph{La Voz del Liceo} decide reaccionar: en lugar de alimentar la polémica, quiere publicar un editorial que ayude a restablecer el diálogo entre los alumnos.
\end{textebox}

\textbf{Traduction de la situation :} Depuis deux semaines, le climat scolaire du Lycée Bilingue de Yaoundé s'est détérioré. Après la défaite de l'équipe de football de la classe de Terminale A2 face à celle de Terminale C lors du tournoi interclasses, de fausses rumeurs ont circulé dans les groupes WhatsApp : on a accusé l'arbitre, élève du lycée, d'avoir favorisé les adversaires. Les insultes se sont multipliées sur les réseaux sociaux et hier, pendant la récréation, deux élèves ont failli se battre devant le réfectoire. Le directeur des études (censeur) a demandé le calme, mais la tension reste présente dans les couloirs. La rédaction du journal scolaire bilingue \emph{La Voz del Liceo} décide de réagir : au lieu d'alimenter la polémique, elle veut publier un éditorial qui aide à rétablir le dialogue entre les élèves.

Votre classe incarne cette rédaction. Par groupes de quatre, vous êtes les journalistes chargés d'écrire l'éditorial du numéro de la semaine, intitulé \esp{« Los medios, mediadores de la paz en nuestro liceo »}. Votre texte sera lu à voix haute devant toute la rédaction (la classe), affiché au tableau, et la classe votera pour l'éditorial le plus convaincant en justifiant son choix en espagnol.

\courssub{Consignes}

\begin{enumerate}
\item \textbf{Relecture ciblée (10 min).} Relisez le texte étudié \esp{« El cuarto poder al servicio de la paz »}. Relevez dans une grille : la thèse du journaliste, deux arguments, deux exemples, et trois connecteurs argumentatifs utilisés. Cette grille servira de modèle de structure pour votre propre texte.

\item \textbf{Remue-méninges en groupe (10 min).} En groupe de quatre, répondez à voix haute, en espagnol, aux questions suivantes : \esp{¿Qué papel jugaron los rumores en esta tensión? ¿Qué puede hacer el periódico escolar por la paz? ¿A quién hay que dirigirse y qué hay que proponer?} Notez vos idées en deux colonnes : \esp{problema} / \esp{soluciones}.

\item \textbf{Plan de l'éditorial (5 min).} Organisez votre texte en quatre parties : (a) présentation du problème sans accuser personne ; (b) rappel du rôle des médias et du journal scolaire comme \esp{mediador} ; (c) propositions concrètes (débat interclasses, match amical de la réconciliation, tribune libre dans le journal) ; (d) appel au \esp{diálogo} et à la \esp{reconciliación}.

\item \textbf{Rédaction (25 min).} Rédigez l'éditorial de 15 à 20 lignes. Contraintes obligatoires : au moins \textbf{six mots du lexique} étudié (soulignez-les), au moins \textbf{trois connecteurs argumentatifs} différents, au moins \textbf{un subjonctif} après un verbe d'appel ou de souhait, et une conclusion qui invite à la \esp{cultura de paz}.

\item \textbf{Relecture et correction mutuelle (10 min).} Échangez votre texte avec le groupe voisin. Vérifiez : les accents, la concordance des temps, la présence des contraintes, la clarté de la thèse. Signalez deux points forts et une erreur à corriger.

\item \textbf{Lecture publique et vote (15 min).} Chaque groupe lit son éditorial à voix haute. La classe vote pour le texte le plus convaincant. Chaque électeur justifie son vote en une phrase en espagnol, par exemple : \esp{« Voto por el grupo tres porque sus propuestas son concretas y su conclusión es muy emotiva. »}
\end{enumerate}

\courssub{Ressources à mobiliser}

\begin{aretenir}[Lexique du quatrième pouvoir et de la paix]
\begin{itemize}
\item \esp{el cuarto poder} : la presse considérée comme un pouvoir qui surveille et influence la société.
\item \esp{la libertad de prensa} : droit d'informer sans \esp{censura}.
\item \esp{la opinión pública} : ce que pense la majorité des citoyens.
\item \esp{el/la mediador(a)} : personne ou institution qui aide deux parties en conflit à dialoguer.
\item \esp{el diálogo}, \esp{la reconciliación}, \esp{la cultura de paz}, \esp{el rumor}, \esp{la calumnia} (calomnie), \esp{informar con objetividad}.
\end{itemize}
\end{aretenir}

\begin{propriete}[Connecteurs argumentatifs utiles]
Pour structurer l'éditorial : \esp{en primer lugar} (d'abord), \esp{además} (de plus), \esp{sin embargo} (cependant), \esp{por lo tanto / por eso} (c'est pourquoi), \esp{en conclusión} (en conclusion), \esp{es decir} (c'est-à-dire), \esp{no solo... sino también} (non seulement... mais aussi).
\end{propriete}

\begin{propriete}[Le subjonctif après les verbes d'appel et de souhait]
Après les verbes qui expriment un souhait, une demande ou un appel (\esp{querer que, pedir que, es necesario que, es imprescindible que, ojalá}), on emploie le \cle{subjonctif} : \esp{Es necesario que todos dialoguemos.} « Il faut que nous dialoguions tous. » \esp{Pedimos que nadie difunda rumores.} « Nous demandons que personne ne répande de rumeurs. »
\end{propriete}

\begin{attention}[Pièges fréquents des francophones]
\begin{itemize}
\item Ne dites pas \esp{*los medios de comunicación son un poder cuarto} : l'expression figée est \esp{el cuarto poder}.
\item \textbf{Faux ami :} \esp{actualmente} signifie « actuellement / de nos jours », et \textbf{non} « en fait / en réalité » (qui se dit \esp{en realidad} ou \esp{de hecho}).
\item \esp{Una noticia falsa} (ordre standard) = une fausse information (fake news). \esp{Una falsa noticia} (adjectif antéposé) = insistance expressive sur le caractère mensonger. Pour « fake news », on privilégie \esp{noticia falsa} ou \esp{bulo}.
\item N'oubliez pas les accents : \esp{opinión, reconciliación, diálogo, crítica, público}.
\item « Défendre la paix » se traduit \esp{defender la paz}, jamais \esp{*defender la paz de} + infinitif (calque du français).
\end{itemize}
\end{attention}

\courssub{Proposition de production et corrigé commenté}

\begin{exoresolu}[Editorial modelo : « Los medios, mediadores de la paz en nuestro liceo »]
Rédigez un éditorial de 15 à 20 lignes respectant toutes les contraintes de la consigne.
\tcblower
\textbf{Production modèle :}

\esp{Desde hace dos semanas, nuestro liceo vive un momento difícil: después del partido de fútbol interclases, los rumores y los insultos en las redes sociales han envenenado el ambiente escolar. En primer lugar, recordemos que un rumor no es una noticia: informar con objetividad es la primera misión de cualquier medio, y también la nuestra. Además, como cuarto poder a pequeña escala, el periódico escolar no debe alimentar la polémica, sino actuar como mediador entre los alumnos enfrentados. Sin embargo, la libertad de prensa no significa decirlo todo sin reflexionar: significa buscar la verdad y servir a la comunidad. Por eso, proponemos tres acciones concretas: un debate interclases sobre el respeto en el deporte, un partido amistoso de la reconciliación y una tribuna libre en nuestro próximo número. Es necesario que todos participemos en este diálogo y que nadie difunda calumnias. En conclusión, la paz no cae del cielo: se construye cada día, con palabras justas. Hagamos de nuestro liceo un ejemplo de cultura de paz.}

\begin{textebox}[Traduction de l'éditorial modèle]
« Depuis deux semaines, notre lycée traverse un moment difficile : après le match de football interclasses, les rumeurs et les insultes sur les réseaux sociaux ont empoisonné le climat scolaire. En premier lieu, rappelons-nous qu'une rumeur n'est pas une information : informer avec objectivité est la première mission de tout média, et aussi la nôtre. De plus, en tant que quatrième pouvoir à petite échelle, le journal scolaire ne doit pas alimenter la polémique, mais agir en médiateur entre les élèves en conflit. Cependant, la liberté de presse ne signifie pas tout dire sans réfléchir : cela signifie chercher la vérité et servir la communauté. C'est pourquoi nous proposons trois actions concrètes : un débat interclasses sur le respect dans le sport, un match amical de la réconciliation et une tribune libre dans notre prochain numéro. Il est nécessaire que nous participions tous à ce dialogue et que personne ne répande de calomnies. En conclusion, la paix ne tombe pas du ciel : elle se construit chaque jour, avec des mots justes. Faisons de notre lycée un exemple de culture de paix. »
\end{textebox}

\textbf{Commentaire des choix de langue :}
\begin{itemize}
\item \textbf{Lexique (8 mots) :} \esp{rumores, informar con objetividad, cuarto poder, mediador, libertad de prensa, reconciliación, diálogo, cultura de paz} — la contrainte des six mots est dépassée.
\item \textbf{Connecteurs (5) :} \esp{en primer lugar, además, sin embargo, por eso, en conclusión} — la progression problème / rôle / propositions / appel est clairement balisée.
\item \textbf{Subjonctifs :} \esp{es necesario que todos participemos... y que nadie difunda} — deux subjonctifs corrects après une tournure impersonnelle d'obligation (\esp{participemos} = 1ère personne du pluriel, \esp{difunda} = 3ème personne du singulier).
\item \textbf{Procédés de persuasion :} antithèse (\esp{un rumor no es una noticia}), énumération de trois propositions, phrase brève finale à effet (\esp{se construye cada día}), impératif collectif \esp{hagamos} qui associe le lecteur.
\item \textbf{Neutralité :} l'éditorial ne nomme ni l'arbitre ni les élèves en conflit, ce qui illustre la déontologie du \esp{mediador}.
\end{itemize}
\end{exoresolu}

\courssub{Bilan de l'activité}

Cette activité a permis de vérifier l'acquisition des trois savoir-faire de la leçon : \textbf{lire et commenter} un texte sur la presse (grille d'analyse et reprise de sa structure), \textbf{employer le lexique} du quatrième pouvoir dans un contexte authentique, et \textbf{défendre la promotion de la paix} par une argumentation écrite et orale. Les élèves ont également pratiqué l'écoute active (lecture des éditoriaux) et l'expression d'un jugement justifié en espagnol (vote argumenté).

Pour prolonger le travail à la maison, chaque élève pourra rédiger individuellement une \esp{carta al director} (lettre au rédacteur en chef) de dix lignes réagissant à l'éditorial gagnant, en exprimant son accord ou son désaccord avec au moins deux connecteurs et un subjonctif. Les meilleurs textes seront publiés dans le prochain numéro de \emph{La Voz del Liceo}, transformant ainsi l'exercice scolaire en véritable acte de communication.

\newpage

\coursec{Méthodes et exercices résolus}

\coursec{Lectura y comentario : el cuarto poder y la promoción de la cultura de paz}

\courssub{Texte support : El cuarto poder al servicio de la paz}

\begin{textebox}[Article d'opinion original — Contexte : sortie de crise post-électorale au Cameroun]
\esp{En una sociedad democrática, el cuarto poder no es un lujo, es una necesidad vital. Los medios de comunicación —prensa escrita, radio, televisión y portales digitales— tienen la misión de informar con rigor, verificar los hechos y dar voz a quienes no la tienen.}\par
\esp{Sin embargo, cuando la censura, la autocensura o la manipulación política se imponen, la prensa deja de ser mediadora y se convierte en instrumento de división. Los titulares sensacionalistas y las \emph{fake news} siembran la desconfianza; en cambio, el periodismo de investigación y la ética profesional construyen puentes.}\par
\esp{En nuestro país, como en toda América Latina y en Guinea Ecuatorial, la reconciliación nacional pasa necesariamente por una prensa libre, independiente y responsable. Es imprescindible que los periodistas ejerzan su oficio sin miedo, que los ciudadanos desarrollen un pensamiento crítico y que el Estado garantice el derecho a la información.}\par
\esp{Solo así la opinión pública se formará en la verdad y la cultura de paz dejará de ser un eslogan para convertirse en realidad cotidiana.}
\end{textebox}

\begin{exemplebox}[Traduction et glossaire]
\textbf{Traduction :} « Dans une société démocratique, le quatrième pouvoir n'est pas un luxe, c'est une nécessité vitale. Les médias de communication — presse écrite, radio, télévision et portails numériques — ont la mission d'informer avec rigueur, de vérifier les faits et de donner la parole à ceux qui ne l'ont pas. Cependant, lorsque la censure, l'autocensure ou la manipulation politique s'imposent, la presse cesse d'être médiatrice et devient un instrument de division. Les titres sensationnalistes et les \emph{fake news} sèment la méfiance ; en revanche, le journalisme d'investigation et l'éthique professionnelle construisent des ponts. Dans notre pays, comme dans toute l'Amérique latine et en Guinée équatoriale, la réconciliation nationale passe nécessairement par une presse libre, indépendante et responsable. Il est indispensable que les journalistes exercent leur métier sans peur, que les citoyens développent une pensée critique et que l'État garantisse le droit à l'information. Ce n'est qu'ainsi que l'opinion publique se formera dans la vérité et que la culture de paix cessera d'être un slogan pour devenir une réalité quotidienne. »\par\smallskip
\textbf{Glossaire :} \esp{el cuarto poder} (le quatrième pouvoir) ; \esp{la prensa escrita} (la presse écrite) ; \esp{el rigor} (la rigueur) ; \esp{la autocensura} (l'autocensure) ; \esp{la manipulación} (la manipulation) ; \esp{el titular sensacionalista} (le titre sensationnaliste/putaclic) ; \esp{la desconfianza} (la méfiance) ; \esp{el periodismo de investigación} (le journalisme d'investigation) ; \esp{la ética profesional} (la déontologie professionnelle) ; \esp{la reconciliación nacional} (la réconciliation nationale) ; \esp{el pensamiento crítico} (l'esprit critique) ; \esp{el derecho a la información} (le droit à l'information) ; \esp{el eslogan} (le slogan) ; \esp{cotidiana} (quotidienne).
\end{exemplebox}

\courssub{Compréhension du texte}

\begin{methode}[Comprendre un texte argumentatif sur le \esp{cuarto poder}]
\begin{enumerate}
\item \textbf{Lecture globale.} Lis le texte deux fois : une première fois sans dictionnaire pour saisir le thème général (le rôle des médias dans la paix), une seconde fois en soulignant le lexique-clé : \esp{libertad de prensa}, \esp{censura}, \esp{autocensura}, \esp{opinión pública}, \esp{mediador}, \esp{reconciliación}, \esp{ética profesional}, \esp{pensamiento crítico}.
\item \textbf{Repérer la thèse.} Un texte argumentatif défend une idée. Cherche la phrase qui la résume (souvent au début ou à la fin) et les connecteurs logiques : \esp{sin embargo} (cependant), \esp{por lo tanto} (donc), \esp{además} (de plus), \esp{en conclusión}, \esp{es imprescindible que}.
\item \textbf{Distinguer faits et opinions.} Un fait : \esp{Los medios tienen la misión de informar.} Une opinion : \esp{La prensa debe ser independiente.} Les verbes comme \esp{deber}, \esp{pensar}, \esp{creer}, \esp{considerar} ou les expressions impersonnelles \esp{es necesario que}, \esp{es imprescindible que} signalent l'opinion ou l'évaluation.
\item \textbf{Répondre aux questions.} Reformule avec tes propres mots, cite une phrase du texte comme preuve (entre guillemets \esp{« ... »}), et réponds toujours par des phrases complètes en espagnol, jamais par un simple \esp{sí} ou \esp{no}.
\item \textbf{Vérifier.} Relis : accord sujet-verbe, accents (\esp{opinión}, \esp{información}, \esp{diálogo}, \esp{reconciliación}), tilde du \esp{ñ}, ponctuation \esp{¿...?} et \esp{¡...!}.
\end{enumerate}
\end{methode}

\begin{exoresolu}[Questions de compréhension — type Bac]
\textbf{Texte support (extrait du texte ci-dessus) :} \esp{En una sociedad democrática, el cuarto poder no es un lujo, es una necesidad vital. Los medios de comunicación tienen la misión de informar con rigor, verificar los hechos y dar voz a quienes no la tienen. Sin embargo, cuando la censura se impone, la prensa deja de ser mediadora y se convierte en instrumento de división. Por lo tanto, la libertad de prensa es condición de la reconciliación nacional.}\par\smallskip
\textbf{Énoncé :} 1) ¿Cuál es la tesis del autor? 2) Según el texto, ¿cuáles son las tres misiones de los medios? 3) ¿Qué ocurre cuando la censura se impone? 4) ¿Estás de acuerdo con la tesis? Justifica tu respuesta con un ejemplo de nuestro contexto camerounais.
\tcblower
\textbf{Raisonnement :} La question 1 porte sur l'idée principale (synthèse première et dernière phrase) ; les questions 2 et 3 demandent de repérer des informations explicites (énumération, conséquence) ; la question 4 est une question d'opinion personnelle : il faut donner un avis argumenté et l'ancrer dans le contexte local (Cameroun).\par\smallskip
\textbf{Réponses rédigées :}\par
1) \esp{La tesis del autor es que la prensa no es un lujo sino una necesidad vital para la democracia: la libertad de prensa es condición indispensable para la reconciliación nacional.}\par
2) \esp{Según el texto, las tres misiones son: informar con rigor, verificar los hechos y dar voz a quienes no la tienen.}\par
3) \esp{Cuando la censura se impone, la prensa deja de ser mediadora y se convierte en instrumento de división.}\par
4) \esp{Estoy de acuerdo con el autor. En Camerún, durante los periodos electorales, los rumores en las redes sociales a veces provocan tensiones étnicas; una prensa libre y profesional, que verifica sus fuentes, permite apaciguar los ánimos y favorece el diálogo entre las comunidades.}\par\smallskip
\textbf{Conclusion méthodologique :} Chaque réponse est une phrase complète, introduite par un connecteur (\esp{según el texto}, \esp{en efecto}, \esp{por ejemplo}), et s'appuie sur le texte ou sur un argument personnel clair et contextualisé.
\end{exoresolu}

\courssub{Lexique thématique : le quatrième pouvoir et la paix}

\begin{propriete}[Champs lexicaux et familles de mots]
\textbf{1. Champ de la presse et des médias}\par
\begin{center}
\begin{tabularx}{\linewidth}{|X|X|X|}\hline
\rowcolor{popblueL} \textbf{Nom / Verbe} & \textbf{Famille / Dérivés} & \textbf{Collocations utiles} \\ \hline
\esp{el periódico / el diario} & \esp{periodista, periodismo, periodístico} & \esp{leer el periódico, la portada del diario} \\ \hline
\esp{el medio (de comunicación)} & \esp{mediático, mediatizar} & \esp{los medios de comunicación, un medio digital} \\ \hline
\esp{la redacción} & \esp{redactar, redactor, redactado} & \esp{la reunión de redacción, redactar un artículo} \\ \hline
\esp{la fuente} & \esp{fidedigno, verificar} & \esp{verificar las fuentes, fuente fidedigna} \\ \hline
\esp{publicar} & \esp{publicación, público, públicamente} & \esp{publicar una noticia, de dominio público} \\ \hline
\esp{informar} & \esp{información, informativo, informado} & \esp{informar con rigor, boletín informativo} \\ \hline
\esp{la censura} & \esp{censurar, censurado, autocensura} & \esp{ejercer la censura, miedo a la censura} \\ \hline
\end{tabularx}
\end{center}

\textbf{2. Champ du pouvoir et de la citoyenneté}\par
\begin{center}
\begin{tabularx}{\linewidth}{|X|X|X|}\hline
\rowcolor{poporangeL} \textbf{Nom / Verbe} & \textbf{Famille / Dérivés} & \textbf{Collocations utiles} \\ \hline
\esp{la libertad de prensa} & \esp{libre, liberar, libremente} & \esp{ejercer la libertad de prensa, prensa libre} \\ \hline
\esp{la independencia} & \esp{independiente, independizarse} & \esp{prensa independiente, poder judicial independiente} \\ \hline
\esp{la opinión pública} & \esp{opinar, opinionado} & \esp{formar la opinión pública, manipular la opinión} \\ \hline
\esp{el ciudadano} & \esp{ciudadanía, ciudadanizar} & \esp{derechos del ciudadano, participación ciudadana} \\ \hline
\esp{el Estado} & \esp{estatal, estatalizar} & \esp{garantizar el Estado de derecho} \\ \hline
\end{tabularx}
\end{center}

\textbf{3. Champ de la paix et de la médiation}\par
\begin{center}
\begin{tabularx}{\linewidth}{|X|X|X|}\hline
\rowcolor{popgreenL} \textbf{Nom / Verbe} & \textbf{Famille / Dérivés} & \textbf{Collocations utiles} \\ \hline
\esp{el mediador} & \esp{mediar, mediación, mediático} & \esp{actuar como mediador, proceso de mediación} \\ \hline
\esp{la reconciliación} & \esp{reconciliar, reconciliador} & \esp{favorecer la reconciliación, reconciliación nacional} \\ \hline
\esp{el diálogo} & \esp{dialogar, dialogante} & \esp{abrir un diálogo, mesa de diálogo} \\ \hline
\esp{la tolerancia} & \esp{tolerar, tolerante} & \esp{promover la tolerancia, tolerancia cero} \\ \hline
\esp{el conflicto} & \esp{conflictivo, conflictuar} & \esp{resolver un conflicto, zona conflictiva} \\ \hline
\esp{la cultura de paz} & \esp{pacífico, pacificar, pacifista} & \esp{promover la cultura de paz, educación para la paz} \\ \hline
\end{tabularx}
\end{center}
\end{propriete}

\begin{methode}[Réutiliser le vocabulaire : stratégies pour l'écrit et l'oral]
\begin{enumerate}
\item \textbf{Classer par champs sémantiques} (voir tableaux ci-dessus) pour activer la mémoire à long terme.
\item \textbf{Construire des familles de mots} pour varier les formes : \esp{publicar / publicación / público} ; \esp{informar / información / informativo} ; \esp{censurar / censura / censurado} ; \esp{reconciliar / reconciliación / reconciliador} ; \esp{mediar / mediador / mediación}.
\item \textbf{Mémoriser les collocations (associations figées)} : \esp{ejercer la libertad de prensa}, \esp{formar la opinión pública}, \esp{verificar una fuente}, \esp{promover la cultura de paz}, \esp{desempeñar el papel de mediador}, \esp{garantizar el derecho a la información}.
\item \textbf{Éviter les calques du français} :
\begin{itemize}
\item « Un journal » $\to$ \esp{un periódico} ou \esp{un diario} (jamais \esp{un journal}).
\item « L'actualité » $\to$ \esp{la actualidad}.
\item « Un média » $\to$ \esp{un medio (de comunicación)}.
\item « Un reportage » $\to$ \esp{un reportaje}.
\item « La rédaction (équipe) » $\to$ \esp{la redacción}.
\end{itemize}
\item \textbf{Surveiller l'orthographe des verbes en -zar / -cer} : \esp{rechazar} $\to$ \esp{rechacé} (z $\to$ c devant e) ; \esp{ejercer} $\to$ \esp{ejerzo} (c $\to$ z devant o/a) ; \esp{vencer} $\to$ \esp{venzo}.
\end{enumerate}
\end{methode}

\begin{exoresolu}[Lexique en contexte — type Bac]
\textbf{Énoncé :} Complète el texto con las palabras siguientes : \esp{censura — mediadores — opinión pública — libertad de prensa — reconciliación}.\par\smallskip
\esp{Después del conflicto postelectoral, los periodistas actuaron como (1) \ldots\ entre las dos comunidades. Gracias a la (2) \ldots, pudieron informar sin miedo a la (3) \ldots. Sus reportajes formaron una (4) \ldots\ favorable al diálogo y abrieron el camino hacia la (5) \ldots\ nacional.}
\tcblower
\textbf{Raisonnement :} On identifie la fonction grammaticale et le sens de chaque trou : (1) après \esp{como}, nom de fonction/pluriel accordé avec \esp{los periodistas} ; (2) après \esp{gracias a la}, nom positif (condition) ; (3) après \esp{miedo a la}, nom négatif (obstacle) ; (4) après \esp{una}, nom féminin singulier (collocation \esp{formar la opinión pública}) ; (5) après \esp{la}, but final (collocation \esp{reconciliación nacional}).\par\smallskip
\textbf{Solution :}\par
(1) \esp{mediadores} — accord pluriel ; « agir comme médiateurs ».\par
(2) \esp{libertad de prensa} — condition positive.\par
(3) \esp{censura} — seul obstacle négatif de la liste.\par
(4) \esp{opinión pública} — collocation figée \esp{formar la opinión pública}.\par
(5) \esp{reconciliación} — collocation \esp{la reconciliación nacional}.\par\smallskip
\textbf{Texte complété :} \esp{Después del conflicto postelectoral, los periodistas actuaron como mediadores entre las dos comunidades. Gracias a la libertad de prensa, pudieron informar sin miedo a la censura. Sus reportajes formaron una opinión pública favorable al diálogo y abrieron el camino hacia la reconciliación nacional.}
\end{exoresolu}

\begin{savaistu}[Le Cameroun, la Guinée équatoriale et la liberté de presse]
La Guinée équatoriale (seul pays africain hispanophone) et le Cameroun (pays bilingue français-anglais avec une forte communauté hispanophone) partagent des défis communs : concentration des médias publics, pressions économiques sur la presse privée, lois sur la cybercriminalité parfois utilisées pour museler les journalistes. Cependant, la société civile, les associations de journalistes (comme le SYNJIC au Cameroun) et la coopération internationale (UNESCO, OIF) œuvrent pour la \esp{ley de acceso a la información} et la protection des \esp{defensores de derechos humanos}. Connaître ce contexte enrichit votre argumentation au Bac.
\end{savaistu}

\begin{popfigure}
\begin{center}\tcbox[colback=popink,colframe=popink,coltext=white,arc=9pt,boxsep=4pt,left=6pt,right=6pt]{\bfseries\small Cuarto Poder}\end{center}
\vspace{-2pt}
\begin{tcbraster}[raster columns=2,raster equal height=rows,raster column skip=6pt,raster row skip=6pt,enhanced,blankest]
\begin{tcolorbox}[enhanced,colback=white,colframe=poporange!60,boxrule=0.9pt,arc=3pt,title={Prensa},fonttitle=\bfseries\scriptsize,coltitle=white,colbacktitle=poporange!60,left=4pt,right=4pt,top=2pt,bottom=2pt,boxsep=1pt,toptitle=1.5pt,bottomtitle=1.5pt,halign=left]
\begin{itemize}[nosep,leftmargin=*,itemsep=1pt,font=\scriptsize]
\item Periódico
\item Radio
\item Televisión
\item Digital
\end{itemize}
\end{tcolorbox}
\begin{tcolorbox}[enhanced,colback=white,colframe=popgreen!60,boxrule=0.9pt,arc=3pt,title={Libertad},fonttitle=\bfseries\scriptsize,coltitle=white,colbacktitle=popgreen!60,left=4pt,right=4pt,top=2pt,bottom=2pt,boxsep=1pt,toptitle=1.5pt,bottomtitle=1.5pt,halign=left]
\begin{itemize}[nosep,leftmargin=*,itemsep=1pt,font=\scriptsize]
\item Expresión
\item Información
\item Independencia
\end{itemize}
\end{tcolorbox}
\begin{tcolorbox}[enhanced,colback=white,colframe=poppurple!60,boxrule=0.9pt,arc=3pt,title={Paz},fonttitle=\bfseries\scriptsize,coltitle=white,colbacktitle=poppurple!60,left=4pt,right=4pt,top=2pt,bottom=2pt,boxsep=1pt,toptitle=1.5pt,bottomtitle=1.5pt,halign=left]
\begin{itemize}[nosep,leftmargin=*,itemsep=1pt,font=\scriptsize]
\item Mediación
\item Reconciliación
\item Diálogo
\end{itemize}
\end{tcolorbox}
\begin{tcolorbox}[enhanced,colback=white,colframe=poppink!60,boxrule=0.9pt,arc=3pt,title={Responsabilidad},fonttitle=\bfseries\scriptsize,coltitle=white,colbacktitle=poppink!60,left=4pt,right=4pt,top=2pt,bottom=2pt,boxsep=1pt,toptitle=1.5pt,bottomtitle=1.5pt,halign=left]
\begin{itemize}[nosep,leftmargin=*,itemsep=1pt,font=\scriptsize]
\item Ética
\item Verificación
\item Crítica
\end{itemize}
\end{tcolorbox}
\end{tcbraster}

\legende{Carte mentale : le quatrième pouvoir au cœur de la démocratie et de la paix}
\end{popfigure}

\courssub{Méthode du commentaire de texte}

\begin{methode}[La fiche du commentaire de texte en espagnol (Bac LV2)]
\begin{enumerate}
\item \textbf{Introduction} (3-4 lignes) :
\begin{itemize}
\item Présenter le document : \esp{El texto que vamos a comentar es un artículo de opinión / un fragmento periodístico que trata de\ldots}
\item Auteur / Source / Date (si connus) : \esp{Publicado en\ldots, el autor denuncia / defiende\ldots}
\item Problématique : \esp{El autor plantea la siguiente cuestión: ¿\ldots?} ou \esp{El tema central es\ldots}
\item Annonce du plan : \esp{Para analizar este texto, veremos primero\ldots, después\ldots\ y finalmente\ldots}
\end{itemize}
\item \textbf{Développement} (2 ou 3 parties équilibrées) :
\begin{itemize}
\item \textbf{A) Analyse des idées (le \emph{quoi}) :} Thèse, arguments, exemples, données. Citer le texte : \esp{El autor afirma que «\ldots»}, \esp{Argumenta diciendo que\ldots}
\item \textbf{B) Analyse des procédés (le \emph{comment}) :} Connecteurs (\esp{sin embargo, por tanto, en efecto}), champ lexical, figures de style (métaphore, antithèse, énumération), ton (indigné, espéré, objectif), temps verbaux (présent de vérité, subjonctif d'exhortation).
\item \textbf{C) Opinion personnelle argumentée (le \emph{moi}) :} \esp{En mi opinión\ldots}, \esp{Desde mi punto de vista\ldots}, \esp{Comparto / no comparto la opinión del autor porque\ldots} Apporter un exemple personnel ou camerounais.
\end{itemize}
\item \textbf{Conclusion} (2-3 lignes) :
\begin{itemize}
\item Bilan : \esp{En conclusión, este texto nos muestra que\ldots}
\item Ouverture : \esp{Esta reflexión nos lleva a preguntarnos\ldots} / \esp{Queda por ver si\ldots}
\end{itemize}
\item \textbf{Réflexes de langue obligatoires} :
\begin{itemize}
\item Présent de vérité générale pour les idées du texte.
\item Subjonctif après \esp{es importante que}, \esp{es necesario que}, \esp{no creo que}, \esp{para que}, \esp{quiero que}.
\item Connecteurs variés (\esp{en primer lugar, además, por consiguiente, en cambio, por ejemplo}).
\item Aucune phrase copiée sans guillemets \esp{« ... »}.
\end{itemize}
\end{enumerate}
\end{methode}

\begin{exoresolu}[Commentaire guidé — type Bac]
\textbf{Énoncé :} Comenta el siguiente texto en 15-20 líneas (introducción, análisis de las ideas, opinión personal, conclusión).\par\smallskip
\esp{En una sociedad dividida, el periodista no es un espectador: es un constructor de paz. Cada titular falso siembra odio; cada noticia verificada siembra confianza. Es importante que los medios respeten la verdad, porque sin verdad no hay reconciliación posible.}
\tcblower
\textbf{Raisonnement :} Texte argumentatif bref, riche en procédés (métaphore du semeur, antithèse \esp{odio/confianza}, subjonctif après \esp{es importante que}). Le commentaire suit le plan en quatre étapes.\par\smallskip
\textbf{Commentaire rédigé (modèle) :}\par
\esp{El texto que vamos a comentar es un fragmento argumentativo sobre el papel de la prensa en la cultura de paz. El autor defiende la idea de que el periodista es un constructor de paz y no un simple espectador.}\par
\esp{En primer lugar, el autor afirma su tesis con una metáfora agrícola: el periodista «siembra» odio o confianza según la calidad de su información. Esta imagen del sembrador muestra que cada noticia produce frutos en la sociedad. En segundo lugar, el texto utiliza una antítesis muy clara: «titular falso / noticia verificada», «odio / confianza». Esta oposición resume la responsabilidad moral de los medios. Finalmente, la última frase relaciona la verdad con la reconciliación mediante una estructura impersonnelle + subjonctif : \esp{es importante que los medios respeten\ldots} ; sans vérité, la paix est impossible.}\par
\esp{En mi opinión, el autor tiene razón: en nuestros países, los rumores difundidos por las redes sociales provocan a veces conflictos violentos. Por eso, es necesario que los periodistas verifiquen sus fuentes y que los ciudadanos lean con espíritu crítico.}\par
\esp{En conclusión, este texto demuestra que la libertad de prensa va unida a la responsabilidad. Nos lleva a preguntarnos: ¿cómo formar a los jóvenes para que distingan la información verdadera de la falsa?}\par\smallskip
\textbf{Justifications grammaticales :} \esp{es necesario que los periodistas verifiquen} : subjonctif présent (3\up{e} pl.) obligatoire après expression impersonnelle de nécessité ; \esp{para que distingan} : subjonctif après \esp{para que} (but) ; \esp{tiene razón} : expression figée (pas d'article) ; \esp{lean} : subjonctif présent (3\up{e} pl. de \esp{leer}) après \esp{que} subordonné à \esp{es necesario que} (implicite).
\end{exoresolu}

\courssub{Traduction : thème français $\to$ espagnol}

\begin{methode}[Réussir le thème : étapes et pièges]
\begin{enumerate}
\item \textbf{Analyser la phrase française} : repère sujet, verbe, temps, compléments, avant de traduire mot à mot.
\item \textbf{Choisir le bon temps du passé} :
\begin{itemize}
\item Action datée, ponctuelle, achevée $\to$ \esp{Pretérito indefinido} (\esp{Ayer publicaron}).
\item Description, habitude, action en cours dans le passé $\to$ \esp{Pretérito imperfecto} (\esp{Cuando era niño, leía}).
\item Antériorité dans le passé $\to$ \esp{Pluscuamperfecto} (\esp{Había publicado}).
\end{itemize}
\item \textbf{Constructions spéciales} :
\begin{itemize}
\item « Il faut que + subjonctif » $\to$ \esp{Es necesario / imprescindible / importante que + subjuntivo}.
\item « On + verbe » $\to$ \esp{Se + verbe} (passif réfléchi) ou 3\up{e} personne pluriel : \esp{Se habla español} / \esp{Hablan español}.
\item « Venir de + infinitif » $\to$ \esp{Acabar de + infinitivo} : \esp{Acaban de llegar}.
\item « Faire + infinitif » (causatif) $\to$ \esp{Hacer + infinitivo} : \esp{Hizo publicar el artículo}.
\end{itemize}
\item \textbf{\esp{Ser} ou \esp{Estar} ?} \esp{Ser} = identité, définition, origine, caractéristique essentielle, heure, possession (\esp{La prensa es libre}). \esp{Estar} = localisation, état temporaire, résultat d'action (\esp{El periódico está en la mesa}, \esp{Está censurado}).
\item \textbf{\esp{Por} ou \esp{Para} ?} \esp{Por} = cause, moyen, durée, échange, mouvement vague. \esp{Para} = but, destination, destinataire, échéance, opinion (\esp{Para mí, es así}).
\item \textbf{Relire systématiquement} : accents (\esp{opinión, información, diálogo, reconciliación, promoción, nación, televisión, también}), \esp{ñ}, signes \esp{¿ ¡}.
\end{enumerate}
\end{methode}

\begin{exoresolu}[Thème — type Bac]
\textbf{Énoncé :} Traduis en espagnol :
1) La presse joue un rôle important dans la promotion de la paix.
2) Hier, les journalistes ont publié un article sur la réconciliation.
3) Il faut que les médias respectent la vérité.
4) Quand j'étais enfant, mon père lisait le journal chaque matin.
5) Les fausses nouvelles sont dangereuses pour la démocratie.
\tcblower
\textbf{Solutions détaillées :}\par
1) \esp{La prensa desempeña un papel importante en la promoción de la paz.} — « Jouer un rôle » = \esp{desempeñar / jugar un papel} ; \esp{promoción} (mot aigu en -n) porte accent.\par
2) \esp{Ayer, los periodistas publicaron un artículo sobre la reconciliación.} — « Hier » (\esp{ayer}) impose le \esp{pretérito indefinido} : \esp{publicaron} (3\up{e} pl., régulier). Pas de \esp{pretérito perfecto} (\esp{han publicado}) avec un marqueur temporel passé.\par
3) \esp{Es necesario que los medios respeten la vérité.} — « Il faut que » $\to$ \esp{Es necesario que} + \textbf{subjonctif présent} : \esp{respeten} (3\up{e} pl., \emph{non} \esp{respetan}). « Médias » = \esp{los medios (de comunicación)}.\par
4) \esp{Cuando era niño, mi padre leía el periódico cada mañana.} — Habitude et description au passé $\to$ \esp{pretérito imperfecto} : \esp{era} (irrégulier \esp{ser}), \esp{leía} (accent sur le \emph{í} pour 1\up{re}/3\up{e} pers. sing. de l'imparfait en -\emph{er}/\emph{-ir}). « Le journal » = \esp{el periódico}.\par
5) \esp{Las noticias falsas son peligrosas para la democracia.} — « Fausses nouvelles » = \esp{noticias falsas} (ou \esp{bulos} / \esp{noticias falsificadas}) ; \esp{peligrosas} (accord fém. pl.) ; \esp{para} (but/destination du danger).\par\smallskip
\textbf{Conclusion :} Au thème, le choix du temps (indefinido vs imperfecto) et l'emploi du subjonctif sont les deux points les plus sanctionnés : toujours justifier mentalement le temps avant d'écrire.
\end{exoresolu}

\courssub{Production guidée : défendre la culture de la paix}

\begin{methode}[Rédiger un paragraphe argumentatif pour la paix (écrit / oral)]
\begin{enumerate}
\item \textbf{Structurer (squelette IAE : Idée - Argument - Exemple) :}
\esp{La prensa contribuye a la paz. (Idée) En efecto, informa con verdad. (Argument) Por ejemplo, durante las elecciones, los medios serios evitan los rumores. (Exemple) Así, protegen la convivencia. (Mini-conclusion)}
\item \textbf{Utiliser les connecteurs logiques} : \esp{en primer lugar}, \esp{además}, \esp{sin embargo}, \esp{por eso / por tanto}, \esp{en conclusión / en definitiva}.
\item \textbf{Exprimer l'opinion} :
\begin{itemize}
\item \esp{Pienso / Creo / Considero que + INDICATIF} (certitude).
\item \esp{No creo / No pienso que + SUBJONCTIF} (doute/négation).
\item \esp{Me parece que + INDICATIF} (impression).
\end{itemize}
\item \textbf{Exprimer la nécessité / l'obligation} :
\begin{itemize}
\item \esp{Hay que + INFINITIF} (impersonnel) : \esp{Hay que verificar las fuentes}.
\item \esp{Es necesario / imprescindible / importante que + SUBJONCTIF} : \esp{Es necesario que los medios respeten}.
\item \esp{Debemos / Tenemos que + INFINITIF} (engagement personnel) : \esp{Debemos promover la tolerancia}.
\end{itemize}
\item \textbf{À l'oral (épreuve orale / EOC)} : Prépare 5 phrases-clés (thèse, 2 arguments, exemple local, conclusion). Parle lentement, articule les voyelles, prononce \esp{c} (devant e/i) comme « th » (norme péninsulaire) ou « s » (norme américaine/équatoriale), \esp{ll} comme « y » (yeísmo généralisé). Conclus par un appel collectif : \esp{¡Construyamos juntos la cultura de paz!} (subjonctif 1\up{re} pl. = impératif \emph{nosotros}).
\end{enumerate}
\end{methode}

\begin{exoresolu}[Production écrite guidée — type Bac]
\textbf{Énoncé :} Escribe un párrafo de 8-10 líneas sobre el tema : « El papel de los medios de comunicación en la promoción de la cultura de paz en nuestro país ».
\tcblower
\textbf{Raisonnement :} On applique le plan IAE, avec au moins trois mots du lexique de la leçon (\esp{cuarto poder}, \esp{opinión pública}, \esp{reconciliación}), un subjonctif justifié, des connecteurs variés et un ancrage camerounais.\par\smallskip
\textbf{Production modèle :}\par
\esp{En nuestra sociedad camerunesa, los medios de comunicación desempeñan un papel fundamental en la promoción de la cultura de paz. En primer lugar, el cuarto poder informa a los ciudadanos y forma la opinión pública: cuando los periodistas verifican sus fuentes, evitan los rumores que dividen a la población. Además, la radio comunitaria y la televisión pueden servir de mediadoras entre las comunidades, dando la palabra a todas las partes después de un conflicto. Por ejemplo, en las elecciones municipales, una prensa libre e independiente ayuda a aceptar los resultados y favorece la reconciliación nacional. Sin embargo, es imprescindible que los medios respeten la verdad y rechacen el discurso de odio. En conclusión, pienso que sin libertad de prensa no hay paz duradera: por eso debemos defender a los periodistas honestos y construir juntos una cultura de paz.}\par\smallskip
\textbf{Points de langue vérifiés :} \esp{es imprescindible que\ldots respeten} (subjonctif présent 3\up{e} pl.) ; \esp{pienso que\ldots hay} (indicatif après opinion affirmative) ; \esp{duradera} (adjectif accordé avec \esp{paz}, fém. sg.) ; accents : \esp{opinión}, \esp{televisión}, \esp{reconciliación}, \esp{camerunesa}, \esp{municipales}, \esp{honestos}.
\end{exoresolu}

\begin{experience}[Activité orale : « Tribuna por la paz » — 20 min]
\textbf{Consigne :} En groupes de 4. Chaque groupe tire un rôle : 1) \esp{Periodista de investigación}, 2) \esp{Dirigente comunitario}, 3) \esp{Estudiante universitario}, 4) \esp{Representante del Ministerio de Comunicación}.
\textbf{Sujet :} \esp{« ¿Cómo pueden los medios locales (radio CRTV, presse privée, réseaux sociaux) contribuer à la paix dans la région de l'Extrême-Nord ? »}
\textbf{Déroulement :} 5 min préparation (notes mots-clés, lexique de la leçon) ; 10 min débat modéré par l'enseignant ; 5 min restitution : un rapporteur par groupe résume les 3 propositions concrètes en espagnol. \textbf{Critères :} utilisation de \esp{mediador, reconciliación, libertad de prensa, verificar fuentes, opinión pública} ; subjonctif après \esp{es necesario que} ; connecteurs ; prononciation soignée.
\end{experience}

\courssub{Points de vigilance pour le Bac}

\begin{attention}[Les 5 erreurs qui coûtent des points au Bac]
\begin{enumerate}
\item \textbf{Oublier les accents obligatoires} : \esp{opinion}, \esp{informacion}, \esp{dialogo}, \esp{reconciliacion}, \esp{promocion}, \esp{nacion}, \esp{television}, \esp{tambien} sont \textbf{FAUX}. Écrire : \esp{opinión}, \esp{información}, \esp{diálogo}, \esp{reconciliación}, \esp{promoción}, \esp{nación}, \esp{televisión}, \esp{también}. Règle : mot aigu (accent tonique sur dernière syllabe) terminé par \textbf{-n, -s ou voyelle} $\to$ accent écrit.
\item \textbf{Les faux amis fréquents} :
\begin{itemize}
\item « Assister à » = \esp{asistir a} (jamais \esp{ayudar} = aider).
\item « Réaliser (un projet) » = \esp{realizar / llevar a cabo} ; « Réaliser (comprendre) » = \esp{darse cuenta de}.
\item « Éxito » = succès (pas \esp{exit} = sortie).
\item « Largo » = long (pas \esp{large} = ancho).
\item « Constipado » = enrhumé (pas \esp{estreñido} = constipé).
\end{itemize}
\item \textbf{Le calque du passé composé français avec marqueur de passé} : « Hier, ils ont publié » $\neq$ \esp{Ayer han publicado}. \textbf{OBLIGATOIRE :} \esp{Ayer publicaron} (\esp{pretérito indefinido}). Le \esp{pretérito perfecto} (\esp{han publicado}) s'utilise pour un passé récent \textbf{sans} marqueur de temps passé (ex. \esp{Hoy han publicado}).
\item \textbf{L'indicatif au lieu du subjonctif} : Après \esp{es importante que}, \esp{es necesario que}, \esp{es imprescindible que}, \esp{no creo que}, \esp{para que}, \esp{quiero que}, \esp{espero que} $\to$ \textbf{SUBJONCTIF} obligatoire. Ex. : \esp{Es necesario que los medios respeten} (pas \esp{respetan}).
\item \textbf{Les accords (genre / nombre) oubliés} : \esp{la opinión pública} (fém.), \esp{una prensa libre e independiente} (fém. sg., \esp{e} eufonique devant \emph{i}-), \esp{los medios respetan} (sujet pl. $\to$ verbe pl.), \esp{noticias falsas} (fém. pl.). Vérifie \textbf{chaque} adjectif et \textbf{chaque} verbe avant de rendre la copie.
\end{enumerate}
\end{attention}

\newpage

\coursec{Exercices}

\courssub{Je vérifie mes connaissances}

\begin{exercice}[QCM : le quatrième pouvoir]
Pour chaque question, choisis la bonne réponse (a, b ou c).
\begin{enumerate}
\item ¿Por qué se llama a la prensa \esp{el cuarto poder} ?
\begin{itemize}
\item[a)] Porque es el cuarto poder del Estado, después del legislativo, el ejecutivo y el judicial.
\item[b)] Porque existen cuatro tipos de medios de comunicación.
\item[c)] Porque fue el cuarto invento del siglo XX.
\end{itemize}
\item \esp{La libertad de prensa} significa :
\begin{itemize}
\item[a)] el derecho de los periodistas a escribir lo que quieran sin ninguna ley.
\item[b)] el derecho de informar sin censura, dentro del respeto de la ley y de la verdad.
\item[c)] el derecho de los políticos a controlar los periódicos.
\end{itemize}
\item \esp{La censura} es :
\begin{itemize}
\item[a)] la opinión de los lectores sobre un artículo.
\item[b)] la prohibición o el control de la información por el poder.
\item[c)] la publicidad en los medios.
\end{itemize}
\item Un \esp{mediador} en un conflicto es :
\begin{itemize}
\item[a)] una persona que toma partido por uno de los dos bandos.
\item[b)] una persona neutra que ayuda a las partes a dialogar y a reconciliarse.
\item[c)] un periodista que entrevista a los políticos.
\end{itemize}
\end{enumerate}
\end{exercice}

\begin{exercice}[Vrai ou faux, justifié]
D'après le texte support \esp{El cuarto poder al servicio de la paz}, dis si les affirmations suivantes sont vraies (\esp{verdadero}) ou fausses (\esp{falso}) et justifie ta réponse en citant une phrase du texte.
\begin{enumerate}
\item La prensa solo transmite noticias de deportes y de espectáculos.
\item Los medios de comunicación pueden influir en la opinión pública.
\item La censura favorece la cultura de paz.
\item Un periodista responsable verifica sus fuentes antes de publicar.
\item Los medios no pueden contribuir a la reconciliación después de un conflicto.
\item La independencia de la prensa es necesaria para la democracia.
\end{enumerate}
\end{exercice}

\begin{exercice}[Vocabulaire : trouve l'intrus]
Dans chaque série, un mot n'appartient pas au champ lexical. Trouve l'intrus et justifie ton choix en français.
\begin{enumerate}
\item el periódico — la revista — la emisora de radio — el ayuntamiento
\item la libertad de prensa — la censura — la independencia — la veracidad
\item el mediador — la reconciliación — el diálogo — la guerra
\item la opinión pública — los lectores — la audiencia — el ejército
\item el periodista — el redactor — el corresponsal — el carpintero
\end{enumerate}
\end{exercice}

\begin{exercice}[Conjugaison : complète le texte]
Complète le texte suivant en conjuguant les verbes entre parenthèses au présent de l'indicatif ou au présent du subjonctif, selon le cas.
\begin{quote}
\esp{Es importante que la prensa (informar) con objetividad. Cuando los periodistas (verificar) sus fuentes, la opinión pública (confiar) más en ellos. No creo que la censura (ser) una buena solución : es necesario que los medios (poder) trabajar con independencia. Así, la prensa (contribuir) a la reconciliación y a la cultura de paz.}
\end{quote}
\end{exercice}

\courssub{Je m'entraîne}

\begin{exercice}[Traduction : thème]
Traduis les phrases suivantes en espagnol.
\begin{enumerate}
\item La presse joue un rôle essentiel dans la promotion de la paix.
\item Il faut que les journalistes respectent la vérité et vérifient leurs informations.
\item Au Cameroun, les médias peuvent aider à la réconciliation entre les communautés.
\item Sans liberté de presse, il n'y a pas de vraie démocratie.
\item L'opinion publique est influencée par ce que publient les journaux.
\end{enumerate}
\end{exercice}

\begin{exercice}[Transformation de phrases]
Transforme les phrases suivantes selon la consigne entre parenthèses.
\begin{enumerate}
\item \esp{El gobierno censura los periódicos.} (mets à la voix passive : \esp{Los periódicos...})
\item \esp{El periodista dijo : « Verifico siempre mis fuentes.»} (mets au discours indirect au passé : \esp{El periodista dijo que...})
\item \esp{Los medios informan. La población quiere participar en la vida pública.} (relie les deux phrases avec \esp{para que} et fais les changements nécessaires)
\item \esp{Es posible que la prensa promueva la paz.} (transforme en phrase affirmative avec \esp{es cierto que})
\item \esp{La radio es un medio importante. Muchos cameruneses escuchan la radio cada mañana.} (relie avec un pronom relatif : \esp{que})
\end{enumerate}
\end{exercice}

\begin{exercice}[Texte à trous : un article de presse]
Complète l'article suivant avec les mots de la liste : \esp{cuarto poder — censura — opinión pública — mediador — reconciliación — libertad de prensa}.
\begin{textebox}[Un artículo a completar]
En toda democracia, la prensa es considerada el \dots\dots\dots\dots porque vigila a los gobernantes e informa a los ciudadanos. La \dots\dots\dots\dots permite a los periodistas publicar sus investigaciones sin miedo. Sin embargo, en algunos países, la \dots\dots\dots\dots impide la difusión de informaciones que molestan al poder. Cuando estalla un conflicto, un buen periódico puede actuar como \dots\dots\dots\dots entre las partes, dando la palabra a todos. Gracias a un periodismo responsable, la \dots\dots\dots\dots se forma con datos verificados y no con rumores. Así, los medios contribuyen a la \dots\dots\dots\dots de las comunidades divididas.
\end{textebox}
\end{exercice}

\begin{exercice}[Appariements : mots et définitions]
Associe chaque mot espagnol (colonne A) à sa définition espagnole (colonne B).
\begin{center}
\begin{tabularx}{\linewidth}{|X|X|}
\hline
\rowcolor{popblueL} \textbf{A. Palabra} & \textbf{B. Definición} \\
\hline
1. el titular & a) persona que escribe artículos en un periódico \\
\hline
2. el periodista & b) título de un artículo en letras grandes \\
\hline
3. la noticia falsa & c) conjunto de personas que escuchan un programa \\
\hline
4. la audiencia & d) información inventada para engañar al público \\
\hline
5. el reportaje & e) trabajo periodístico sobre un tema, con entrevistas \\
\hline
6. la rueda de prensa & f) reunión en la que un responsable responde a los periodistas \\
\hline
\end{tabularx}
\end{center}
\end{exercice}

\courssub{Je me prépare au Bac}

\begin{exercice}[Compréhension et questions de langue (type Bac)]
Lis le texte suivant, puis réponds aux questions.

\begin{textebox}[La libertad de prensa en África]
En muchos países africanos, la prensa ha conquistado, poco a poco, un espacio de libertad que no existía hace treinta años. Los periódicos independientes, las radios comunitarias y, más recientemente, los medios digitales informan cada día a millones de ciudadanos. Sin embargo, esta libertad sigue siendo frágil. Algunos periodistas son amenazados o detenidos por sus investigaciones ; otros medios cierran por falta de dinero, porque la publicidad es escasa.

A pesar de estas dificultades, la prensa africana desempeña un papel fundamental en la construcción de la paz. En los periodos electorales, las radios comunitarias explican a los electores cómo votar y llaman a la calma. Después de los conflictos, los programas de diálogo dan la palabra a las víctimas y a los antiguos combatientes, favoreciendo así la reconciliación. En Camerún, por ejemplo, muchas emisoras difunden mensajes de tolerancia en varias lenguas nacionales.

Pero el periodista también tiene deberes. Debe verificar sus fuentes, respetar a las personas y rechazar los discursos de odio. Una noticia falsa puede provocar violencia en pocas horas, sobre todo en las redes sociales. Por eso, la formación de los periodistas y la educación de los ciudadanos a los medios son indispensables.

El cuarto poder, cuando es libre y responsable, no es un enemigo del Estado : es un aliado de la democracia y de la paz.
\end{textebox}

\textbf{I. Comprensión del texto} (réponds en espagnol, avec des phrases complètes)
\begin{enumerate}
\item ¿Qué tres tipos de medios de comunicación menciona el texto ?
\item ¿Por qué dice el autor que la libertad de prensa en África es « frágil » ? Da dos razones.
\item ¿Qué hacen las radios comunitarias en los periodos electorales ?
\item ¿Qué deberes tiene el periodista, según el texto ?
\item Explica con tus palabras la última frase del texto : \esp{« es un aliado de la democracia y de la paz »}.
\end{enumerate}

\textbf{II. Pregunta de opinión personal} (5 à 8 lignes, en espagnol)

¿Crees que las redes sociales son un peligro o una oportunidad para la libertad de prensa en tu país ? Justifica tu respuesta con ejemplos.

\textbf{III. Preguntas de lengua}
\begin{enumerate}
\item Busca en el texto dos sinónimos de \esp{informar} y un antónimo de \esp{libertad}.
\item Conjuga los verbos entre paréntesis : \esp{Es necesario que los medios (respetar) la verdad ; no creo que una noticia falsa (ayudar) a la reconciliación.}
\item Pon a la voz pasiva : \esp{Las radios comunitarias difunden mensajes de tolerancia.}
\end{enumerate}
\end{exercice}

\begin{exercice}[Thème et version (type Bac)]
\textbf{A. Thème.} Traduis en espagnol le texte suivant.

\begin{quote}
Dans mon pays, la presse est de plus en plus libre. Les journalistes peuvent critiquer le gouvernement, mais ils doivent respecter la vérité. Il est important que les médias ne publient pas de fausses nouvelles, car elles peuvent diviser la population. Quand il y a un conflit, les journalistes doivent jouer le rôle de médiateurs et promouvoir la réconciliation. Je crois que sans presse libre, la paix est impossible.
\end{quote}

\textbf{B. Version.} Traduis en français le texte suivant.

\begin{textebox}[Para traducir]
La opinión pública se forma gracias a la información que recibe cada día. Si los medios son independientes, los ciudadanos pueden conocer la verdad y participar en la vida de su país. Pero cuando existe la censura, nadie sabe lo que ocurre realmente. Por eso los periodistas defienden la libertad de prensa : saben que informar bien es servir a la paz.
\end{textebox}
\end{exercice}

\begin{exercice}[Rédaction et expression orale (type Bac)]
\textbf{A. Rédaction guidée} (12 à 15 lignes, en espagnol)

Sujet : \esp{¿Es la prensa el cuarto poder en tu país ?}

Respecte le plan imposé suivant :
\begin{enumerate}
\item \textbf{Introducción} : presenta el tema y anuncia tu opinión.
\item \textbf{Desarrollo} : da dos argumentos con ejemplos concretos (medios de tu país, papel en las elecciones, conflictos, reconciliación).
\item \textbf{Conclusión} : resume tu posición y propone una recomendación para que la prensa sirva mejor a la cultura de paz.
\end{enumerate}
Utilise au moins cinq mots du lexique de la leçon : \esp{cuarto poder, libertad de prensa, censura, opinión pública, mediador, reconciliación, independencia}.

\textbf{B. Expression orale} (présentation de 3 minutes)

Prépare une présentation orale sur le sujet suivant : \esp{El papel de un medio de comunicación camerunés o hispanohablante en la promoción de la cultura de paz.} Tu présenteras le média choisi (radio, télévision, journal ou média en ligne), ses actions concrètes en faveur de la paix, puis tu donneras ton avis personnel. Après ta présentation, le jury te posera deux questions : sois prêt à défendre ton point de vue et à donner des exemples.
\end{exercice}

\newpage

\coursec{Corrigés des exercices}

\begin{corrige}[Exercice 1 — QCM : le quatrième pouvoir]
\begin{enumerate}
\item \textbf{Réponse a)}. La presse est appelée \esp{el cuarto poder} parce qu'elle s'ajoute aux trois pouvoirs classiques de l'État : le législatif (\esp{el poder legislativo}), l'exécutif (\esp{el poder ejecutivo}) et le judiciaire (\esp{el poder judicial}). Elle n'en fait pas officiellement partie, mais elle les surveille (\esp{vigila}) et influence la vie politique.
\item \textbf{Réponse b)}. \cle{La libertad de prensa} est le droit d'informer sans censure, dans le respect de la loi et de la vérité. La réponse a) est fausse : la liberté n'est pas l'anarchie, il existe des lois (diffamation \esp{la difamación}, respect de la vie privée \esp{el derecho a la intimidad}).
\item \textbf{Réponse b)}. \cle{La censura} est l'interdiction ou le contrôle préventif de l'information par le pouvoir politique.
\item \textbf{Réponse b)}. Un \cle{mediador} est une personne neutre et impartiale qui aide les parties en conflit à dialoguer pour trouver une solution pacifique (\esp{la reconciliación}). La réponse a) décrit un partisan (\esp{un partidario}), pas un médiateur.
\end{enumerate}
\textbf{Points de vigilance} : ne pas confondre \esp{mediador} (médiateur) et \esp{periodista} (journaliste) ; attention au faux ami \esp{censura}, qui signifie bien « censure » (contrôle politique) et non « critique de film » (\esp{crítica}). \esp{El cuarto poder} s'écrit en minuscules (sauf en début de phrase), ce n'est pas un nom propre.
\end{corrige}

\begin{corrige}[Exercice 2 — Vrai ou faux, justifié]
\begin{enumerate}
\item \textbf{Falso.} La presse transmet toutes sortes d'informations : politique, économie, société, culture, et pas seulement des nouvelles de sports et de spectacles. (Justification attendue : citer une phrase du texte indiquant que la presse informe sur la vie publique et surveille les gouvernants, ex. \esp{« La prensa informa a los ciudadanos sobre la vida pública y vigila a los gobernantes »}.)
\item \textbf{Verdadero.} Les médias influencent l'opinion publique (\esp{la opinión pública}) : c'est précisément pour cela qu'on les appelle \esp{el cuarto poder}.
\item \textbf{Falso.} La censure empêche la libre circulation de l'information ; elle favorise les rumeurs (\esp{los rumores}) et les tensions, pas la culture de paix (\esp{la cultura de paz}).
\item \textbf{Verdadero.} Un journaliste responsable (\esp{un periodista responsable}) vérifie ses sources (\esp{verifica sus fuentes}) avant de publier : c'est un devoir professionnel essentiel (\esp{deber profesional}).
\item \textbf{Falso.} Au contraire, les médias peuvent contribuer à la réconciliation (\esp{la reconciliación}) après un conflit, en donnant la parole à toutes les parties (\esp{dar la palabra a todas las partes}) et en diffusant des messages de tolérance (\esp{mensajes de tolerancia}).
\item \textbf{Verdadero.} L'indépendance de la presse (\esp{la independencia de la prensa}) est une condition de la démocratie : sans presse libre (\esp{sin prensa libre}), les citoyens ne peuvent pas connaître la vérité.
\end{enumerate}
\textbf{Méthode} : pour justifier, cite toujours une phrase du texte entre guillemets espagnols, introduite par \esp{El texto dice: « ... »} ou \esp{Como afirma el autor: « ... »}.
\end{corrige}

\begin{corrige}[Exercice 3 — Vocabulaire : trouve l'intrus]
\begin{enumerate}
\item \textbf{el ayuntamiento} : c'est la mairie, une institution politique locale ; les trois autres sont des médias : \esp{el periódico} (journal), \esp{la revista} (magazine), \esp{la emisora de radio} (station de radio).
\item \textbf{la censura} : c'est une atteinte à la liberté ; les trois autres mots désignent des valeurs positives du journalisme : \esp{la libertad}, \esp{la independencia}, \esp{la veracidad}.
\item \textbf{la guerra} : c'est le contraire de la paix ; les trois autres appartiennent au champ lexical de la résolution pacifique des conflits : \esp{el mediador}, \esp{la reconciliación}, \esp{el diálogo}.
\item \textbf{el ejército} : c'est une institution militaire ; les trois autres désignent le public des médias : \esp{la opinión pública}, \esp{los lectores}, \esp{la audiencia}.
\item \textbf{el carpintero} : c'est un menuisier/charpentier (métier du bois) ; les trois autres sont des métiers du journalisme : \esp{el periodista}, \esp{el redactor}, \esp{el corresponsal}.
\end{enumerate}
\textbf{Points de vigilance} : \esp{el carpintero} est un faux ami partiel : il signifie « menuisier/charpentier », pas « charpentier » au sens figuré d'architecte. \esp{La audiencia} signifie « l'audience, le public (radio/TV) » mais aussi « l'audience » d'un tribunal (audience judiciaire).
\end{corrige}

\begin{corrige}[Exercice 4 — Conjugaison : complète le texte]
Texte complété :
\begin{quote}
\esp{Es importante que la prensa \textbf{informe} con objetividad. Cuando los periodistas \textbf{verifican} sus fuentes, la opinión pública \textbf{confía} más en ellos. No creo que la censura \textbf{sea} una buena solución: es necesario que los medios \textbf{puedan} trabajar con independencia. Así, la prensa \textbf{contribuye} a la reconciliación y a la cultura de paz.}
\end{quote}
Traduction : « Il est important que la presse informe avec objectivité. Quand les journalistes vérifient leurs sources, l'opinion publique leur fait davantage confiance. Je ne crois pas que la censure soit une bonne solution : il est nécessaire que les médias puissent travailler avec indépendance. Ainsi, la presse contribue à la réconciliation et à la culture de paix. »

Justifications grammaticales :
\begin{itemize}
\item \esp{informe} : subjonctif présent (\esp{informe}) après l'expression impersonnelle \esp{es importante que} (jugement de valeur, obligation morale).
\item \esp{verifican} : indicatif présent après \esp{cuando} pour exprimer une habitude, un fait réel et récurrent (pas de futur ni d'incertitude).
\item \esp{confía} : indicatif présent, constatation réelle. Verbe \esp{confiar} : accent graphique sur le \textbf{í} (\esp{confío, confías, confía...}) pour rompre le diphtongue (hiatus \emph{i} + \emph{a}), pas de diphtongaison radicale.
\item \esp{sea} : subjonctif présent de \esp{ser} (irrégulier : \esp{sea, seas, sea, seamos, seáis, sean}) après \esp{no creo que} (négation de la pensée $\rightarrow$ doute/subjonctif).
\item \esp{puedan} : subjonctif présent de \esp{poder} (diphtongaison radicale \emph{o} $\rightarrow$ \emph{ue} : \esp{pueda, puedas, pueda, podamos, podáis, puedan}) après \esp{es necesario que} (impersonnel de nécessité).
\item \esp{contribuye} : indicatif présent, conséquence réelle. Verbe \esp{contribuir} (régulier en \emph{-ir}) : \esp{contribuyo, contribuyes, contribuye, contribuimos, contribuís, contribuyen}. Attention : pas de diphtongaison, mais ajout d'un \emph{y} graphique aux 3 personnes du singulier et 3\textsuperscript{e} pluriel pour éviter \emph{ii} (\esp{contribuye}, \esp{contribuyen}).
\end{itemize}
\end{corrige}

\begin{corrige}[Exercice 5 — Traduction : thème]
\begin{enumerate}
\item \esp{La prensa desempeña un papel esencial en la promoción de la paz.} (On peut aussi dire : \esp{juega un papel esencial}, mais \esp{desempeñar un papel} est plus soutenu/administratif.)
\item \esp{Es necesario que los periodistas respeten la verdad y verifiquen sus informaciones.} Subjonctif obligatoire après \esp{es necesario que}. Attention : \esp{verifiquen} s'écrit avec \textbf{qu} devant \textbf{e} pour garder le son [k] (règle \emph{c} $\rightarrow$ \emph{qu} devant \emph{e}, \emph{i}).
\item \esp{En Camerún, los medios de comunicación pueden ayudar a la reconciliación entre las comunidades.} Attention à l'accent tonique sur le \textbf{ú} : \esp{Camerún} (oxytone terminée par \emph{n}).
\item \esp{Sin libertad de prensa, no hay verdadera democracia.} (ou \esp{no existe una verdadera democracia}). \esp{Hay} (impersonnel, invariable) = « il y a ».
\item \esp{La opinión pública está influenciada por lo que publican los periódicos.} On peut aussi écrire \esp{es influenciada} (passif d'état), mais \esp{estar + participio} insiste sur l'état résultant. \esp{Lo que} (neutre) traduit « ce que » (antécédent indéfini).
\end{enumerate}
\textbf{Points de vigilance} : « les médias » se traduit \cle{los medios} ou \cle{los medios de comunicación}, \textbf{jamais} \esp{las medias} (= les bas, les chaussettes) ; « une fausse nouvelle » = \cle{una noticia falsa}, pas \esp{una nueva falsa} (\esp{nueva} = neuve, adjectif qualificatif).
\end{corrige}

\begin{corrige}[Exercice 6 — Transformation de phrases]
\begin{enumerate}
\item \esp{Los periódicos \textbf{son censurados por} el gobierno.} Voix passive \emph{propre} : \esp{ser} (conjugué au même temps que l'actif, ici présent) + participe passé accordé en genre et nombre (\esp{censurados}, masc. pl. comme \esp{periódicos}) + \esp{por} + complément d'agent.
\item \esp{El periodista dijo que \textbf{verificaba} siempre sus fuentes.} Discours indirect au passé : la principale (\esp{dijo}, ind. passé) attire le verbe de la subordonnée : présent \esp{verifico} $\rightarrow$ imparfait \esp{verificaba}. Suppression des guillemets, changement de personne (1\textsuperscript{e} $\rightarrow$ 3\textsuperscript{e}).
\item \esp{Los medios informan \textbf{para que} la población \textbf{pueda} participer en la vida pública.} \esp{Para que} exprime le but (finalité) et exige le subjonctif ; on ajoute le verbe \esp{poder} pour donner un sujet grammatical à la subordonnée (sujets différents : \esp{los medios} / \esp{la población}). \esp{Pueda} = subj. présent 3\textsuperscript{e} sg de \esp{poder}.
\item \esp{\textbf{Es cierto que} la prensa \textbf{promueve} la paz.} \esp{Es cierto que} exprime la certitude $\rightarrow$ indicatif (\esp{promueve}) ; \esp{es posible que} exprimait la possibilité $\rightarrow$ subjonctif (\esp{promueva}) dans la phrase de départ.
\item \esp{La radio es un medio importante \textbf{que} muchos cameruneses escuchan cada mañana.} Le pronom relatif \esp{que} reprend \esp{un medio importante} (COD de \esp{escuchan}). Attention : \esp{cameruneses} (gentilé correct, sans accent) ; \esp{muchos} accordé au masculin pluriel.
\end{enumerate}
\end{corrige}

\begin{corrige}[Exercice 7 — Texte à trous : un article de presse]
Texte complété :
\begin{quote}
\esp{En toda democracia, la prensa es considerada el \textbf{cuarto poder} porque vigila a los gobernantes e informa a los ciudadanos. La \textbf{libertad de prensa} permite a los periodistas publicar sus investigaciones sin miedo. Sin embargo, en algunos países, la \textbf{censura} impide la difusión de informaciones que molestan al poder. Cuando estalla un conflicto, un buen periódico puede actuar como \textbf{mediador} entre las partes, dando la palabra a todos. Gracias a un periodismo responsable, la \textbf{opinión pública} se forma con datos verificados y no con rumores. Así, los medios contribuyen a la \textbf{reconciliación} de las comunidades divididas.}
\end{quote}
Traduction : « Dans toute démocratie, la presse est considérée comme le quatrième pouvoir parce qu'elle surveille les gouvernants et informe les citoyens. La liberté de presse permet aux journalistes de publier leurs enquêtes sans peur. Cependant, dans certains pays, la censure empêche la diffusion d'informations qui dérangent le pouvoir. Quand un conflit éclate, un bon journal peut agir comme médiateur entre les parties, en donnant la parole à tous. Grâce à un journalisme responsable, l'opinion publique se forme avec des données vérifiées et non avec des rumeurs. Ainsi, les médias contribuent à la réconciliation des communautés divisées. »

\textbf{Indices de résolution} : « vigila a los gobernantes » $\rightarrow$ \esp{cuarto poder} ; « sin miedo » $\rightarrow$ \esp{libertad de prensa} ; « impide la difusión » $\rightarrow$ \esp{censura} ; « entre las partes » $\rightarrow$ \esp{mediador} ; « se forma » $\rightarrow$ \esp{opinión pública} ; « comunidades divididas » $\rightarrow$ \esp{reconciliación}.
\end{corrige}

\begin{corrige}[Exercice 8 — Appariements : mots et définitions]
\begin{center}
\begin{tabularx}{\linewidth}{|X|X|}
\hline
\rowcolor{popblueL} \textbf{A. Palabra} & \textbf{B. Definición correcta} \\
\hline
1. el titular & b) título de un artículo en letras grandes \\
\hline
2. el periodista & a) persona que escribe artículos en un periódico \\
\hline
3. la noticia falsa & d) información inventada para engañar al público \\
\hline
4. la audiencia & c) conjunto de personas que escuchan un programa \\
\hline
5. el reportaje & e) trabajo periodístico sobre un tema, con entrevistas \\
\hline
6. la rueda de prensa & f) reunión en la que un responsable responde a los periodistas \\
\hline
\end{tabularx}
\end{center}
Traductions des mots : \esp{el titular} = le titre (d'un article, en gros caractères) ; \esp{la noticia falsa} = la fausse nouvelle (\emph{fake news}) ; \esp{la audiencia} = l'audience (public) ; \esp{el reportaje} = le reportage ; \esp{la rueda de prensa} = la conférence de presse (litt. « ronde de presse »).
\end{corrige}

\begin{corrige}[Exercice 9 — Compréhension et questions de langue (type Bac)]
\textbf{Barème indicatif (sur 20)} : compréhension 8 points (1,5 point par question 1 à 4, 2 points pour la question 5) ; opinion personnelle 6 points ; questions de langue 6 points (2 points par question).

\textbf{I. Comprensión del texto} — réponses modèles :
\begin{enumerate}
\item \esp{El texto menciona tres tipos de medios: los periódicos independientes, las radios comunitarias y los medios digitales.}
\item \esp{El autor dice que la libertad de prensa es frágil por dos razones: primero, algunos periodistas son amenazados o detenidos por sus investigaciones; segundo, algunos medios cierran por falta de dinero, porque la publicidad es escasa.}
\item \esp{En los periodos electorales, las radios comunitarias explican a los electores cómo votar y llaman a la calma.}
\item \esp{Según el texto, el periodista debe verificar sus fuentes, respetar a las personas y rechazar los discursos de odio.}
\item \esp{Esta frase significa que una prensa libre y responsable ayuda a la democracia, porque informa a los ciudadanos y vigila a los gobernantes, y ayuda a la paz, porque favorece el diálogo y la reconciliación. Por eso no es una enemiga del Estado, sino una aliada.}
\end{enumerate}

\textbf{II. Pregunta de opinión personal} — réponse modèle :
\begin{quote}
\esp{En mi opinión, las redes sociales son a la vez un peligro y una oportunidad para la libertad de prensa en mi país. Por un lado, son una oportunidad porque permiten a los ciudadanos informarse rápidamente y expresar sus opiniones libremente, incluso cuando los medios tradicionales están controlados. Por otro lado, son un peligro porque las noticias falsas se difunden muy rápido y pueden provocar violencia, sobre todo en los periodos electorales. En Camerún, por ejemplo, circulan a veces rumores que dividen a las comunidades. Por eso creo que es necesario educar a los jóvenes en el uso crítico de las redes sociales y verificar siempre las informaciones antes de compartirlas.}
\end{quote}
Traduction : « À mon avis, les réseaux sociaux sont à la fois un danger et une opportunité pour la liberté de presse dans mon pays. D'un côté, ils sont une opportunité car ils permettent aux citoyens de s'informer rapidement et d'exprimer leurs opinions librement, même quand les médias traditionnels sont contrôlés. De l'autre, ils sont un danger car les fausses nouvelles se diffusent très vite et peuvent provoquer des violences, surtout en période électorale. Au Cameroun, par exemple, circulent parfois des rumeurs qui divisent les communautés. C'est pourquoi je crois qu'il faut éduquer les jeunes à un usage critique des réseaux sociaux et toujours vérifier les informations avant de les partager. »

\textbf{III. Preguntas de lengua} :
\begin{enumerate}
\item Deux synonymes d'\esp{informar} dans le texte : \esp{difundir} et \esp{explicar} (on accepte aussi \esp{comunicar}). Un antonyme de \esp{libertad} : \esp{censura} (sous-entendu dans le texte) ; on accepte aussi \esp{represión} ou \esp{control} si le candidat justifie.
\item \esp{Es necesario que los medios \textbf{respeten} la verdad; no creo que una noticia falsa \textbf{ayude} a la reconciliación.} Subjonctif présent (\esp{respeten}, \esp{ayude}) après \esp{es necesario que} et après \esp{no creo que}.
\item \esp{Mensajes de tolerancia \textbf{son difundidos por} las radios comunitarias.} Voix passive : participe passé \esp{difundidos} accordé au masculin pluriel (comme \esp{mensajes}).
\end{enumerate}
\textbf{Points de vigilance} : répondre par des phrases complètes en reprenant les mots de la question ; ne pas recopier tout un paragraphe ; pour l'opinion personnelle, annoncer clairement sa position dès la première phrase (\esp{En mi opinión...}, \esp{Yo creo que...}).
\end{corrige}

\begin{corrige}[Exercice 10 — Thème et version (type Bac)]
\textbf{Barème indicatif} : thème 10 points (2 points par phrase, moins 0,5 par faute de grammaire, 0,25 par faute d'accent ou de vocabulaire) ; version 10 points (fidélité 6, qualité du français 4).

\textbf{A. Thème} — traduction modèle :
\begin{quote}
\esp{En mi país, la prensa es cada vez más libre. Los periodistas pueden criticar al gobierno, pero deben respetar la verdad. Es importante que los medios no publiquen noticias falsas, porque pueden dividir a la población. Cuando hay un conflicto, los periodistas deben desempeñar el papel de mediadores y promover la reconciliación. Creo que sin prensa libre, la paz es imposible.}
\end{quote}
Points de grammaire vérifiés :
\begin{itemize}
\item « de plus en plus libre » = \esp{cada vez más libre} (locution figée ; \emph{ne pas dire} \esp{más y más libre}, tournure non idiomatique, calque de l'anglais \emph{more and more}).
\item « critiquer le gouvernement » = \esp{criticar \textbf{al} gobierno} : \textbf{a personale} obligatoire devant COD personne (ou animal domestique/entité personifiée).
\item « il est important que les médias ne publient pas » = \esp{es importante que los medios no \textbf{publiquen}} : subjonctif présent (\esp{publiquen}, 3\textsuperscript{e} pl.) obligatoire ; \esp{qu} pour garder [k] devant \emph{e}.
\item « promouvoir » = \esp{promover} (verbe régulier en \emph{-er}, radical \esp{promov-}) ; ne pas confondre avec \esp{promocionar} (promouvoir un produit/événement).
\item « je crois que » + indicatif : \esp{creo que... la paz \textbf{es} imposible} (certitude $\rightarrow$ indicatif).
\end{itemize}

\textbf{B. Version} — traduction modèle :
\begin{quote}
L'opinion publique se forme grâce à l'information qu'elle reçoit chaque jour. Si les médias sont indépendants, les citoyens peuvent connaître la vérité et participer à la vie de leur pays. Mais quand la censure existe, personne ne sait ce qui se passe réellement. C'est pourquoi les journalistes défendent la liberté de presse : ils savent que bien informer, c'est servir la paix.
\end{quote}
Points de vigilance : \esp{lo que ocurre} = « ce qui se passe » (pronom neutre \esp{lo que} = « ce qui/ce que ») ; \esp{por eso} = « c'est pourquoi » ; \esp{informar bien es servir a la paz} (infinitif sujet) se traduit élégamment par une infinitive française : « bien informer, c'est servir la paix ».
\end{corrige}

\begin{corrige}[Exercice 11 — Rédaction et expression orale (type Bac)]
\textbf{Barème indicatif (rédaction, sur 20)} : respect du plan et de la longueur (150--180 mots) 4 points ; richesse du vocabulaire (au moins 5 mots du lexique de la leçon) 4 points ; correction grammaticale et orthographique 8 points ; cohérence et force des arguments 4 points.

\textbf{A. Rédaction guidée} — proposition de rédaction modèle :
\begin{quote}
\esp{En mi opinión, la prensa merece ser llamada el cuarto poder en mi país, aunque su libertad todavía tiene límites. Voy a explicar mi posición con dos argumentos.

En primer lugar, la prensa camerunesa informa cada día a la opinión pública sobre la vida política, económica y social. Los periódicos, las radios y las televisiones comentan las decisiones del gobierno y dan la palabra a los ciudadanos. Durante las elecciones, por ejemplo, los periodistas explican el proceso de voto y llaman a la calma, lo que demuestra su influencia real.

En segundo lugar, la prensa puede actuar como mediador en los conflictos y contribuir a la reconciliación. Cuando hay tensiones entre comunidades, algunos programas de radio difunden mensajes de tolerancia en varias lenguas nacionales. Sin embargo, la libertad de prensa sigue siendo frágil: a veces existe la censura y algunos periodistas no pueden trabajar con total independencia.

En conclusión, la prensa es un verdadero cuarto poder en mi país, pero debe ser más libre y más responsable. Para servir mejor a la cultura de paz, recomiendo que los periodistas verifiquen siempre sus fuentes y que el Estado garantice la independencia de los medios.}
\end{quote}
Traduction : « À mon avis, la presse mérite d'être appelée le quatrième pouvoir dans mon pays, même si sa liberté a encore des limites. Je vais expliquer ma position avec deux arguments. Premièrement, la presse camerounaise informe chaque jour l'opinion publique sur la vie politique, économique et sociale. Les journaux, les radios et les télévisions commentent les décisions du gouvernement et donnent la parole aux citoyens. Pendant les élections, par exemple, les journalistes expliquent le processus de vote et appellent au calme, ce qui démontre leur influence réelle. Deuxièmement, la presse peut agir comme médiateur dans les conflits et contribuer à la réconciliation. Quand il y a des tensions entre communautés, certains programmes de radio diffusent des messages de tolérance en plusieurs langues nationales. Cependant, la liberté de presse reste fragile : la censure existe parfois et certains journalistes ne peuvent pas travailler en toute indépendance. En conclusion, la presse est un véritable quatrième pouvoir dans mon pays, mais elle doit être plus libre et plus responsable. Pour mieux servir la culture de paix, je recommande que les journalistes vérifient toujours leurs sources et que l'État garantisse l'indépendance des médias. »

Lexique de la leçon employé (7 mots, minimum exigé : 5) : \cle{cuarto poder}, \cle{opinión pública}, \cle{mediador}, \cle{reconciliación}, \cle{libertad de prensa}, \cle{censura}, \cle{independencia}.

Connecteurs utilisés : \esp{en primer lugar, en segundo lugar, sin embargo, en conclusión, por ejemplo}. Pense à les employer dans ta propre copie.

\textbf{B. Expression orale} — trame de présentation modèle (3 minutes) :
\begin{enumerate}
\item \textbf{Presentación del medio} : \esp{He elegido hablar de una radio comunitaria de mi región. Es una emisora que difunde programas en español, en francés y en lenguas nacionales.}
\item \textbf{Acciones en favor de la paz} : \esp{Esta radio organiza debates entre jóvenes de diferentes comunidades, difunde mensajes de tolerancia durante las elecciones y da la palabra a las víctimas de los conflictos para favorecer la reconciliación.}
\item \textbf{Opinión personal} : \esp{En mi opinión, este medio desempeña un papel esencial porque llega a la población rural y lucha contra los rumores y las noticias falsas. Creo que es un verdadero mediador de la paz.}
\item \textbf{Conclusión} : \esp{Para terminar, diría que una radio libre y responsable es una aliada de la democracia y de la cultura de paz.}
\end{enumerate}
Conseils pour l'oral : parle lentement, regarde le jury, utilise des connecteurs (\esp{primero, además, sin embargo, para terminar}), et prépare des exemples concrets (noms d'émissions, faits précis) pour répondre aux questions du jury. Évite de réciter par cœur : apprends le plan et les mots-clés, pas les phrases toutes faites.
\end{corrige}

\newpage

\coursec{Fiche bilan}

\begin{popfigure}
\begin{center}\tcbox[colback=poporange,colframe=poporange,coltext=white,arc=9pt,boxsep=4pt,left=6pt,right=6pt]{\bfseries\small El cuarto poder y la cultura de paz}\end{center}
\vspace{-2pt}
\begin{tcbraster}[raster columns=3,raster equal height=rows,raster column skip=6pt,raster row skip=6pt,enhanced,blankest]
\begin{tcolorbox}[enhanced,colback=poporange,colframe=poporange,coltext=white,boxrule=0.9pt,arc=3pt,left=4pt,right=4pt,top=3pt,bottom=3pt,boxsep=1pt,halign=left]\bfseries\scriptsize Lexique clé \textit{cuarto poder, censura, mediador}\end{tcolorbox}
\begin{tcolorbox}[enhanced,colback=popgreen,colframe=popgreen,coltext=white,boxrule=0.9pt,arc=3pt,left=4pt,right=4pt,top=3pt,bottom=3pt,boxsep=1pt,halign=left]\bfseries\scriptsize Textes argumentatifs \textit{opinión, tesis, argumentos}\end{tcolorbox}
\begin{tcolorbox}[enhanced,colback=poppurple,colframe=poppurple,coltext=white,boxrule=0.9pt,arc=3pt,left=4pt,right=4pt,top=3pt,bottom=3pt,boxsep=1pt,halign=left]\bfseries\scriptsize Liberté de presse \textit{independencia, pluralismo}\end{tcolorbox}
\begin{tcolorbox}[enhanced,colback=poppink,colframe=poppink,coltext=white,boxrule=0.9pt,arc=3pt,left=4pt,right=4pt,top=3pt,bottom=3pt,boxsep=1pt,halign=left]\bfseries\scriptsize Censure et indépendance \textit{autocensura, ética}\end{tcolorbox}
\begin{tcolorbox}[enhanced,colback=popgold,colframe=popgold,coltext=white,boxrule=0.9pt,arc=3pt,left=4pt,right=4pt,top=3pt,bottom=3pt,boxsep=1pt,halign=left]\bfseries\scriptsize Médiateur et réconciliation \textit{diálogo, consenso}\end{tcolorbox}
\begin{tcolorbox}[enhanced,colback=popblue,colframe=popblue,coltext=white,boxrule=0.9pt,arc=3pt,left=4pt,right=4pt,top=3pt,bottom=3pt,boxsep=1pt,halign=left]\bfseries\scriptsize Culture de la paix \textit{no violencia, tolerancia}\end{tcolorbox}
\begin{tcolorbox}[enhanced,colback=popmuted,colframe=popmuted,coltext=white,boxrule=0.9pt,arc=3pt,left=4pt,right=4pt,top=3pt,bottom=3pt,boxsep=1pt,halign=left]\bfseries\scriptsize Commentaire de texte \textit{structure, analyse}\end{tcolorbox}
\end{tcbraster}

\legende{Carte mentale de la leçon E24 : nœuds principaux et vocabulaire associé}
\end{popfigure}

\begin{bilan}

\textbf{Lexique clé (espagnol / français)} \\
\begin{itemize}
\item \esp{el cuarto poder} — le quatrième pouvoir (médias)
\item \esp{la libertad de prensa} — la liberté de la presse
\item \esp{la independencia (editorial)} — l'indépendance (éditoriale)
\item \esp{la censura / la autocensura} — la censure / l'autocensure
\item \esp{la opinión pública} — l'opinion publique
\item \esp{el mediador / la mediación} — le médiateur / la médiation
\item \esp{la reconciliación} — la réconciliation
\item \esp{la cultura de paz} — la culture de paix
\item \esp{la no violencia} — la non-violence
\item \esp{el pluralismo informativo} — le pluralisme de l'information
\item \esp{la ética periodística} — l'éthique journalistique
\item \esp{la posverdad / las fake news} — la post-vérité / les infox
\end{itemize}

\textbf{Points de grammaire à maîtriser (révision ciblée)} \\
\begin{center}
\begin{tabularx}{\linewidth}{|X|X|X|}
\hline
\rowcolor{popblueL}
\textbf{Notion} & \textbf{Forme / Emploi} & \textbf{Exemple} \\
\hline
Connecteurs argumentatifs & \esp{por tanto, en consecuencia, sin embargo, no obstante, además, en primer lugar} & \esp{La prensa es libre; \textbf{por tanto}, debe informar con rigor.} \\
\hline
Subjonctif dans l'opinion & \esp{es importante que + subj., es necesario que + subj., dudo que + subj.} & \esp{Es fundamental que los medios \textbf{sean} independientes.} \\
\hline
Voix passive réflexe (impersonnelle) & \esp{se + verbe (3e pers.)} pour objectivité & \esp{En el artículo \textbf{se denuncia} la censura.} \\
\hline
Discours rapporté (style indirect) & Changement de temps, pronoms, adverbes & \esp{«La prensa es libre» $\rightarrow$ Dijo que la prensa \textbf{era} libre.} \\
\hline
Gérondif (\textit{gerundio}) pour cause/moyen & \esp{al + infinitif, al + gérondif, mediante + gérondif} & \esp{\textbf{Al promover} el diálogo, los medios fomentan la paz.} \\
\hline
\end{tabularx}
\end{center}

\textbf{Méthode express : Commentaire de texte (4 étapes)} \\
\begin{enumerate}
\item \textbf{Lecture active} : identifier le thème, la thèse, le ton, la source, le public visé.
\item \textbf{Analyse structurelle} : repérer l'introduction (accroche, thèse), le développement (arguments, exemples, connecteurs), la conclusion (ouverture, appel).
\item \textbf{Analyse linguistique} : lexique thématique, figures de style, temps verbaux, modalisateurs, connecteurs.
\item \textbf{Rédaction} : introduction (présentation + thèse + plan), développement (2-3 parties équilibrées), conclusion (bilan + ouverture personnelle).
\end{enumerate}

\textbf{Méthode express : Production orale « Défendre la culture de la paix »} \\
\begin{enumerate}
\item Accroche : citation, chiffre, anecdote (ex. \esp{«La paz no es ausencia de conflicto, es la capacidad de resolverlo sin violencia»}).
\item Thèse claire : les médias comme acteurs de paix.
\item 2 arguments illustrés (ex. : journalisme de solutions, médias communautaires au Cameroun, rôle de la radio en Guinée équatoriale).
\item Contre-argument + réfutation (ex. : risque de propagande $\rightarrow$ éthique déontologique).
\item Conclusion engageante : appel à l'action citoyenne.
\end{enumerate}

\end{bilan}

\begin{aretenir}[Checklist avant le Bac]
\begin{itemize}
\item $\square$ Je connais le lexique du quatrième pouvoir (12 mots minimum) et je sais l'employer dans une phrase.
\item $\square$ Je sais définir \esp{cuarto poder}, \esp{libertad de prensa}, \esp{censura}, \esp{mediador}, \esp{cultura de paz} en espagnol.
\item $\square$ Je reconnais la structure d'un texte argumentatif (thèse, arguments, connecteurs, conclusion).
\item $\square$ J'utilise correctement les connecteurs logiques (\esp{por tanto, sin embargo, además, en cambio}).
\item $\square$ Je maîtrise le subjonctif après expressions d'opinion/nécessité (\esp{es vital que, exijo que, dudo que}).
\item $\square$ Je transforme le style direct en style indirect sans erreur de concordance des temps.
\item $\square$ J'emploie la voix passive réflexe (\esp{se + verbe}) pour donner de l'objectivité.
\item $\square$ Je sais rédiger une introduction de commentaire (présentation, thèse, plan annoncé).
\item $\square$ Je sais produire un paragraphe argumenté avec exemple concret (Cameroun, Espagne, Amérique latine).
\item $\square$ Je peux présenter oralement une défense de la culture de paix en 2-3 minutes (structure : accroche, thèse, 2 arguments, conclusion).
\item $\square$ Je relis mon travail pour éviter les faux amis (\esp{actualidad $\neq$ actualité, pretender $\neq$ prétendre, asistir $\neq$ assister}) et les accents manquants.
\end{itemize}
\end{aretenir}

\findecours

\end{document}
